\documentclass[11pt]{article}
\usepackage[T1]{fontenc}
\usepackage[utf8]{inputenc}
\usepackage{lmodern}
\usepackage{microtype}
\usepackage[a4paper,margin=30mm]{geometry}
\usepackage{amsmath,amssymb,amsthm,mathtools}
\usepackage{booktabs,array,tabularx}
\usepackage{enumitem}
\usepackage{xcolor}
\usepackage{aliascnt}
\usepackage{hyperref}
\usepackage[nameinlink,noabbrev]{cleveref}

\hypersetup{
  colorlinks=true,
  linkcolor=blue!55!black,
  citecolor=blue!55!black,
  urlcolor=blue!55!black,
  pdftitle={Finite Gaussian-Rational Approximation: Denominator Heights and a Constructive Hurwitz--Lakein Theorem},
  pdfauthor={Rich Lawrence},
  pdfsubject={Finite residual bounds, Hurwitz convergents, and the Lakein comparison}
}
\setlist{nosep,leftmargin=*}
\setlength{\parskip}{0.35em}
\setlength{\parindent}{1.25em}
\setlength{\emergencystretch}{3em}

\newtheorem{theorem}{Theorem}[section]
\newaliascnt{lemma}{theorem}
\newtheorem{lemma}[lemma]{Lemma}
\aliascntresetthe{lemma}
\newaliascnt{proposition}{theorem}
\newtheorem{proposition}[proposition]{Proposition}
\aliascntresetthe{proposition}
\newaliascnt{corollary}{theorem}
\newtheorem{corollary}[corollary]{Corollary}
\aliascntresetthe{corollary}
\newaliascnt{definition}{theorem}
\newtheorem{definition}[definition]{Definition}
\aliascntresetthe{definition}
\newaliascnt{remark}{theorem}
\newtheorem{remark}[remark]{Remark}
\aliascntresetthe{remark}
\newaliascnt{example}{theorem}
\newtheorem{example}[example]{Example}
\aliascntresetthe{example}

\crefname{theorem}{theorem}{theorems}
\Crefname{theorem}{Theorem}{Theorems}
\crefname{lemma}{lemma}{lemmas}
\Crefname{lemma}{Lemma}{Lemmas}
\crefname{proposition}{proposition}{propositions}
\Crefname{proposition}{Proposition}{Propositions}
\crefname{corollary}{corollary}{corollaries}
\Crefname{corollary}{Corollary}{Corollaries}
\crefname{definition}{definition}{definitions}
\Crefname{definition}{Definition}{Definitions}
\crefname{remark}{remark}{remarks}
\Crefname{remark}{Remark}{Remarks}
\crefname{example}{example}{examples}
\Crefname{example}{Example}{Examples}

\newcommand{\Qi}{\mathbb{Q}(i)}
\newcommand{\Zi}{\mathbb{Z}[i]}
\newcommand{\Q}{\mathbb{Q}}
\newcommand{\N}{\mathbb{N}}
\newcommand{\Z}{\mathbb{Z}}
\newcommand{\Norm}[1]{\operatorname{N}\!\left(#1\right)}
\newcommand{\Ree}{\operatorname{Re}}
\newcommand{\Imm}{\operatorname{Im}}
\newcommand{\diamq}{\operatorname{diag}^{2}}
\newcommand{\eps}{\varepsilon}
\newcommand{\leanref}[1]{\par\smallskip\noindent
  \emph{Formal verification:}\par\nobreak
  {\raggedright\nolinkurl{#1}.\par}}

\title{Finite Gaussian-Rational Approximation\\
\large Denominator-height bounds and a constructive Hurwitz--Lakein theorem}
\author{Rich Lawrence}
\date{August 22, 2026}

\begin{document}
\maketitle

\begin{abstract}
The finite rational-box Fundamental Theorem of Algebra supplies arbitrarily
small polynomial residuals at displayed Gaussian rationals without first
postulating completed complex roots.  We first make the arithmetic cost of
those samples explicit.  For a displayed common-denominator fraction
$q=(a+ib)/D$, put $H(q)=D^2$.  If $f\in\Qi[Z]$ is nonzero, let $S_f$ be the
least positive integer clearing every reduced coefficient-coordinate
denominator.  For $f(q)\ne0$,
\[
                 1\le S_f^2H(q)^{\deg f}\Norm{f(q)}.
\]
The scaled residual is exactly the quadrance of a computed Gaussian integer,
and the family $f(Z)=Z^d/S$, $q=1/D$ attains equality for every positive
$S,D$.
An intrinsic version executes Gaussian Euclidean reduction, canonicalizes
the remaining unit, and replaces $H(q)$ by the resulting
presentation-independent denominator quadrance.  Conversely, nearest-grid
rounding constructs a fraction of
exact height $(\lfloor(1/2)/r\rfloor+1)^2$ inside the returned Newton band;
this is the least denominator giving the required uniform strict half-mesh.
It proves the improved bound
$\Norm{f(q)}\le2L_1(f,A)\eps^2$, and a bounded square-threshold scan chooses
$\eps$ at the expected square-root scale for any displayed residual target.
The same argument is extended to protected zero and pole tracks of a rational
function, where the nonvanishing opposite polynomial contributes computable
upper and lower quadrance bounds.

We then prove a finite Hurwitz--Lakein comparison.  Nearest-Gaussian-integer
division gives a terminating continued-fraction trace for every displayed
Gaussian fraction $z$.  If $p_n/q_n$ is any generated convergent and
$p/q$ is any displayed Gaussian-integral competitor with $q\ne0$ and
$\Norm q<\Norm{q_n}$, then
\[
              \Norm{q_nz-p_n}\le\Norm{qz-p}.
\]
The proof uses only the finite Euclidean trace, adjacent-continuant
identities, exact fundamental-square inequalities, and a finite Hermitian-form
classification of the exceptional quadrance-five case.

We also recover a finite rational analogue of the all-best result.  The
strict-below/weak-at-equal-height definition of a best approximation is
proved equivalent to minimality of the displayed vector $(qz-p,q)$ in its
closed error--denominator cylinder.  For $z=N/D$, the exact pair $(N,D)$ and
the nearest-integer half-square give explicit global denominator and
numerator bounds.  A terminating finite selector therefore returns all and
only the best displayed pairs.  A second finite filter returns exactly the
pairs at the next occupied best denominator level above any supplied lower
pair.  Applied to every finite response of any RAT returned by the RFAP-1
FTA or by one of its signed rational-function sites, this becomes a
RAT-facing program which emits only full best approximants to the displayed
rational anchor of that response.  From any supplied best pair, it returns a
representative of the next occupied best level, or proves that the pair
already represents the response exactly.  Exact linear FTA and
rational-function sites receive the same arithmetic package directly at
their displayed terminal point.  This is deliberately a statement about
finite observations, not unnamed limiting roots.
Each consecutive pair supplies a finite Dirichlet-window certificate
for every natural denominator-quadrance request between its endpoint levels.
A normalized nearest-Gaussian-integer calculation gives a rational quadratic
boundary, without naming either root.  The displayed certificate $102/1189$
then bounds the exact window product by $1189/102$ times its
coefficient-determinant quadrance; that determinant quadrance is exactly one
or two.  A second finite calculation rotates the normalized error ratio into
one of eight rational coordinate octants and tests five small coefficient
combinations.  Their four-region polynomial certificates force the
normalized determinant gap to be at least $205/256$, so an index-one window
is at most $256/205$.  In index two, four centered unit certificates first
force both normalized coordinate quadrances to be at least $15/56$; an
oriented coordinate calculation with the affine minorant $2(U+1)/7$ then
raises the gap to $205/128$ and again bounds the window by $256/205$.  Thus
the universal squared window bound is $256/205$.  More exactly, every
displayed rational $\tau$ satisfying
\[
 \frac34\le\tau<1,\qquad \tau^2-12\tau+9\ge0
\]
obeys the root-free universal certificate
$\tau\mathcal D_z(c_-,c_+)\le1$.  These rational cuts approach Chevallier's
optimal squared constant from the finite side without constructing the
quadratic endpoint.  Conversely, every displayed rational $\tau$ on the
sharper side of that boundary and every natural scale determine, by an
explicit Pell recurrence, a displayed rational source and a consecutive
index-one window beyond the requested scale for which
$1<\tau\mathcal D$.  More generally, any requested finite number of such
windows is realized at strictly increasing scales by one displayed rational
source, and one proper positive-content rational box preserves every window
for every displayed rational source inside it.  On the universal side, a
finite maximum over the complete best-pair table proves that every displayed
rational source and every natural squared request $R>1$ has a nonzero
denominator below $R$ satisfying
$\tau R\Norm{qz-p}^{2}\le1$ for every admissible rational cut $\tau$.
The source used by the converse construction is still allowed to vary with
the requested scale and finite run length.
An exact parity calculation also identifies the index-two ambient
coefficient pairs with the ordinary span together with its centered coset;
two putative successive index-two transitions are placed in an explicit
unit-and-divisible-coefficient normal form whose remaining Gaussian step has
quadrance one, two, or four.  Exact normalization and a finite cylinder
calculation exclude all twelve values.  Hence an index-two transition is
necessarily followed by an index-one transition.  Two adjacent window
products obey $\tau(\mathcal D_0+\mathcal D_1)\le2$, and a chain of $n$
windows obeys $\tau\sum_j\mathcal D_j\le n$ for every such $\tau$;
$512/205$ and $(256/205)n$ remain convenient explicit consequences.  From
one supplied adjacent
pair, a coefficient-local generator contains every pair at the next best
level.  Exact denominator and error comparisons within that local list return
exactly that level, and its candidate list has an explicit quadratic size
bound independent of the global numerator cutoff.  The terminal Hurwitz pair
satisfies the stronger best predicate, although the
ordinary Hurwitz trace need not enumerate the whole returned list.

The results assume neither an irrational source nor an infinite continued
fraction.  All constructions are executable and formally verified in Lean
with axiom closure \texttt{[propext]}, without either classical choice or
\texttt{Quot.sound}.
\end{abstract}

\tableofcontents

\section{Introduction}
\label{sec:introduction}

The companion paper~\cite{RFAPOne} constructs the finite objects that replace
the usual completed zeros and poles of a polynomial or rational function.  A
higher-degree site is represented by a root-approximation track (RAT): every
displayed positive rational accuracy request produces a smaller rational box
with clear boundary, winding one, and a coefficient-controlled residual at
its rational anchor.  That result establishes arbitrarily accurate finite
responses, but it leaves open two arithmetic questions.  How large a
denominator is necessary or sufficient for a requested residual?  And does a
natural finite continued-fraction calculation select samples that compare
well with all competitors having smaller denominators?

This paper answers both questions for displayed Gaussian-rational data.  The
first answer is a coefficient-scaled denominator-height theory.  Its
coefficient scale is the least common multiple of the reduced coordinate
denominators, and its lower obstruction is an exact Gaussian-integral norm
identity.  The common-denominator bound is attained by an explicit family;
an executable Gaussian Euclidean reduction strengthens it to a
presentation-independent denominator height.  In the
other direction, nearest-grid rounding constructs a uniformly least mesh
denominator from an RFAP-1 Newton band, retains a sharper residual constant,
and converts a displayed residual target into a square-root-scaled request.
Protected B\'ezout bounds extend the same statements to zero and pole tracks
of a rational function.

The second answer is a finite Hurwitz--Lakein theorem.  Starting with an
uncancelled fraction over the Gaussian integers, deterministic
nearest-Gaussian-integer division produces a finite Euclidean trace and its
adjacent continuants.  Every generated convergent has no larger
division-free residual quadrance than any displayed competitor of strictly
smaller denominator quadrance.  The usual language of an infinite continued
fraction for a complex number is not needed: the source is finite fraction
data, the trace terminates, and the comparison concerns one displayed stage
and one displayed competitor at a time.

The stronger question is then answered for the same finite input class.  We
adopt the usual two-tier definition of a best approximation, identify it
coefficientwise with Gaussian-lattice cylinder minimality, and compute the
complete finite collection for a displayed rational source.  This is not a
claim that ordinary Hurwitz digits find them all.  Instead the exact source
pair supplies a global denominator cutoff, a cleared-error identity supplies
a global numerator cutoff, and a bounded comparison pass decides the full
predicate.  A further bounded pass decides the consecutive-level relation
between best pairs, and each returned successor gives an exact finite
Dirichlet-window product.  Normalizing the endpoint errors and denominators,
then testing one nearest Gaussian integer, bounds that product by $1189/102$
times the coefficient index; an exact Cauchy calculation retains the
underlying root-free quadratic boundary.  Testing five small coefficient
combinations in a finitely selected coordinate octant strengthens the gap to
$205/256$ and bounds index-one windows by $256/205$.  For index two, four
centered unit certificates and an oriented affine-minorant calculation raise
the normalized gap to $205/128$ and give the same window bound.  The
same calculations, parameterized by a displayed rational $\tau$ with
$3/4\le\tau<1$ and $\tau^2-12\tau+9\ge0$, prove the exact root-free family
$\tau\mathcal D\le1$.  Hence two adjacent windows satisfy
$\tau(\mathcal D_0+\mathcal D_1)\le2$, and a supplied chain of $n$ windows
satisfies $\tau\sum_j\mathcal D_j\le n$.  The concrete bounds $256/205$,
$512/205$, and $(256/205)n$ remain simple rational instances.  If the
quadratic polynomial is instead negative, an explicit Pell family constructs,
beyond every supplied finite scale, a displayed rational source and
consecutive window with $1<\tau\mathcal D$.  The terminal Hurwitz continuant
is among these best pairs.

The final statement is compact; making its customary proof obey the same
finite discipline was not.  References to complete quotients, limiting
values, geometric pictures, or cases ``following by symmetry'' conceal
important proof data, especially on the half-open boundary chosen by a
deterministic tie rule.  The formal proof replaces those shortcuts by exact
continuant identities, rational fundamental-square inequalities, and a
complete finite classification of exceptional Hermitian forms.  The
resulting translation difficulties are foreshadowed here and discussed
directly in \cref{sec:discussion}.

After this introduction states the terminology and main results, the
remaining sections prove them.  Sections~\ref{sec:lower}--\ref{sec:quotients}
develop the denominator-height theory.  Sections~\ref{sec:hurwitz-trace}--
\ref{sec:lakein-proof} construct the finite Hurwitz trace, its succession
geometry, the Hermitian-form invariant, and the Lakein comparison.
Section~\ref{sec:all-best} proves the minimal-vector correspondence, the
finite all-best selector, and its exact consecutive-level successor filter.
Section~\ref{sec:arithmetic-envelopes} gives explicit arithmetic envelopes
for the RFAP-2 calculations, \cref{sec:cubic} retains the exact cubic
calculation, and
\cref{sec:discussion} closes with the constructive and formal lessons of the
proof.

\subsection{Statement, terminology, and constructive discipline}
\label{sec:discipline}

Every mathematical input in this paper is displayed finite data.  Integers,
rationals, Gaussian integers, Gaussian rationals, polynomial coefficient
lists, rational boxes, continued-fraction digits, and proof certificates are
all records on which the stated algorithms operate.  For a displayed
$z=x+iy\in\Qi$, its \emph{quadrance} is the rational number
\begin{equation}
\label{eq:quadrance}
                         \Norm z=x^2+y^2.
\end{equation}
The residual quadrance of a displayed polynomial $f$ at a displayed sample
$q$ is therefore $\Norm{f(q)}$, computed by finite Horner evaluation followed
by two rational squares.

The polynomial conventions are those of RFAP-1.  A displayed polynomial is
its normalized finite coefficient list, and polynomial equality means
coefficientwise equality of those normalized lists.  Evaluation preserves a
proved polynomial identity, but agreement at evaluated values is never used
to define polynomial equality.  The letter $f$ denotes the polynomial under
the main approximation argument; $P,Q$ denote a numerator and denominator;
$A$ is a fixed ambient rational box; $B$ is a current or returned box; and
$q$ is a displayed Gaussian-rational sample.

Two fraction presentations occur and must be distinguished.
\begin{itemize}
\item A \emph{common-denominator fraction} is written
$q=(a+ib)/D$ with $a,b\in\Z$ and $D\in\N_{>0}$.  Its presentation-dependent
height is $H(q)=D^2$.

\item A \emph{Gaussian fraction} is written $z=N/D$ with $N,D\in\Zi$ and
$D\ne0$.  Its \emph{displayed Gaussian-denominator height} is the
quadrance $h_D(z)=\Norm D\in\N$.
\end{itemize}
The first presentation is convenient for polynomial coefficient clearing and
rational grids.  The second is the natural input to Euclidean division over
$\Zi$.  The code applies an executable extended Euclidean calculation to
$N,D$, divides by its displayed common factor, and rotates the result by one
of the four Gaussian units so that the reduced denominator has positive real
part and nonnegative imaginary part.  If the resulting canonical fraction is
$N_*/D_*$, the \emph{intrinsic Gaussian-denominator height} is

\[
                              h(z)=\Norm{D_*}.
\]

This height is proved independent of the supplied presentation.  Thus no
minimization over a quotient is hidden in the notation: cancellation,
unit-normalization, and the accompanying B\'ezout certificate are finite
operations performed by the stated program.

No completed zero is introduced in the height argument.  Its lower bound is
an integer-divisibility statement about the displayed value $f(q)$, not a
claim about the distance to a postulated algebraic point.  Its upper bound
constructs a rational grid point inside a finite Newton band, rather than
appealing to density in a completed plane.  Similarly, the Hurwitz argument
does not introduce an irrational source or an infinite digit stream.  A
source Gaussian fraction has a finite Euclidean trace; a \emph{stage} is one
finite prefix together with its current complete quotient and continuants;
and a comparison is an exact rational inequality between displayed
quadrances.

The same convention applies to geometric wording:
\begin{itemize}
\item membership in a rational box means four decidable rational
inequalities;
\item a statement about every $q\in B$ is conditional on a supplied displayed
$q$ and its membership inequalities;
\item membership in the Hurwitz fundamental square means four displayed
rational coordinate inequalities; and
\item a finite family, scan, trace, or classification is an explicit list or
natural recursion, not a completed collection or an actually infinite
process.
\end{itemize}

\subsubsection{Related work}
\label{sec:related-work}

The foundational comparison with constructive FTA work is developed more
fully in RFAP-1~\cite{RFAPOne}.  Richman~\cite{Richman} constructs a
choice-free completion and represents a monic polynomial by a completed limit
multiset of approximate factors.  The present height theory instead consumes
the finite rational-box interface before any completion is formed.  Perrucci
and Roy~\cite{PerrucciRoyQuantitative} obtain quantitative degree bounds for
an effective real-algebraic FTA.  Our lower height obstruction is narrower and
more elementary: the least common denominator of the displayed coefficient
coordinates turns its scaled residual exactly into the quadrance of one
Gaussian integer.  The sufficient bound, however, depends essentially on the
finite Newton-band construction of RFAP-1.

Hurwitz~\cite{HurwitzComplexCF} introduced the nearest-Gaussian-integer
continued fraction and its square fundamental region.  LeVeque's two-part
study~\cite{LeVequeGaussianCF} develops continued fractions and approximation
in $\Q(i)$.  Lakein~\cite{LakeinApproximation} isolates the approximation
properties needed here: Hurwitz convergents satisfy a second-kind comparison
against Gaussian-rational competitors, even though the complex theory does
not reproduce every familiar real best-approximation property.  The headline
inequality of this paper is the finite, calculation-based counterpart of that
comparison.

Later work places Hurwitz's algorithm in broader geometric and dynamical
settings.  Dani and Nogueira~\cite{DaniNogueira} treat general complex
continued-fraction algorithms and their relation to binary quadratic forms.
Ei, Ito, Nakada, and Natsui~\cite{EiItoNakadaNatsui} construct the natural
extension of the Hurwitz map and analyze its fundamental-region geometry.
Those developments explain why admissible digit succession is subtler than
the single bound $\Norm\rho\le1/2$.  They work with completed complex inputs,
infinite expansions, or dynamical spaces; our proof extracts only the finite
successor inequalities needed for a displayed rational trace.

Hurwitz convergents need not enumerate every best complex approximation.
Chevallier~\cite{ChevallierGaussLattices} identifies best approximations with
minimal vectors in a two-dimensional Gaussian lattice and gives a
lattice-successor algorithm reaching all of them.  The present paper proves
the corresponding minimal-vector equivalence directly for displayed
rational data, gives a complete bounded selector, and filters it into the
exact consecutive occupied denominator levels.  This recovers the all-best
conclusion and its finite successor relation for that input class, including
the local empty-cylinder characterization.  It also supplies a coefficientwise
Cramer certificate and proves constructively that the determinant quadrance
of consecutive pairs is exactly one or two.  This recovers Chevallier's sharp
index dichotomy for displayed rational sources without his Minkowski or
module-basis route.  We also prove the centered-coset description in the
index-two case by displayed coordinate parity.  Four explicit centered-coset
candidates then give a square-root-free error-ratio bound and reduce the
coefficient in a putative repeated-index transition to twelve displayed
Gaussian integers.  Normalizing the two error and denominator coordinates
then reduces these twelve values, by displayed unit associates, to the three
coefficients $1+i$, $2$, and $2+2i$.  A finite exact cylinder calculation
excludes all three, so an index-two transition is followed by an index-one
transition.  This proves the displayed-rational counterpart of Chevallier's
no-repeated-index lemma.  We also formalize its coefficient-local candidate
generation consequence, give it a coefficient-local validation pass, and
prove that the resulting executable update skips no next best level.  It is
a bounded coefficient enumeration rather than Chevallier's more economical
structural formula.  We do not treat arbitrary
completed complex inputs or an infinite minimal-vector sequence.
For displayed rational inputs, however, five small coefficient tests in one
of eight finite coordinate octants, together with four oriented centered
unit tests in index two, give the exact rational-cut family
\[
 \frac34\le\tau<1,\quad \tau^2-12\tau+9\ge0
 \quad\Longrightarrow\quad \tau\mathcal D\le1.
\]
The smaller boundary of this cut is conventionally $6-3\sqrt3$, whose
reciprocal is Chevallier's optimal squared constant
$(2+\sqrt3)/3\approx1.24$.  Thus the present finite result recovers the full
upper boundary without constructing that endpoint.  In the reverse
direction, for every displayed rational $\tau$ satisfying the same range and
$\tau^2-12\tau+9<0$, and every natural scale, the Pell recurrence produces
an explicit displayed rational source and a consecutive window beyond that scale with
$1<\tau\mathcal D$.  The strengthened finite form returns an arbitrary
finite same-source run and a proper positive-content rational box preserving
it.  For every displayed rational source in that box, the bounded
coefficient-local successor table gives a stopped regenerative kernel:
genuine successor edges strictly increase denominator quadrance, a terminal
edge resets to the selected critical seed, and the seed is then held.
Natural-rank absorption therefore gives probability one of a bad-window hit
after any burn-in and one positive finite block horizon.  The source and box
may depend on the requested scale and finite run length, and the stop and
reset are explicit boundary rules rather than Haar invariance.  This does not
recover Chevallier's almost-everywhere fixed-source sharpness statement, nor
his formulation for arbitrary completed complex inputs and an infinite
minimal-vector sequence.
Separately,
our Hurwitz theorem
still claims only that each generated convergent is at least as good as every
competitor of strictly smaller denominator quadrance.

Formal continued-fraction developments also exist.  Eberl's Isabelle
formalization~\cite{EberlContinuedFractions} treats real continued fractions,
periodic quadratic irrationals, approximation, and executable applications
such as Pell equations.  The present formalization has a different boundary:
it treats Gaussian fractions and the Hurwitz--Lakein comparison, and it audits
the public theorem specifically to exclude classical choice and completed
infinite objects from its proof dependency.

\subsection{Machine verification}
\label{sec:machine-verification}

The complete machinery described in this paper has been implemented and
verified in Lean.  The accompanying development contains the executable
coefficient-clearing, grid, Gaussian Euclidean, continuant, and stage-trace
programs together with machine-checked proofs of the paper-facing guarantees.
The \emph{Formal verification} markers throughout the paper name the
corresponding Lean declarations.

This is a particularly appropriate use of an automated theorem prover.
Every construction is finite and explicitly checkable: recursive calls
decrease a displayed natural measure, searches traverse displayed finite
lists, polynomial identities are coefficientwise, and geometric assertions
reduce to integer or rational equalities and inequalities.  The proof
assistant can check the algorithms at the same level at which the theorem is
stated, including their termination and the exact dependency closure of the
result.

RFAP paper~1~\cite{RFAPOne} records the same formal-verification problems.
In particular, apparently elementary generic Mathlib routes can introduce
\texttt{Classical.choice} through an elaborated instance or inherited
theorem.  The present development again uses project-owned constructive
rational and Gaussian algebra, direct order lemmas, and focused
\texttt{\#print axioms} manifests.  The public RFAP-2 declarations have the
audited closure \texttt{[propext]}; they contain neither
\texttt{Classical.choice}, \texttt{Quot.sound}, nor a new native decision
axiom.  We do not repeat RFAP-1's general formalization appendix.  The
paper-specific problem---turning customary Hurwitz--Lakein prose into finite
machine-checkable certificates---belongs to the mathematical discussion in
\cref{sec:discussion}.

\leanref{RFAP.ApproximationPaper.scaledHeightResidual_eq_integralNorm;
RFAP.ApproximationPaper.one_le_scaled_gaussianHeight_mul_residual;
RFAP.ApproximationPaper.denominatorHeightBound_isSharp;
RFAP.ApproximationPaper.approximateWithTarget_residual_le;
RFAP.ApproximationPaper.hurwitzConvergent_atLeastAsGoodAs;
RFAP.ApproximationPaper.mem_allBestGaussianApproximations_iff}

\subsection{Definitions and theorem statements}
\label{sec:statements}

We now give, in dependency order, the definitions needed to read the main
results.  This subsection is deliberately a fast path to the main theorem
statements.  It records the essential finite interfaces and defers intuition,
derivation, and the larger calculations to the later proof sections.

\subsubsection{Height data and finite RFAP input}

\begin{definition}[least coefficient scale and displayed heights]
\label{def:intro-height}
For a displayed polynomial
$f(Z)=\sum_{j=0}^{d}c_jZ^{d-j}\in\Qi[Z]$, let $s(c_j)$ be the least common
multiple of the reduced denominators of $\Ree c_j$ and $\Imm c_j$.  The
\emph{coefficient scale} is
\[
                         S_f=\mathop{\rm lcm}_j s(c_j).
\]
Thus $S_fc_j\in\Zi$ for every $j$, and $S_f$ divides every natural scale
with this property.  For a common-denominator fraction $q=(a+ib)/D$, put
$H(q)=D^2$.  For a displayed Gaussian fraction $q=U/V$, put
$h_V(q)=\Norm V$.  The extended Gaussian Euclidean calculation divides
$U,V$ by their displayed common factor, and a deterministic choice among
$1,-1,i,-i$ makes the reduced denominator $V_*$ have positive real part and
nonnegative imaginary part.  Write the canonical result as
$\operatorname{red}(q)=U_*/V_*$ and put $h(q)=\Norm{V_*}$.  With
$A_j=S_fc_j\in\Zi$, define the integral homogeneous numerator
\[
             W_f(U,V)=\sum_{j=0}^{d}A_jU^{d-j}V^j\in\Zi.
\]
\end{definition}

\leanref{RFAP.DenominatorHeight.coefficientScale;
RFAP.DenominatorHeight.coefficientScale_dvd_of_clears;
RFAP.ConstructiveGaussianHeight.Fraction.height;
RFAP.HurwitzRational.Fraction.height;
RFAP.ReducedGaussianFraction.reduceCanonical;
RFAP.GaussianDenominatorHeight.intrinsicHeight;
RFAP.ApproximationPaper.intrinsicGaussianDenominatorHeight}

The geometric input imported from RFAP-1 is the following finite interface.
Its hypotheses are displayed coefficient data for a squarefree polynomial
$f$ of degree at least two, a certified winding state inside an ambient box
$A$, and the finite axis-elimination data required for boundary-clear cuts.

\begin{proposition}[RFAP-1 Newton-band interface]
\label{prop:rfap-input}
For every displayed $\eps>0$, finite Newton extraction computes a centre
$c=x+iy\in\Qi$, a radius $r>0$, and a rational box $B$ such that
\begin{enumerate}[label=(\roman*)]
\item the left, right, bottom, and top sides of $B$ lie respectively at
coordinate distances between $r$ and $2r$ from $c$;
\item $B$ has clear boundary, winding one, and
$\diamq(B)\le\tfrac12\eps^2$; and
\item every displayed $q\in B$ satisfies
\begin{equation}
\label{eq:box-residual}
                \Norm{f(q)}\le2L_1(f,A)\eps^2,
\end{equation}
where $L_1(f,A)>0$ is computed from the displayed coefficients and endpoints
by finite Horner recursions.
\end{enumerate}
\end{proposition}

The returned box, centre, radius, winding calculation, and coefficient bound
are output data.  The proposition does not quantify over an unshown point.

\leanref{RFAP.NewtonBand.extract;
RFAP.NewtonBand.Output.diagonalSq_le_half;
RFAP.CoprimeBounds.normSq_eval_le_four_lipschitz_diagonal;
RFAP.HeightApproximation.approximate}

\begin{theorem}[finite denominator-height package]
\label{thm:height-package}
Let $f\in\Qi[Z]$ be displayed and nonzero, of degree $d$.
\begin{enumerate}[label=(\roman*)]
\item For $q=u/D$ in common-denominator form,
\[
 S_f^2H(q)^d\Norm{f(q)}=\Norm{W_f(u,D)}.
\]
If $f(q)\ne0$, this natural number is at least one.
\item For $q=U/V$ in displayed Gaussian-denominator form,
\[
 S_f^2h_V(q)^d\Norm{f(q)}=\Norm{W_f(U,V)},
\]
and again a nonzero residual makes the common value at least one.
\item For the canonical reduction
$\operatorname{red}(q)=U_*/V_*$,
\[
 S_f^2h(q)^d\Norm{f(q)}=\Norm{W_f(U_*,V_*)}.
\]
The height $h(q)$ is positive, no larger than $h_V(q)$, and equal for any
two displayed fractions representing the same exact value.  A nonzero
residual again makes the common value at least one.
\item For every $d\in\N$ and positive $S,D$, the displayed family
$f(Z)=Z^d/S$, $q=1/D$ has $S_f=S$, $H(q)=D^2$, nonzero residual, and makes
the common value in (i) exactly one.  The unit constant and the displayed
powers of $S_f$ and $H(q)$ are therefore jointly sharp.
\item At a protected numerator or denominator site of a rational function,
the lower obstruction transfers to the shrinking quotient residual after
including the displayed upper majorant of the opposite polynomial.
\end{enumerate}
\end{theorem}

\leanref{RFAP.ApproximationPaper.scaledHeightResidual_eq_integralNorm;
RFAP.ApproximationPaper.one_le_scaled_height_mul_residual;
RFAP.ApproximationPaper.one_le_scaled_gaussianHeight_mul_residual;
RFAP.ApproximationPaper.intrinsicGaussianDenominatorHeight_le_displayed;
RFAP.ApproximationPaper.intrinsicGaussianDenominatorHeight_eq_of_value_eq;
RFAP.ApproximationPaper.scaledIntrinsicGaussianHeightResidual_eq_integralNorm;
RFAP.ApproximationPaper.one_le_scaled_intrinsicGaussianHeight_mul_residual;
RFAP.ApproximationPaper.denominatorHeightBound_isSharp;
RFAP.ApproximationPaper.one_le_scaled_height_mul_quotientResidual}

\begin{theorem}[finite sufficient-height package]
\label{thm:sufficient-height-package}
In the Newton extraction of Proposition~\ref{prop:rfap-input}, put
\[
 D_r=\left\lfloor\frac{1/2}{r}\right\rfloor+1.
\]
Nearest-grid rounding returns a displayed $q_r\in B$ with
\[
 H(q_r)=D_r^2,
 \qquad \Norm{f(q_r)}\le2L_1(f,A)\eps^2.
\]
The number $D_r$ is the least positive denominator whose uniform half-mesh
is strictly smaller than $r$.  For every displayed target $\delta>0$, a
bounded scan computes a reciprocal-integer request $\eps$ at square-root
scale and makes the last residual at most $\delta$.  At a protected
rational-function site the same construction applies to the shrinking
quotient residual after division by the positive protection margin of the
opposite polynomial.
\end{theorem}

\leanref{RFAP.HeightApproximation.gridDenominator_minimal;
RFAP.ApproximationPaper.approximateWithHeight;
RFAP.ApproximationPaper.approximateWithTarget_residual_le;
RFAP.ApproximationPaper.approximateQuotientWithHeight;
RFAP.ApproximationPaper.approximateQuotientWithTarget_residual_le}

The two packages are proved in
\cref{sec:lower,sec:upper,sec:quotients}.

\subsubsection{Finite Hurwitz data}

\begin{definition}[nearest Gaussian integer and fundamental square]
\label{def:nearest-gaussian}
For $x\in\Q$, put
\[
                  [x]_{\rm n}=\left\lfloor x+\frac12\right\rfloor.
\]
Thus a half-integer tie is sent upward.  For $z=x+iy\in\Qi$, define
\[
 [z]_{\Zi}=[x]_{\rm n}+i[y]_{\rm n},
 \qquad
 \rho(z)=z-[z]_{\Zi}.
\]
The chosen half-open fundamental square is
\begin{equation}
\label{eq:fundamental-square}
 \mathfrak F=
 \left\{u+iv:-\frac12\le u<\frac12,
                  -\frac12\le v<\frac12\right\}.
\end{equation}
Finite floor inequalities certify $\rho(z)\in\mathfrak F$ and
$\Norm{\rho(z)}\le1/2$.
\end{definition}

\leanref{RFAP.HurwitzNearest.nearestCoordinate;
RFAP.HurwitzNearest.nearest;
RFAP.HurwitzNearest.InFundamentalSquare;
RFAP.HurwitzNearest.remainder_mem;
RFAP.HurwitzNearest.remainder_normSq_le_half}

\begin{definition}[finite Hurwitz transition and encoding]
\label{def:finite-hurwitz-transition}
Let $z=N/D$ be a displayed Gaussian fraction.  Its digit is
$a=[z]_{\Zi}$ and its integral remainder is
\[
                           R=N-aD\in\Zi.
\]
If $R=0$, the trace is terminal.  If $R\ne0$, the next displayed fraction is
\[
                         T(z)=\frac{D}{R}
                             =\rho(z)^{-1}.
\]
The exact norm identity
\[
                   \Norm R=\Norm{\rho(z)}\Norm D
\]
and $\Norm{\rho(z)}\le1/2$ give $\Norm R<\Norm D$.  Recursion on this
strictly decreasing natural denominator quadrance therefore returns a finite,
nonempty digit list and an exact reconstruction certificate for $z$.
\end{definition}

\leanref{RFAP.HurwitzRational.digit;
RFAP.HurwitzRational.remainder;
RFAP.HurwitzRational.next;
RFAP.HurwitzRational.remainder_norm_lt_denominator;
RFAP.HurwitzRational.encodeResult;
RFAP.HurwitzRational.encode_reconstruct}

\begin{definition}[continuants, generated stages, and denominator ratio]
\label{def:hurwitz-stage}
For finite digits $a_0,\ldots,a_n\in\Zi$, define adjacent continuants by
\[
 (p_{-1},q_{-1})=(1,0),\qquad (p_0,q_0)=(a_0,1),
\]
and
\begin{equation}
\label{eq:continuant-recurrence}
 (p_{j+1},q_{j+1})=
 a_{j+1}(p_j,q_j)+(p_{j-1},q_{j-1}).
\end{equation}
A \emph{Hurwitz stage} records one displayed complete quotient, the adjacent
pairs $(p_{j-1},q_{j-1})$ and $(p_j,q_j)$, their unit-determinant
certificate, and the exact reconstruction identity for the source.  A stage
is \emph{generated} when it is the initial stage or is obtained from a
generated stage by one nonterminal transition.  Its \emph{denominator ratio}
is the displayed Gaussian rational
\begin{equation}
\label{eq:denominator-ratio}
                         r_j=\frac{q_{j-1}}{q_j}.
\end{equation}
\end{definition}

\leanref{RFAP.HurwitzRational.ConvergentPair;
RFAP.HurwitzTrace.advancePair;
RFAP.HurwitzTrace.Stage;
RFAP.HurwitzGeometry.Generated;
RFAP.HurwitzGeometry.Stage.denominatorRatio}

\begin{definition}[division-free error and Lakein comparison]
\label{def:lakein-comparison}
For a displayed source $z$ and an integral pair $(p,q)\in\Zi^2$, define its
division-free error and error quadrance by
\[
             E_z(p,q)=zq-p,
             \qquad \mathcal E_z(p,q)=\Norm{E_z(p,q)}.
\]
A convergent pair $(p_n,q_n)$ is \emph{at least as good as} a competitor
$(p,q)$ when
\[
                      \mathcal E_z(p_n,q_n)
                      \le\mathcal E_z(p,q).
\]
This is a second-kind comparison: denominator size is used only to select the
competitors, while the compared quantity has been cleared of division.
\end{definition}

\leanref{RFAP.HurwitzApproximation.error;
RFAP.HurwitzApproximation.AtLeastAsGoodAs}

\begin{theorem}[finite Hurwitz--Lakein inequality]
\label{thm:finite-lakein}
Let $z=N/D$ be a displayed Gaussian fraction, let a generated stage have
current continuant $(p_n,q_n)$, and let $(p,q)\in\Zi^2$ be any displayed
competitor.  If
\[
                         q\ne0,
              \qquad    \Norm q<\Norm{q_n},
\]
then
\begin{equation}
\label{eq:finite-lakein}
                        \Norm{zq_n-p_n}
                     \le\Norm{zq-p}.
\end{equation}
At a terminal stage the left side is zero.  Every earlier generated stage is
covered by the finite fundamental-square and Hermitian-form invariant proved
below.
\end{theorem}

\leanref{RFAP.HurwitzGeometry.Stage.atLeastAsGoodAs_of_generated;
RFAP.ApproximationPaper.hurwitzConvergent_atLeastAsGoodAs}

The proof is completed in \cref{sec:lakein-proof}.

\begin{definition}[best approximation and source-lattice minimality]
\label{def:best-approximation}
Let $c=(p,q)\in\Zi^2$ with $q\ne0$.  It is a \emph{best approximation} to
the displayed source $z$ when every $d=(a,b)\in\Zi^2$ with $b\ne0$
satisfies the two-tier comparison
\begin{align}
 \Norm b<\Norm q
   &\Longrightarrow \mathcal E_z(c)<\mathcal E_z(d),
   \label{eq:best-strict}\\
 \Norm b\le\Norm q
   &\Longrightarrow \mathcal E_z(c)\le\mathcal E_z(d).
   \label{eq:best-weak}
\end{align}
The associated displayed source-lattice vector is $(E_z(c),q)$.  It is
\emph{minimal} when every nonzero displayed pair $d$ in its closed cylinder,
\[
       \mathcal E_z(d)\le\mathcal E_z(c),\qquad
       \Norm b\le\Norm q,
\]
has equality in both coordinates.  This is a definition on coefficient
records; it does not first form a completed lattice carrier.

Two best pairs $c_-$ and $c_+$ occupy \emph{consecutive best levels} when
\[
  \Norm{q_-}<\Norm{q_+}
\]
and no best pair has denominator quadrance strictly between these two values.
All unit associates and any other ties remain at the same level.  For a
displayed rational source, this is equivalently the local condition that no
nonzero coefficient pair $c$ satisfies both
\[
 \mathcal E_z(c)<\mathcal E_z(c_-),
 \qquad \Norm{q_c}<\Norm{q_+}.
\]
\end{definition}

\leanref{RFAP.BestApproximation.IsBest;
RFAP.BestApproximation.sourceVector;
RFAP.BestApproximation.sourceVector_injective;
RFAP.BestApproximation.pairNonzero_iff_sourceVector_ne_zero;
RFAP.BestApproximation.InClosedCylinder;
RFAP.BestApproximation.IsGaussLatticeMinimal;
RFAP.BestApproximation.AreConsecutiveBest}

\begin{theorem}[finite minimal-vector correspondence and all-best selector]
\label{thm:finite-all-best}
For every displayed Gaussian fraction $z=N/D$ and every displayed pair $c$,
the following are equivalent:
\begin{enumerate}
\item $c$ is a best approximation in the sense of
      \eqref{eq:best-strict}--\eqref{eq:best-weak};
\item the source-lattice vector $(E_z(c),q)$ is minimal in its closed
      error--denominator cylinder;
\item $c$ occurs in the executable finite list $\operatorname{AllBest}(z)$.
\end{enumerate}
If $B=\Norm D$ and
\[
                    P=2(\Norm N\,B+B),
\]
the selector need inspect only denominators whose two component magnitudes
are at most $B$ and numerators whose two component magnitudes are at most
$P$.  Its rectangular candidate list has exactly
\[
                         16(B+1)^2(P+1)^2
\]
entries before filtering; duplicated zero components are harmless.  The
terminal pair of every generated Hurwitz trace is best.  Every generated
pair, terminal or not, retains the weaker conclusion of
\cref{thm:finite-lakein}.

For any supplied lower pair $c_-$, a second executable list
$\operatorname{NextBest}(z,c_-)$ satisfies
\begin{equation}
\label{eq:next-best-membership}
 c_+\in\operatorname{NextBest}(z,c_-)
 \quad\Longleftrightarrow\quad
 c_-,c_+\text{ occupy consecutive best levels}.
\end{equation}
Thus the finite all-best collection also has an exact computed successor
relation.  This first implementation filters the global rectangle; the later
coefficient-local update computes the same relation from an adjacent state.

For the rational anchor $z_R$ carried by any finite response of an RFAP-1 RAT,
the same calculation emits all and only best approximants to $z_R$.
Selecting the first member of the next occupied level defines a total optional
step: if the supplied pair $c$ is best, the step returns a consecutive best
pair, or else $\mathcal E_{z_R}(c)=0$.  Conversely, the option is empty exactly
when this error is zero.  Thus no provisional or merely nearby candidate can
be emitted by the RAT-response interface.  The target here is the displayed
rational response carried by the finite call; no completed root is introduced.
This applies to every degree-indexed RAT returned by the finite FTA and to
every higher-degree signed numerator or denominator site returned by RFAP-1;
their exact linear branches use the same calculation at the displayed
terminal point.

For every returned consecutive pair, the displayed coefficient determinant
satisfies the sharp finite alternative
\[
       \Norm{p_-q_+-p_+q_-}\in\{1,2\}.
\]
\end{theorem}

\leanref{RFAP.BestApproximation.isBest_iff_isGaussLatticeMinimal;
RFAP.BestApproximation.candidates_length;
RFAP.BestApproximation.mem_allBest_iff_isBest;
RFAP.BestApproximation.mem_consecutiveSuccessors_iff;
RFAP.BestApproximation.AreConsecutiveBest.coefficientIndexQuadrance_eq_one_or_two;
RFAP.HurwitzGeometry.Stage.isWeakLowerBest_of_generated;
RFAP.HurwitzGeometry.Stage.isBest_of_generated_of_terminal;
RFAP.ApproximationPaper.mem_allBestGaussianApproximations_iff;
RFAP.ApproximationPaper.mem_consecutiveBestSuccessors_iff;
RFAP.RATBestApproximation.respondToPoint;
RFAP.RATBestApproximation.respondToTrack;
RFAP.ApproximationPaper.strengthenFTAResult;
RFAP.ApproximationPaper.strengthenRationalSite;
RFAP.ApproximationPaper.consecutiveBest_coefficientIndexQuadrance_eq_one_or_two}

The proof is completed in \cref{sec:all-best}.

\subsection{Proof-route glossary}
\label{sec:proof-route-glossary}

The entries below are a dependency-ordered map of the proof.  They summarize
the formal definitions rather than replacing them.

\paragraph{Displayed data and certificate.}
An object is \emph{displayed} when it is finite integer or rational data on
which the stated calculation runs.  A \emph{certificate} is displayed data
together with finite checks that imply its advertised property; later
algorithms consume certificates and return new ones.

\paragraph{Quadrance and residual.}
The \emph{quadrance} of $x+iy$ is $x^2+y^2$.  A polynomial \emph{residual}
is $\Norm{f(q)}$ at a displayed sample.  A continued-fraction
\emph{division-free residual} is $\Norm{zq-p}$ for a displayed integral pair
$(p,q)$.

\paragraph{Common-denominator height and Gaussian denominator height.}
For $q=(a+ib)/D$ with positive integer $D$, the
\emph{displayed common-denominator height} is $H(q)=D^2$.  For
$z=N/D$ with nonzero Gaussian-integer $D$, the \emph{displayed Gaussian
denominator height} is $h_D(z)=\Norm D$.  Gaussian Euclidean reduction and
a four-unit normalization produce canonical coordinates $N_*/D_*$; the
\emph{intrinsic Gaussian denominator height} is $h(z)=\Norm{D_*}$.
Equal displayed Gaussian fractions have identical canonical coordinates and
hence the same intrinsic height.

\paragraph{Coefficient scale and height obstruction.}
The \emph{coefficient scale} $S_f$ is the least common multiple of all
reduced real and imaginary coefficient-coordinate denominators.  Homogeneous
Horner evaluation then produces a Gaussian integer $W_f$ satisfying
$W_f=D^dS_ff(q)$ for a displayed common denominator $D$, or
$W_f=V^dS_ff(q)$ for a displayed Gaussian denominator $V$.  Applying the
same calculation to the canonical reduced denominator $D_*$ replaces the
presentation-dependent height by $h(q)$.  Its quadrance is exactly the
corresponding scaled residual.  If the displayed residual is nonzero,
$\Norm{W_f}\ge1$ is the \emph{height obstruction}.

\paragraph{RAT, Newton band, and grid fraction.}
A \emph{RAT} is the RFAP-1 finite request-and-response program returning
strictly refining, boundary-clear, winding-one rational boxes.  A
\emph{Newton band} is its coefficient-certified box around a displayed Newton
centre.  The \emph{grid fraction} rounds both centre coordinates with the
nearest-integer rule and denominator $\lfloor(1/2)/r\rfloor+1$, and lies
strictly inside that band.  This denominator is least for the uniform strict
half-mesh condition.  A \emph{square-threshold request} uses a bounded scan
for the least reciprocal-integer $\eps$ sufficient for a displayed quadratic
residual target.

\paragraph{Protected quotient residual.}
At a numerator site for $P/Q$, a protected RFAP-1 state supplies a positive
lower bound and a finite upper majorant for $\Norm{Q(q)}$.  These calibrate the
\emph{quotient residual} $\Norm{P(q)}/\Norm{Q(q)}$.  At a pole the roles are
exchanged, so the shrinking quantity is the reciprocal quotient residual.

\paragraph{Nearest Gaussian integer and fundamental square.}
The \emph{nearest Gaussian integer} rounds each rational coordinate by
$\lfloor x+1/2\rfloor$.  Its residual lies in the half-open
\emph{fundamental square} $[-1/2,1/2)^2$.  Retaining which sides are open and
closed is essential to the strict successor inequalities.

\paragraph{Integral remainder, complete quotient, and finite trace.}
For $z=N/D$ and digit $a$, the \emph{integral remainder} is $R=N-aD$.  If it
is nonzero, the next \emph{complete quotient} is $D/R$.  The strict decrease
$\Norm R<\Norm D$ makes iteration a terminating \emph{finite trace}.

\paragraph{Continuant and generated stage.}
A \emph{continuant} is an integral pair of a numerator and denominator
produced by the recurrence \eqref{eq:continuant-recurrence}.  A \emph{generated stage}
retains two adjacent continuants, their unit determinant, the current complete
quotient, and an exact reconstruction of the original source.

\paragraph{Competitors and Lakein factor.}
The adjacent unit determinant expresses every displayed competitor $(p,q)$
with explicit Gaussian-integer coefficients $s,t$.  Its error factors as the
current convergent error times an explicit \emph{Lakein factor}.  Proving the
competitor is no better therefore reduces to proving that factor has
quadrance at least one.

\paragraph{Denominator growth and denominator ratio.}
The \emph{denominator-growth certificate} is
$\Norm{q_{j-1}}<\Norm{q_j}$.  The \emph{denominator ratio}
$r_j=q_{j-1}/q_j$ is consequently a displayed Gaussian rational of quadrance
less than one.  Its finite backward succession is the state variable used by
the exceptional proof.

\paragraph{Successor geometry and forbidden pair.}
The fundamental-square inequalities imply exact exterior-disc conditions for
successive digits.  A \emph{forbidden pair} is one of the finitely classified
small adjacent-digit pairs contradicted by those conditions, including the
direction of every half-open boundary tie.

\paragraph{Hermitian root form and pullback.}
The remaining quadrance-five comparison is one of eight displayed integral
\emph{Hermitian root forms} evaluated at $r_j$.  A digit \emph{pulls back} a
form by substituting the exact backward-ratio transformation and clearing its
positive denominator quadrance.  Coefficientwise calculation proves the
pullback identity.

\paragraph{Exceptional form and root-form invariant.}
Most pulled-back root forms have an immediate nonnegative slack certificate.
The remaining twenty-four \emph{exceptional forms} are checked across one
more finite predecessor step; every unsafe result would require a forbidden
digit pair.  Induction on the displayed generated-stage proof yields the
\emph{root-form invariant}.

\paragraph{Finite Lakein inequality.}
The root-form invariant supplies the last quadrance-five Lakein-factor case;
the smaller coefficient cases follow from elementary fundamental-square
geometry and the larger cases from denominator estimates.  The factor has
quadrance at least one, so multiplying by the current error proves
\eqref{eq:finite-lakein}.  A terminal trace instead gives exact zero error.

\paragraph{Best approximation, minimal vector, and finite successors.}
A \emph{best approximation} is strictly better than every lower-height pair
and weakly better than every pair of no greater height.  Its displayed
source-lattice vector is \emph{minimal} when no nonzero vector enters both
the interior error range and the no-greater denominator range of its closed
cylinder.  For a rational source, the \emph{all-best selector} checks this
condition over one source-bounded rectangle; its output contains precisely
the best pairs, including the terminal Hurwitz pair but not only Hurwitz
pairs.  Two best pairs occupy \emph{consecutive best levels} when their
denominator quadrances are strictly ordered and no best pair has an
intermediate denominator quadrance.  The \emph{finite successor selector}
returns every pair at the consecutive level above a supplied lower pair.  A
finite lexicographic descent proves that this global level condition is
equivalent to emptiness of the corresponding \emph{mixed open cylinder}.
The coefficient determinant of two consecutive pairs has quadrance exactly
one or two; a computed B\'ezout complement and one reduced Gaussian remainder
prove this without constructing a quotient lattice.

\paragraph{Centered coefficient coset.}
When the consecutive determinant $\delta$ has quadrance two, divisibility by
$\delta$ is exactly even parity of the sum of the displayed real and
imaginary coordinates.  The \emph{centered coefficient coset} consists of
the pairs $x$ for which
$\delta x=a u+bv$ with both cleared coefficients nondivisible by $\delta$.
Together with the ordinary integral span it contains every ambient displayed
pair, and every pair of nondivisible coefficients constructs such an $x$.
For two putative successive index-two transitions, this classification gives
the \emph{repeated-index normal form}
$w=\varepsilon u+\delta k v$ with $\Norm\varepsilon=1$ and
$\Norm k\in\{1,2,4\}$.  The bound on $k$ comes from four explicit centered
candidates and the square-root-free error estimate
$\mathcal E_z(u)<4\mathcal E_z(v)$.  Dividing the two recurrences by the
displayed nonzero error and denominator references gives \emph{normalized
cylinder data}.  Unit rotation reduces its transition coefficient to
$1+i$, $2$, or $2+2i$; the corresponding finite cylinder exclusions rule
out all three cases.

\paragraph{Coefficient-local successor generator.}
Given adjacent best pairs $u,v$, the \emph{local successor candidates} are
computed from the eight leading coefficients of quadrance one or two and a
bounded trailing coefficient.  They use either the integral form $au+bv$ or
the centered form obtained by exact quadrance-two division.  Every actual
next best pair occurs in this list.  The \emph{local successor output} retains
the improving candidates having least denominator quadrance in this list and,
among candidates on that level, least error quadrance.  This is a local finite
calculation, and its output is proved to be exactly the next best level.

\paragraph{Dirichlet window and finite sharpness witness.}
For consecutive best pairs $c_-,c_+$, the \emph{squared Dirichlet-window
product} is
$\mathcal D_z(c_-,c_+)=\Norm{q_+}\mathcal E_z(c_-)$.  A displayed rational
threshold $\tau$ on the certified side of the quadratic boundary satisfies
$\tau\mathcal D\le1$ for every such window.  A \emph{finite sharpness
witness} for a threshold on the other side consists of an explicitly
computed displayed rational source and consecutive pair whose denominator
quadrance exceeds a supplied natural scale and for which
$1<\tau\mathcal D$.  The witness source is output data and may change when
the requested scale changes.

\paragraph{Arithmetic envelope.}
An \emph{arithmetic envelope} is a natural-number recurrence computed from
the displayed input before the calculation it bounds.  It limits recursive
clauses or the bit widths of the integral records produced by the grid,
homogeneous-evaluation, and Hurwitz programs.  It is not a maximum recovered
by inspecting an already completed trace.

\section{The denominator-height obstruction}
\label{sec:lower}

This section turns the glossary's coefficient scale into the universal
height obstruction.  Finite coefficient clearing and homogeneous Horner
evaluation construct one displayed Gaussian integer.  If the residual is
nonzero, that integer is nonzero and hence has quadrance at least one.  The
entire lower bound is the quadrance of this coefficientwise identity.  We
first make the scale least, then prove both common- and Gaussian-denominator
forms of the identity, and finally exhibit equality cases.

\begin{definition}[least coefficient-denominator scale]
\label{def:least-coefficient-scale}
Write a rational coordinate in its reduced numerator--positive-denominator
form and denote its denominator by $\operatorname{den}(-)$.  For
$c\in\Qi$, put
\[
 s(c)=\operatorname{lcm}\bigl(\operatorname{den}(\Ree c),
                              \operatorname{den}(\Imm c)\bigr).
\]
For the normalized displayed coefficient list $(c_0,\ldots,c_d)$, define
\begin{equation}
\label{eq:least-coefficient-scale}
                     S_f=\operatorname{lcm}_{0\le j\le d}s(c_j).
\end{equation}
The finite fold starts with $1$, so $S_f$ is positive even for an empty
intermediate list.
\end{definition}

\begin{lemma}[coefficient-scale minimality]
\label{lem:coefficient-scale-minimal}
The number $S_f$ clears every coefficient: $S_fc_j\in\Zi$.  If a natural
number $T$ clears every reduced real and imaginary coefficient denominator,
then $S_f\mid T$.
\end{lemma}

\begin{proof}
For one coefficient, each coordinate denominator divides its least common
multiple.  At each step of the finite list fold, the current least common
multiple and the remaining-list scale both divide the new one, proving that
$S_f$ clears every coordinate.  Conversely, if $T$ is divisible by the two
coordinate denominators of every coefficient, it is divisible by each
$s(c_j)$.  Induction through the same fold and the defining universal
property of the least common multiple give $S_f\mid T$.
\end{proof}

\leanref{RFAP.DenominatorHeight.coefficientDenominator;
RFAP.DenominatorHeight.lcmScale;
RFAP.DenominatorHeight.lcmScale_clears;
RFAP.DenominatorHeight.coefficientScale_dvd_of_clears}

\begin{definition}[homogeneous residual numerator]
\label{def:homogeneous-numerator}
Let
\[
 f(Z)=\sum_{j=0}^{d}c_jZ^{d-j},
\]
and put $A_j=S_fc_j\in\Zi$.  For $U,V\in\Zi$, the \emph{homogeneous residual
numerator} is the Gaussian integer
\begin{equation}
\label{eq:homogeneous-numerator}
 W_f(U,V)=\sum_{j=0}^{d}A_jU^{d-j}V^j\in\Zi.
\end{equation}
It is computed by a high-to-low homogeneous Horner recursion on the displayed
coefficient list; the sum is notation for the resulting finite fold.
\end{definition}

\leanref{RFAP.DenominatorHeight.homogeneousEval;
RFAP.DenominatorHeight.residualNumerator;
RFAP.GaussianDenominatorHeight.homogeneousEval;
RFAP.GaussianDenominatorHeight.residualNumerator}

\begin{lemma}[homogeneous evaluation identity]
\label{lem:homogeneous-evaluation}
For a common-denominator fraction $q=u/D$, with $D$ regarded as the Gaussian
integer $D+0i$, the displayed numerator satisfies
\begin{equation}
\label{eq:homogeneous-evaluation}
                        W_f(u,D)=D^dS_ff(q)
\end{equation}
after embedding both sides in $\Qi$.  More generally, a Gaussian fraction
$q=U/V$ with $V\ne0$ satisfies
\begin{equation}
\label{eq:gaussian-homogeneous-evaluation}
                        W_f(U,V)=V^dS_ff(q).
\end{equation}
\end{lemma}

\begin{proof}
Coefficient clearing first proves by induction on the normalized coefficient
list that ordinary Horner evaluation of the cleared list is $S_ff(q)$.
Independently, induction on that same list proves that homogeneous Horner
evaluation at $(U,V)$ is $V^d$ times ordinary evaluation at $U/V$.  Composing
the two coefficientwise identities gives
\eqref{eq:gaussian-homogeneous-evaluation}; taking $V=D+0i$ gives
\eqref{eq:homogeneous-evaluation}.  No conclusion is inferred from equality
on a collection of evaluated points.
\end{proof}

\leanref{RFAP.DenominatorHeight.eval_lcmClear;
RFAP.DenominatorHeight.homogeneousEval_cons_eq_scaled;
RFAP.DenominatorHeight.residualNumerator_identity;
RFAP.GaussianDenominatorHeight.residualNumerator_identity}

\begin{lemma}[nonzero residual gives nonzero integral numerator]
\label{lem:residual-numerator-nonzero}
If $f$ is nonzero, $q=u/D$ is in common-denominator form, and the displayed
evaluation $f(q)$ is nonzero, then $W_f(u,D)\ne0$ in $\Zi$.
\end{lemma}

\begin{proof}
The positive integers $D$ and $S_f$ embed as nonzero Gaussian rationals.
Their product with the nonzero displayed value $f(q)$ is therefore nonzero by
direct Gaussian multiplication.  The identity
\eqref{eq:homogeneous-evaluation} then prevents $W_f(u,D)$ from being zero.
\end{proof}

\leanref{RFAP.DenominatorHeight.residualNumerator_ne_zero}

\begin{lemma}[Gaussian-integral quadrance floor]
\label{lem:gaussian-integer-floor}
For every displayed $w\in\Zi$, if $w\ne0$, then $1\le\Norm w$.
\end{lemma}

\begin{proof}
Write $w=x+iy$ with $x,y\in\Z$.  Its quadrance $x^2+y^2$ is a natural
number.  If it were zero, the two nonnegative integer squares would both be
zero and hence $w=0$.  Positivity of the nonzero natural quadrance gives the
stated lower bound.
\end{proof}

\leanref{RFAP.ConstructiveGaussianHeight.GaussianInt.normNat_pos}

\begin{theorem}[exact common-denominator height identity]
\label{lem:height-lower}
If $f$ is nonzero and $q=u/D$ is a displayed common-denominator fraction,
then
\begin{equation}
\label{eq:height-identity}
 S_f^2H(q)^d\Norm{f(q)}=\Norm{W_f(u,D)}.
\end{equation}
If the displayed evaluation $f(q)$ is nonzero, then
\begin{equation}
\label{eq:height-lower}
                 1\le S_f^2H(q)^d\Norm{f(q)}.
\end{equation}
Consequently a nonzero residual is at least
$1/(S_f^2H(q)^d)$, with equality exactly when $\Norm{W_f(u,D)}=1$.
\end{theorem}

\begin{proof}
Taking quadrances in \eqref{eq:homogeneous-evaluation} gives the exact
calculation
\[
 \Norm{W_f(u,D)}=D^{2d}S_f^2\Norm{f(q)}
   =S_f^2H(q)^d\Norm{f(q)},
\]
which is \eqref{eq:height-identity}.  If $f(q)\ne0$, then
\cref{lem:residual-numerator-nonzero,lem:gaussian-integer-floor} give
$1\le\Norm{W_f(u,D)}$, proving \eqref{eq:height-lower}.  The equality
criterion follows from the exact identity.
\end{proof}

\leanref{RFAP.DenominatorHeight.scaled_height_mul_residual_eq_normNat;
RFAP.DenominatorHeight.scaled_height_mul_residual_eq_one_iff;
RFAP.DenominatorHeight.one_le_scaled_height_mul_residual;
RFAP.ApproximationPaper.scaledHeightResidual_eq_integralNorm;
RFAP.ApproximationPaper.one_le_scaled_height_mul_residual}

\begin{theorem}[displayed Gaussian-denominator height identity]
\label{thm:displayed-gaussian-height-lower}
Let $q=U/V$ be a displayed Gaussian fraction with $V\ne0$, and put
$h_V(q)=\Norm V$.  If $f$ is nonzero, then
\begin{equation}
\label{eq:displayed-gaussian-height-identity}
 S_f^2h_V(q)^d\Norm{f(q)}=\Norm{W_f(U,V)}.
\end{equation}
If $f(q)\ne0$, the common value is at least one, and hence
\begin{equation}
\label{eq:displayed-gaussian-height-lower}
                         \Norm{f(q)}\ge
                         \frac1{S_f^2h_V(q)^d}.
\end{equation}
\end{theorem}

\begin{proof}
Take quadrances in \eqref{eq:gaussian-homogeneous-evaluation}.  The rules
$\Norm{ab}=\Norm a\Norm b$ and $\Norm{V^d}=\Norm V^d$ give
\eqref{eq:displayed-gaussian-height-identity}.  If $f(q)\ne0$, the nonzero factors
$V^d$ and $S_f$ make $W_f(U,V)$ nonzero.  Its Gaussian-integral quadrance is
therefore at least one by \cref{lem:gaussian-integer-floor}.
\end{proof}

\leanref{RFAP.GaussianDenominatorHeight.scaled_height_mul_residual_eq_normNat;
RFAP.GaussianDenominatorHeight.one_le_scaled_height_mul_residual;
RFAP.ApproximationPaper.one_le_scaled_gaussianHeight_mul_residual}

\begin{definition}[canonical Gaussian reduction and intrinsic height]
\label{def:intrinsic-gaussian-height}
For a displayed Gaussian fraction $q=U/V$, run the extended Gaussian
Euclidean algorithm on $(U,V)$.  Its output includes a displayed common
factor $G$, quotients $U_0,V_0$, and B\'ezout coefficients satisfying
\[
 U=GU_0,\qquad V=GV_0,\qquad 1=A U_0+B V_0.
\]
Multiply $U_0,V_0$ by the uniquely selected unit in
$\{1,-1,i,-i\}$ for which the new denominator $V_*$ has
$\Ree(V_*)>0$ and $\Imm(V_*)\ge0$.  Write the resulting fraction as
\[
                 \operatorname{red}(q)=U_*/V_*
\]
and define its \emph{intrinsic Gaussian-denominator height} by
\[
                              h(q)=\Norm{V_*}.
\]
\end{definition}

The calculation preserves the exact value and gives
\[
       \Norm V=\Norm G\,h(q),\qquad 0<h(q)\le\Norm V.
\]
If two displayed fractions have equal exact values, cross multiplication and
their two displayed B\'ezout identities show that their reduced denominators
are unit associates.  The chosen denominator sector then forces the unit to
be $1$, after a finite classification of the four quadrance-one Gaussian
integers.  Their canonical numerators and denominators, and therefore their
intrinsic heights, are equal.  This proves presentation independence by an
executable normalization rather than by selecting a representative from a
quotient.

\leanref{RFAP.ReducedGaussianFraction.reduce;
RFAP.ReducedGaussianFraction.reduce_value;
RFAP.ReducedGaussianFraction.reduce_height_factor;
RFAP.ReducedGaussianFraction.reduce_height_le;
RFAP.ReducedGaussianFraction.canonicalize;
RFAP.ReducedGaussianFraction.canonicalize_value;
RFAP.ReducedGaussianFraction.canonicalize_height;
RFAP.ReducedGaussianFraction.reduceCanonical;
RFAP.ReducedGaussianFraction.reduceCanonical_value;
RFAP.ReducedGaussianFraction.reduceCanonical_coordinates_eq_of_value_eq;
RFAP.ReducedGaussianFraction.reduceCanonical_height_eq_of_value_eq;
RFAP.GaussianDenominatorHeight.intrinsicHeight;
RFAP.GaussianDenominatorHeight.intrinsicHeight_le_height;
RFAP.GaussianDenominatorHeight.intrinsicHeight_eq_of_value_eq;
RFAP.ApproximationPaper.canonicalReducedGaussianFraction;
RFAP.ApproximationPaper.intrinsicGaussianDenominatorHeight}

\begin{theorem}[intrinsic Gaussian-denominator height identity]
\label{thm:gaussian-height-lower}
Let $\operatorname{red}(q)=U_*/V_*$ be the canonical reduction from
\cref{def:intrinsic-gaussian-height}.  If $f$ is nonzero, then
\begin{equation}
\label{eq:gaussian-height-identity}
 S_f^2h(q)^d\Norm{f(q)}=\Norm{W_f(U_*,V_*)}.
\end{equation}
If $f(q)\ne0$, the common value is at least one, and hence
\begin{equation}
\label{eq:gaussian-height-lower}
                         \Norm{f(q)}\ge
                         \frac1{S_f^2h(q)^d}.
\end{equation}
\end{theorem}

\begin{proof}
The executable reduction preserves $q$ exactly.  Apply
\cref{thm:displayed-gaussian-height-lower} to its canonical displayed
coordinates $(U_*,V_*)$.  Their displayed denominator height is, by
definition, the intrinsic height $h(q)$.  The exact identity and the
nonzero-residual lower bound follow unchanged.
\end{proof}

\leanref{RFAP.GaussianDenominatorHeight.intrinsicResidualNumerator;
RFAP.GaussianDenominatorHeight.scaled_intrinsicHeight_mul_residual_eq_normNat;
RFAP.GaussianDenominatorHeight.one_le_scaled_intrinsicHeight_mul_residual;
RFAP.ApproximationPaper.scaledIntrinsicGaussianHeightResidual_eq_integralNorm;
RFAP.ApproximationPaper.one_le_scaled_intrinsicGaussianHeight_mul_residual}

\begin{theorem}[sharpness of the common-denominator obstruction]
\label{thm:height-lower-sharp}
For every $d\in\N$ and positive $S,D\in\N$, let
\[
                         f(Z)=\frac1S Z^d,
                 \qquad q=\frac1D.
\]
Then $S_f=S$, $\deg f=d$, $H(q)=D^2$, $f(q)\ne0$, and
\begin{equation}
\label{eq:height-lower-sharp}
                     S_f^2H(q)^d\Norm{f(q)}=1.
\end{equation}
Thus the unit constant, the square on $S_f$, and the degree-th power of
$H(q)$ are jointly sharp for a uniform theorem depending only on these
displayed parameters.
\end{theorem}

\begin{proof}
The only nonzero coefficient is the leading coordinate $1/S$, whose reduced
denominator is $S$; every zero coordinate has denominator $1$.  The finite
least-common-multiple calculation therefore gives $S_f=S$.
Finite monomial evaluation gives $f(1/D)=1/(SD^d)\ne0$.  Taking its
quadrance and substituting $H(q)=D^2$ makes every factor cancel, giving
\eqref{eq:height-lower-sharp}.
\end{proof}

\leanref{RFAP.DenominatorHeightSharp.sharpPolynomial;
RFAP.DenominatorHeightSharp.sharpFraction;
RFAP.DenominatorHeightSharp.sharpPolynomial_coefficientScale;
RFAP.DenominatorHeightSharp.sharpPolynomial_natDegree;
RFAP.DenominatorHeightSharp.scaled_height_mul_residual_eq_one;
RFAP.ApproximationPaper.denominatorHeightBound_isSharp}

The two displayed height identities have an exact relation.  A positive
common denominator $D$ is the Gaussian denominator $D+0i$, with the same
value and displayed height.  Conversely, multiplying a Gaussian fraction
$U/V$ by $\overline V/\overline V$ produces the common-denominator record
\[
                    \frac{U\overline V}{\Norm V},
\]
with the same value; displayed Gaussian height $h_V=\Norm V$ becomes
common-denominator height $h_V^2$.  These remain useful finite conversions.
Canonical Gaussian reduction is stronger: it cancels every displayed common
Gaussian factor, removes the remaining unit ambiguity, and supplies the
presentation-independent height used in
\cref{thm:gaussian-height-lower}.  Scaling an unreduced presentation can
weaken either displayed bound, while the intrinsic bound is unchanged.
Finally, \cref{thm:height-lower-sharp} shows that the common-denominator
identity cannot be uniformly improved from its displayed scale, height, and
degree alone.

\leanref{RFAP.GaussianDenominatorHeight.ofCommon_value;
RFAP.GaussianDenominatorHeight.ofCommon_height;
RFAP.GaussianDenominatorHeight.toCommon_value;
RFAP.GaussianDenominatorHeight.toCommon_height;
RFAP.GaussianDenominatorHeight.intrinsicHeight_le_height}

\section{A constructive sufficient height}
\label{sec:upper}

This section follows the sufficient-height branch of the glossary.  From the
displayed Newton centre and band radius it computes one grid denominator,
rounds the two centre coordinates, and proves strict box membership.  The
RFAP-1 residual certificate can then be applied to that displayed fraction.

\begin{definition}[Newton grid denominator and fraction]
\label{def:newton-grid}
For a displayed centre $c=x+iy\in\Qi$ and radius $r>0$, put
\[
 D_r=\left\lfloor\frac{1/2}{r}\right\rfloor+1,
 \qquad
 q_r=\frac{[D_rx]_{\rm n}+i[D_ry]_{\rm n}}{D_r}.
\]
The same positive denominator is used in both coordinates, so $q_r$ is a
displayed common-denominator fraction of exact height $H(q_r)=D_r^2$.
\end{definition}

\leanref{RFAP.HeightApproximation.gridDenominator;
RFAP.HeightApproximation.gridCoordinate;
RFAP.HeightApproximation.gridFraction;
RFAP.HeightApproximation.gridFraction_height}

\begin{lemma}[strict and least uniform grid mesh]
\label{lem:grid-mesh}
For every displayed $r>0$, the denominator of
\cref{def:newton-grid} satisfies
\[
                         \frac{1/2}{D_r}<r.
\]
Moreover, if $E>0$ is a natural number and $(1/2)/E<r$, then $D_r\le E$.
\end{lemma}

\begin{proof}
Put $t=(1/2)/r$.  The floor certificate gives
$\lfloor t\rfloor\le t<\lfloor t\rfloor+1=D_r$.  Clearing the positive
factors $r$ and $D_r$ in the strict inequality gives $(1/2)/D_r<r$.

Conversely, $(1/2)/E<r$ gives $t<E$.  Since $E$ is integral,
$\lfloor t\rfloor<E$, and hence $D_r=\lfloor t\rfloor+1\le E$.
\end{proof}

\leanref{RFAP.HeightApproximation.gridDenominator_pos;
RFAP.HeightApproximation.half_reciprocal_gridDenominator_lt;
RFAP.HeightApproximation.gridDenominator_minimal}

\begin{lemma}[nearest coordinate bracket]
\label{lem:grid-coordinate}
For every displayed $t\in\Q$,
\[
 t-\frac{1/2}{D_r}
 <\frac{[D_rt]_{\rm n}}{D_r}
 \le t+\frac{1/2}{D_r}.
\]
Consequently the grid fraction satisfies
\[
 x-r<\Ree q_r<x+r,
 \qquad
 y-r<\Imm q_r<y+r.
\]
\end{lemma}

\begin{proof}
The tie-up nearest-integer certificate is
\[
               -\frac12\le D_rt-[D_rt]_{\rm n}<\frac12.
\]
Rearrange it, divide by the positive $D_r$, and retain the resulting strict
lower and weak upper half-mesh endpoints.  The strict mesh inequality of
\cref{lem:grid-mesh} makes both final band inequalities strict, independently
in the real and imaginary coordinates.
\end{proof}

\leanref{RFAP.HurwitzNearest.nearestCoordinate_error_bounds;
RFAP.HeightApproximation.gridCoordinate_bounds}

\begin{lemma}[strict membership in the Newton band]
\label{lem:grid}
The fraction $q_r$ lies strictly inside every Newton band with centre $c$ and
radius $r$ returned by Proposition~\ref{prop:rfap-input}.
\end{lemma}

\begin{proof}
A Newton band places its left and bottom sides at least $r$ below the
corresponding centre coordinate and its right and top sides at least $r$
above it.  The strict coordinate bounds of \cref{lem:grid-coordinate}
therefore place both coordinates strictly between the certified sides.
\end{proof}

\leanref{
RFAP.HeightApproximation.gridFraction_strictly_inside}

\begin{lemma}[exact Newton-band geometric coefficient]
\label{lem:newton-band-half-sharp}
Every Newton band of radius $r$ has
\[
                         \diamq(B)\le32r^2.
\]
The band using all four outer allowed sides has equality.  In particular,
when $\eps=8r$, that band satisfies
\begin{equation}
\label{eq:newton-band-half-sharp}
                         \diamq(B)=\frac12\eps^2.
\end{equation}
Thus the factor $1/2$ retained by Proposition~\ref{prop:rfap-input} cannot be
improved from the Newton-band side constraints and radius budget alone.
\end{lemma}

\begin{proof}
Each side length is at most $4r$, so the sum of their squares is at most
$16r^2+16r^2=32r^2$.  Choosing the left, right, bottom, and top sides at
distance exactly $2r$ from the centre gives both side lengths $4r$ and hence
equality.  Substitution of $\eps=8r$ gives
$32r^2=(1/2)(8r)^2$.
\end{proof}

\leanref{RFAP.NewtonBand.NewtonBandBox.diagonalSq_le;
RFAP.NewtonBand.widestNewtonBand;
RFAP.NewtonBand.widestNewtonBand_diagonalSq;
RFAP.NewtonBand.widestNewtonBand_diagonalSq_eq_half_request}

\begin{theorem}[constructive height--residual calibration]
\label{thm:height-upper}
In the Newton extraction of Proposition~\ref{prop:rfap-input}, the algorithm
constructs the fraction $q_r$ of \cref{def:newton-grid} and certifies
\begin{equation}
\label{eq:height-upper}
 H(q_r)=\left(\left\lfloor\frac{1/2}{r}\right\rfloor+1\right)^2,
 \qquad
 \Norm{f(q_r)}\le2L_1(f,A)\eps^2.
\end{equation}
\end{theorem}

\begin{proof}
Proposition~\ref{prop:rfap-input} returns a Newton band with squared diagonal
at most $(1/2)\eps^2$.  Definition~\ref{def:newton-grid} gives the displayed
height identity, and Lemma~\ref{lem:grid} puts $q_r$ strictly inside the
returned band.  The RFAP-1 bound first gives
\[
                    \Norm{f(q_r)}
          \le4L_1(f,A)\diamq(B).
\]
Substituting $\diamq(B)\le(1/2)\eps^2$ gives the stated factor $2$.
\end{proof}

\leanref{RFAP.HeightApproximation.approximate;
RFAP.HeightApproximation.Output.height_eq;
RFAP.HeightApproximation.Output.residual_le;
RFAP.ApproximationPaper.approximateWithHeight}

\begin{corollary}[a displayed residual target]
\label{cor:height-target}
Let $\delta>0$ be displayed, put $C=2L_1(f,A)>0$, and compute
$n=\lceil C/\delta\rceil$.  Scan the positive integers from $1$ through the
certified endpoint $n+2$, let $m$ be the first for which $n<m^2$, and put
$\eps_\delta=1/m$.
It certifies
\begin{equation}
\label{eq:target-request}
                         C\eps_\delta^2\le\delta.
\end{equation}
Calling Theorem~\ref{thm:height-upper} with this square-root-scaled request
returns a displayed fraction of exact sufficient height and residual
quadrance at most $\delta$.
\end{corollary}

\begin{proof}
The endpoint $n+2$ works because $n<(n+2)^2$, so the finite scan terminates;
checking candidates in ascending order proves its minimality.  The ceiling
certificate and the selected strict square threshold give
\[
                    \frac C\delta\le n<m^2.
\]
Multiplying by the displayed positive $\delta/m^2$ yields
$C/m^2\le\delta$, which is \eqref{eq:target-request}.  The residual
conclusion follows from \eqref{eq:height-upper}.
\end{proof}

\leanref{RFAP.HeightApproximation.squareThreshold;
RFAP.HeightApproximation.squareThreshold_spec;
RFAP.HeightApproximation.squareThreshold_minimal;
RFAP.HeightApproximation.targetDenominator;
RFAP.HeightApproximation.targetEpsilon;
RFAP.HeightApproximation.coefficient_mul_targetEpsilon_sq_le;
RFAP.HeightApproximation.approximateTarget_residual_le;
RFAP.ApproximationPaper.approximateWithTarget_residual_le}

The lower and upper bounds answer different questions.  Theorem~\ref{lem:height-lower}
rules out too small a residual for a specified displayed height.  Theorem~\ref{thm:height-upper}
constructs one sufficient height from the actual band radius.  It does not
claim that this height is least for the particular centre or polynomial.
Lemma~\ref{lem:grid-mesh} proves only the precise uniform statement available:
$D_r$ is least among denominators whose half-mesh is strictly smaller than
$r$.  A particular centre may happen to admit a smaller denominator, and
the nearest grid point need not minimize the polynomial residual at its
height.

The Lean development also retains a checked exhaustive residual minimiser for
a window with both numerator coordinates and the positive common denominator
bounded.  This routine finite-list calculation is not part of the height
theorem and no later result depends on it; unlike Hurwitz--Lakein, it imposes
an independent numerator bound.

\leanref{RFAP.BestResidual.best_is_optimal}

\section{Protected rational-function sites}
\label{sec:quotients}

This section applies the height package to the glossary's protected quotient
residual.  The argument is symmetric: call the polynomial being made small
the \emph{source} $g$ and the other polynomial the \emph{opposite} $h$.  A
zero site of $P/Q$ uses $(g,h)=(P,Q)$; a pole site uses $(g,h)=(Q,P)$.  A
displayed B\'ezout certificate first keeps $h$ uniformly away from zero on a
protected source box.  Its lower margin and upper Horner majorant then
transfer the polynomial height bounds to the quotient.

\begin{definition}[quotient residual]
\label{def:quotient-residual}
For displayed polynomials $g,h\in\Qi[Z]$ and a displayed sample $q$ at which
$h(q)\ne0$, define
\begin{equation}
\label{eq:quotient-residual}
              \mathcal R_{g/h}(q)
                 =\Norm{g(q)/h(q)}
                 =\frac{\Norm{g(q)}}{\Norm{h(q)}}.
\end{equation}
The nonvanishing proof is part of every quantitative use of this expression.
\end{definition}

\leanref{RFAP.RationalResidualChordal.quotientResidualSq;
RFAP.Paper.quotientResidualSq}

\begin{definition}[B\'ezout protection data]
\label{def:bezout-protection}
Let $A$ be a rational ambient box.  A displayed coprimality certificate for
$g,h$ supplies polynomials $U,V\in\Qi[Z]$ with the coefficientwise identity
\begin{equation}
\label{eq:bezout-protection}
                            Ug+Vh=1.
\end{equation}
Let $\mathsf M_A(R)>0$ denote the finite Horner majorant for
$\Norm{R(q)}$ on $A$, and put
\begin{equation}
\label{eq:protection-margins}
 \theta_g=\frac1{16\mathsf M_A(U)},
 \qquad
 \eta_h=\frac1{16\mathsf M_A(V)},
 \qquad
 M_h=\mathsf M_A(h).
\end{equation}
The \emph{protected source state} is obtained by asking the source RAT for
the explicit positive radius
\begin{equation}
\label{eq:protection-radius}
 \rho_{g,h}=\min\left\{1,
   \frac{\theta_g}{8L_1(g,A)},
   \frac{\eta_h}{2L_1(h,A)}\right\}.
\end{equation}
A local quotient context records this certificate, the source RAT, its
current state, and the resulting protected state.
\end{definition}

\leanref{RFAP.PolynomialBounds.denseEvalNormMajorant;
RFAP.CoprimeBounds.leftThreshold;
RFAP.CoprimeBounds.rightMargin;
RFAP.CoprimeBounds.protectionRadius;
RFAP.CoprimeBounds.protectedState;
RFAP.RationalResidualChordal.LocalQuotientContext}

\begin{lemma}[the opposite polynomial is protected]
\label{lem:opposite-protected}
Every displayed $q$ in the protected source state satisfies
\begin{equation}
\label{eq:opposite-bounds}
                    0<\eta_h\le\Norm{h(q)}\le M_h.
\end{equation}
In particular, $h(q)\ne0$ and the quotient residual
\eqref{eq:quotient-residual} is defined there.
\end{lemma}

\begin{proof}
The source residual estimate and the first radius budget in
\eqref{eq:protection-radius} give
$\Norm{g(q)}\le\theta_g$.  Suppose the decidable rational comparison for the
lower bound instead returned $\Norm{h(q)}<\eta_h$.  Evaluating
\eqref{eq:bezout-protection}, using multiplicativity of quadrance, and
applying the Horner bounds would give
\[
 \Norm{U(q)g(q)}\le\mathsf M_A(U)\theta_g=\frac1{16},
 \qquad
 \Norm{V(q)h(q)}\le\mathsf M_A(V)\eta_h=\frac1{16}.
\]
The exact two-term quadrance estimate then calculates
\[
 1=\Norm 1
  =\Norm{U(q)g(q)+V(q)h(q)}
 \le2\left(\frac1{16}+\frac1{16}\right)=\frac14,
\]
which the rational order decision rejects.  Hence
$\eta_h\le\Norm{h(q)}$.  Positivity of $\eta_h$ gives nonvanishing.  The
upper bound is the defining Horner-majorant property for $h$ on $A$, applied
after box refinement.
\end{proof}

\leanref{RFAP.CoprimeBounds.leftThreshold_pos;
RFAP.CoprimeBounds.rightMargin_pos;
RFAP.CoprimeBounds.rightMargin_le_of_left_small;
RFAP.CoprimeBounds.protectedState_left_small;
RFAP.CoprimeBounds.protectedState_right_lower;
RFAP.RationalResidualChordal.LocalQuotientContext.oppositeLower;
RFAP.RationalResidualChordal.LocalQuotientContext.oppositeNonzero;
RFAP.RationalResidualChordal.LocalQuotientContext.oppositeUpper}

\begin{lemma}[protected quotient transfer]
\label{lem:protected-quotient-transfer}
For every displayed $q$ in a state refining the protected source state,
\begin{equation}
\label{eq:quotient-transfer}
 \Norm{g(q)}\le M_h\mathcal R_{g/h}(q),
 \qquad
 \mathcal R_{g/h}(q)\le\frac{\Norm{g(q)}}{\eta_h}.
\end{equation}
\end{lemma}

\begin{proof}
By \cref{lem:opposite-protected}, $h(q)\ne0$.  Taking quadrances in the
displayed identity
\[
                        (g(q)/h(q))h(q)=g(q)
\]
therefore gives
\[
       \mathcal R_{g/h}(q)\Norm{h(q)}=\Norm{g(q)}.
\]
Multiplication by the nonnegative quotient residual preserves the upper
bound $\Norm{h(q)}\le M_h$, proving the first inequality.  Multiplication by
the positive inverse of $\eta_h$ preserves the lower bound
$\eta_h\le\Norm{h(q)}$, proving the second.
\end{proof}

\leanref{RFAP.RationalResidualChordal.LocalQuotientContext.quotientResidualSq_mul_oppositeNorm;
RFAP.RationalResidualChordal.LocalQuotientContext.sourceResidual_le_quotientResidualSq_mul_majorant;
RFAP.RationalResidualChordal.LocalQuotientContext.quotientResidualSq_le_diagonal}

\begin{theorem}[protected quotient height obstruction]
\label{thm:quotient-height-lower}
Let $q$ be a displayed common-denominator fraction in a state refining the
protected $g$-state.  If $g(q)\ne0$, then
\begin{equation}
\label{eq:quotient-lower}
  1\le S_g^2H(q)^{\deg g}M_h\mathcal R_{g/h}(q).
\end{equation}
\end{theorem}

\begin{proof}
Theorem~\ref{lem:height-lower} gives the same lower bound with
$\Norm{g(q)}$ as its last factor.  Substitute the first inequality of
\eqref{eq:quotient-transfer} and multiply by the nonnegative displayed scale
and height factors.
\end{proof}

\leanref{RFAP.RationalDenominatorHeight.one_le_scaled_height_mul_quotientResidual;
RFAP.ApproximationPaper.one_le_scaled_height_mul_quotientResidual}

\begin{theorem}[constructive protected quotient height]
\label{thm:quotient-height-upper}
For every displayed $\eps>0$, Newton-grid extraction within the protected
$g$-state returns a displayed $q_r$ such that
\begin{equation}
\label{eq:quotient-upper}
 H(q_r)=\left(\left\lfloor\frac{1/2}{r}\right\rfloor+1\right)^2,
 \qquad
 \mathcal R_{g/h}(q_r)
 \le\frac{2L_1(g,A)\eps^2}{\eta_h}.
\end{equation}
\end{theorem}

\begin{proof}
Theorem~\ref{thm:height-upper} constructs $q_r$ in the protected source box,
with the displayed height and
$\Norm{g(q_r)}\le2L_1(g,A)\eps^2$.  Apply the second inequality of
\eqref{eq:quotient-transfer}.  The only division is by the explicitly
positive rational margin $\eta_h$.
\end{proof}

\leanref{RFAP.RationalDenominatorHeight.Output;
RFAP.RationalDenominatorHeight.approximate;
RFAP.ApproximationPaper.approximateQuotientWithHeight}

\begin{corollary}[a displayed protected quotient target]
\label{cor:quotient-height-target}
For a displayed target $\delta>0$, put
\[
                  C_{g/h}=\frac{2L_1(g,A)}{\eta_h}.
\]
Apply the bounded square-threshold scan of
\cref{cor:height-target} to $(C_{g/h},\delta)$.  Its reciprocal-integer
request $\eps_\delta$ makes Newton-grid extraction return a protected sample
with
\[
                       \mathcal R_{g/h}(q_r)\le\delta.
\]
\end{corollary}

\begin{proof}
The coefficient $C_{g/h}$ is positive because both the Lipschitz bound and
the protection margin are explicitly positive.  The scan proves
$C_{g/h}\eps_\delta^2\le\delta$, while
\eqref{eq:quotient-upper} gives
$\mathcal R_{g/h}(q_r)\le C_{g/h}\eps_\delta^2$.
\end{proof}

\leanref{RFAP.RationalDenominatorHeight.targetCoefficient;
RFAP.RationalDenominatorHeight.targetCoefficient_pos;
RFAP.RationalDenominatorHeight.approximateTarget;
RFAP.RationalDenominatorHeight.approximateTarget_quotientResidual_le;
RFAP.ApproximationPaper.approximateQuotientWithTarget_residual_le}

\begin{corollary}[zero and pole specialisations]
\label{thm:quotient-height}
For a protected zero track of $P/Q$, take $(g,h)=(P,Q)$ in
\cref{thm:quotient-height-lower,thm:quotient-height-upper,cor:quotient-height-target}.
For a protected pole track take
$(g,h)=(Q,P)$.  The shrinking quantity at a pole is therefore
$\mathcal R_{Q/P}$, with bounds
\[
 1\le S_Q^2H(q)^{\deg Q}M_P\mathcal R_{Q/P}(q),
 \qquad
 \mathcal R_{Q/P}(q_r)
 \le\frac{2L_1(Q,A)\eps^2}{\eta_P}.
\]
\end{corollary}

The opposite polynomial genuinely affects the rational-function estimate.
At a pole the numerator majorant $M_P$ occurs in the necessary lower bound,
while its B\'ezout protection margin $\eta_P$ occurs in the sufficient upper
bound and in the square-root target coefficient
$2L_1(Q,A)/\eta_P$ for the reciprocal residual.  If finite evaluation instead
finds $Q(q_r)=0$, it has found an exact displayed pole.  Otherwise an upper
bound on $\mathcal R_{Q/P}$ is equivalently a lower bound on the ordinary
quotient residual $\mathcal R_{P/Q}$.

\section{Finite Hurwitz expansion and continuants}
\label{sec:hurwitz-trace}

This section constructs the glossary's finite trace and generated-stage
algebra.  Deterministic coordinatewise rounding puts each residual in the
half-open fundamental square.  Clearing the division produces an integral
remainder of strictly smaller denominator quadrance, so recursion terminates
on a natural number.  Along the resulting trace, adjacent continuants retain
a unit determinant and an exact reconstruction of the original displayed
source.  These identities express every competitor in the adjacent basis and
factor its error by the current convergent error.

\begin{definition}[nearest-integer residual]
\label{def:nearest-residual}
For $x\in\Q$, $z=x+iy\in\Qi$, define
\[
 [x]_{\rm n}=\left\lfloor x+\frac12\right\rfloor,
 \qquad [z]_{\Zi}=[x]_{\rm n}+i[y]_{\rm n},
 \qquad \rho(z)=z-[z]_{\Zi}.
\]
The chosen tie rule sends a half-integer upward in each coordinate.  Thus the
associated fundamental square is the half-open region
\[
 \mathfrak F=\{u+iv:-\tfrac12\le u<\tfrac12,
                         -\tfrac12\le v<\tfrac12\}.
\]
\end{definition}

\leanref{RFAP.HurwitzNearest.nearestCoordinate;
RFAP.HurwitzNearest.nearest;
RFAP.HurwitzNearest.remainder;
RFAP.HurwitzNearest.InFundamentalSquare}

\begin{lemma}[fundamental-square remainder]
\label{lem:fundamental-square-remainder}
For every displayed $z\in\Qi$,
\[
              \rho(z)\in\mathfrak F,
              \qquad \Norm{\rho(z)}\le\frac12.
\]
\end{lemma}

\begin{proof}
Put $n=\lfloor x+1/2\rfloor$.  The defining floor inequalities are
\[
                     n\le x+\frac12<n+1.
\]
Subtracting $n+1/2$ gives
$-1/2\le x-n<1/2$.  Apply the same calculation to the imaginary coordinate.
Each coordinate square is at most $1/4$, so their sum is at most $1/2$.
The strict upper sides and weak lower sides are retained exactly; no symmetry
argument changes the tie convention.
\end{proof}

\leanref{RFAP.HurwitzNearest.nearestCoordinate_error_bounds;
RFAP.HurwitzNearest.remainder_mem;
RFAP.HurwitzNearest.remainder_normSq_le_half}

\begin{lemma}[nearest Gaussian integer is quadrance-minimising]
\label{lem:nearest-minimises}
For every $z\in\Qi$ and every displayed $c\in\Zi$,
\begin{equation}
\label{eq:nearest-minimises}
                    \Norm{z-[z]_{\Zi}}\le\Norm{z-c}.
\end{equation}
\end{lemma}

\begin{proof}
It is enough to prove the statement in one rational coordinate and add the
two inequalities.  Let $n=[x]_{\rm n}$, let $e=x-n$, and write an integral
competitor as $m=n+k$.  By \cref{lem:fundamental-square-remainder},
$-1/2\le e<1/2$, while
\[
                     (x-m)^2-(x-n)^2=k(k-2e).
\]
If $k=0$ this is zero.  If $k\ge1$, then $k-2e>0$; if $k\le-1$, both $k$ and
$k-2e$ are nonpositive.  The product is therefore nonnegative in every
decidable integer case.  Apply this to the real and imaginary coordinates
and add.
\end{proof}

\leanref{RFAP.HurwitzNearest.nearestCoordinate_sq_le;
RFAP.HurwitzNearest.remainder_normSq_le_competitor}

\begin{definition}[displayed Gaussian fraction and Hurwitz transition]
\label{def:hurwitz-transition}
A displayed Gaussian fraction is a record
\[
                         z=\frac ND,
              \qquad N,D\in\Zi,\quad D\ne0,
\]
whose denominator height is $\Norm D\in\N_{>0}$.  Put
\begin{equation}
\label{eq:hurwitz-division}
 a=[z]_{\Zi},\qquad R=N-aD,
 \qquad \rho=z-a.
\end{equation}
If $R=0$, the fraction is terminal.  If $R\ne0$, its next complete quotient
is the displayed fraction
\begin{equation}
\label{eq:hurwitz-next}
                         T(z)=\frac D R=\rho^{-1}.
\end{equation}
No cancellation is implicit in either fraction record.
\end{definition}

\leanref{RFAP.HurwitzRational.Fraction;
RFAP.HurwitzRational.Fraction.value;
RFAP.HurwitzRational.Fraction.height;
RFAP.HurwitzRational.digit;
RFAP.HurwitzRational.residual;
RFAP.HurwitzRational.remainder;
RFAP.HurwitzRational.next}

\begin{lemma}[one Hurwitz step decreases denominator quadrance]
\label{lem:hurwitz-decrease}
The data of \cref{def:hurwitz-transition} satisfy
\begin{equation}
\label{eq:hurwitz-remainder-norm}
                 \Norm R=\Norm\rho\,\Norm D.
\end{equation}
Consequently $\Norm R<\Norm D$.
\end{lemma}

\begin{proof}
Cross-multiplication in \eqref{eq:hurwitz-division} gives the displayed
Gaussian identity $R=\rho D$.  Multiplicativity of quadrance proves
\eqref{eq:hurwitz-remainder-norm}.  Since $D\ne0$, its integral quadrance is
positive.  Lemma~\ref{lem:fundamental-square-remainder} therefore gives
\[
       \Norm R\le\frac12\Norm D<\Norm D.
\]
The cast comparison between the two natural quadrances reflects this exact
rational calculation.
\end{proof}

\leanref{RFAP.HurwitzRational.digit_mul_denominator_add_remainder;
RFAP.HurwitzRational.remainder_toGaussian;
RFAP.HurwitzRational.remainder_norm_identity;
RFAP.HurwitzRational.remainder_norm_lt_denominator;
RFAP.HurwitzRational.next_value}

\begin{theorem}[terminating finite Hurwitz encoding]
\label{thm:finite-hurwitz-encoding}
Every displayed Gaussian fraction $z=N/D$ computes a nonempty finite digit
list $[a_0,\ldots,a_m]$ and a displayed Gaussian-integral pair $(p_m,q_m)$
with $q_m\ne0$ such that
\begin{equation}
\label{eq:hurwitz-reconstruct}
                         zq_m=p_m.
\end{equation}
The list is obtained solely by repeating \cref{def:hurwitz-transition} until
the integral remainder is zero.
\end{theorem}

\begin{proof}
Recurse on the natural denominator quadrance $\Norm D$.  A zero remainder
returns the one-element list $[a_0]$, and the integral division identity
$a_0D=N$ proves reconstruction.  A nonzero remainder calls the recursion on
$D/R$; this is structurally permitted by \cref{lem:hurwitz-decrease}.  If the
recursive tail represents $D/R$, prepend $a_0$ and use
$z=a_0+R/D$ to obtain the new continuant identity.  The tail numerator is
nonzero because its represented complete quotient is nonzero, so it becomes
the nonzero denominator of the prepended pair.  The recursion therefore
returns both the list and the certificate \eqref{eq:hurwitz-reconstruct}.
\end{proof}

\leanref{RFAP.HurwitzRational.Encoding;
RFAP.HurwitzRational.encodeResult;
RFAP.HurwitzRational.encode;
RFAP.HurwitzRational.encode_nonempty;
RFAP.HurwitzRational.encode_denominatorNonzero;
RFAP.HurwitzRational.encode_reconstruct}

The theorem should be read as a Euclidean algorithm, not as a truncation of
an already existing infinite expansion.  Its recursive argument is the
visible natural number $\Norm D$; for rational input the terminal digit and
the exact reconstruction are computed data.

\begin{definition}[chronological continuants]
\label{def:chronological-continuants}
For a displayed digit prefix $a_0,\ldots,a_n$, set
\[
 (p_{-1},q_{-1})=(1,0),\qquad (p_0,q_0)=(a_0,1),
\]
and advance chronologically by
\begin{equation}
\label{eq:chronological-continuants}
 (p_{j+1},q_{j+1})=
 a_{j+1}(p_j,q_j)+(p_{j-1},q_{j-1}).
\end{equation}
The finite-list continuant used by the encoder is the equivalent front
recursion with empty tail $(1,0)$ and
\[
           \operatorname{cont}(a::L)
            =\bigl(a p_L+q_L,p_L\bigr),
 \qquad \operatorname{cont}(L)=(p_L,q_L).
\]
\end{definition}

\leanref{RFAP.HurwitzRational.ConvergentPair;
RFAP.HurwitzRational.prepend;
RFAP.HurwitzRational.continuant;
RFAP.HurwitzTrace.advancePair}

\begin{definition}[determinant and adjacent basis]
\label{def:adjacent-basis}
For integral pairs $u=(p,q)$ and $v=(p',q')$, put
\[
                       \Delta(u,v)=pq'-p'q\in\Zi.
\]
They have a \emph{unit determinant} when the displayed alternative
$\Delta(u,v)=1$ or $\Delta(u,v)=-1$ is supplied.
\end{definition}

\leanref{RFAP.HurwitzApproximation.determinant;
RFAP.HurwitzApproximation.HasUnitDeterminant}

\begin{lemma}[adjacent continuants have unit determinant]
\label{lem:continuant-determinant}
Every two adjacent continuants produced by
\eqref{eq:chronological-continuants} have unit determinant.
\end{lemma}

\begin{proof}
The initial determinant is
$\Delta((a_0,1),(1,0))=-1$.  Direct coefficient calculation gives
\[
 \Delta\bigl(a_{j+1}(p_j,q_j)+(p_{j-1},q_{j-1}),
              (p_j,q_j)\bigr)
      =-\Delta((p_j,q_j),(p_{j-1},q_{j-1})).
\]
Thus each displayed step merely flips the explicit sign alternative.
\end{proof}

\leanref{RFAP.HurwitzApproximation.determinant_prepend;
RFAP.HurwitzApproximation.adjacent_hasUnitDeterminant;
RFAP.HurwitzTrace.determinant_advancePair;
RFAP.HurwitzTrace.unitDeterminant_advancePair}

\begin{definition}[Hurwitz stage and finite stage trace]
\label{def:full-hurwitz-stage}
At prefix $n$, let $x_n$ be the current complete quotient, let
$\rho_n=x_n-[x_n]_{\Zi}$, and let the adjacent pairs be
$c_{n-1}=(p_{n-1},q_{n-1})$ and $c_n=(p_n,q_n)$.  A \emph{Hurwitz stage}
stores these displayed objects, a unit-determinant certificate, and the exact
source reconstruction
\begin{equation}
\label{eq:stage-represents}
 z(q_n+\rho_nq_{n-1})=p_n+\rho_np_{n-1}.
\end{equation}
The first stage uses $x_0=z$, $c_{-1}=(1,0)$ and
$c_0=(a_0,1)$.  At a nonterminal stage, advance the complete quotient by
$x_{n+1}=\rho_n^{-1}$ and the pair by
\eqref{eq:chronological-continuants}.  Repeating this transition until the
integral remainder vanishes is the finite stage trace.
\end{definition}

\leanref{RFAP.HurwitzTrace.Stage;
RFAP.HurwitzTrace.firstStage;
RFAP.HurwitzTrace.advanceStage;
RFAP.HurwitzTrace.stagesFrom;
RFAP.HurwitzTrace.stages}

\begin{lemma}[stage reconstruction is preserved]
\label{lem:stage-reconstruction}
The first stage satisfies \eqref{eq:stage-represents}.  Every nonterminal
transition preserves that identity and the unit-determinant certificate.
The stage recursion terminates.
\end{lemma}

\begin{proof}
For the first stage, \eqref{eq:stage-represents} reduces to
$z=a_0+\rho_0$.  At a nonterminal step use
$x_{n+1}\rho_n=1$ and
$x_{n+1}=a_{n+1}+\rho_{n+1}$.  Multiplying the old reconstruction by
$x_{n+1}$ and regrouping gives
\[
 z\bigl((a_{n+1}q_n+q_{n-1})+\rho_{n+1}q_n\bigr)
 = (a_{n+1}p_n+p_{n-1})+\rho_{n+1}p_n,
\]
which is precisely the next reconstruction.  The determinant certificate is
preserved by \cref{lem:continuant-determinant}; termination is again the
strict natural decrease of \cref{lem:hurwitz-decrease}.
\end{proof}

\leanref{RFAP.HurwitzTrace.firstStage;
RFAP.HurwitzTrace.advanceStage;
RFAP.HurwitzTrace.advanceStage_quotient;
RFAP.HurwitzRational.remainder_norm_lt_denominator}

\begin{lemma}[explicit adjacent-basis decomposition]
\label{lem:adjacent-basis-decomposition}
Let $c_n,c_{n-1}$ be the adjacent pairs of a stage, let $c=(p,q)$ be any
displayed Gaussian-integral competitor, and write
$\delta=\Delta(c_n,c_{n-1})\in\{1,-1\}$.  Define
\begin{equation}
\label{eq:competitor-coefficients}
 s=\delta\Delta(c,c_{n-1}),
 \qquad t=\delta\Delta(c_n,c).
\end{equation}
Then, coefficientwise in $\Zi^2$,
\begin{equation}
\label{eq:competitor-decomposition}
                         c=sc_n+tc_{n-1}.
\end{equation}
Consequently its division-free error satisfies
\begin{equation}
\label{eq:competitor-error-decomposition}
            E_z(c)=sE_z(c_n)+tE_z(c_{n-1}).
\end{equation}
\end{lemma}

\begin{proof}
The explicit unit alternative gives $\delta^2=1$.  Expanding the two
determinants in \eqref{eq:competitor-coefficients}, first in the numerator
coordinate and then in the denominator coordinate, cancels the mixed terms
and leaves $\delta^2p=p$ and $\delta^2q=q$.  This proves
\eqref{eq:competitor-decomposition} coefficientwise.  Substitute those two
identities into $E_z(p,q)=zq-p$ and distribute to obtain
\eqref{eq:competitor-error-decomposition}.
\end{proof}

\leanref{RFAP.HurwitzApproximation.determinant_mul_self_eq_one;
RFAP.HurwitzApproximation.currentCoefficient;
RFAP.HurwitzApproximation.previousCoefficient;
RFAP.HurwitzApproximation.competitor_numerator_decomposition;
RFAP.HurwitzApproximation.competitor_denominator_decomposition;
RFAP.HurwitzApproximation.competitor_error_decomposition}

\begin{lemma}[stage error recurrence and Lakein factor]
\label{lem:lakein-factor}
At a nonterminal stage,
\begin{equation}
\label{eq:stage-error-recurrence}
 E_z(c_n)+\rho_nE_z(c_{n-1})=0,
 \qquad
 E_z(c_{n-1})=-x_{n+1}E_z(c_n).
\end{equation}
For the coefficients $s,t$ of
\cref{lem:adjacent-basis-decomposition},
\begin{equation}
\label{eq:lakein-factor}
                         E_z(c)=(s-tx_{n+1})E_z(c_n).
\end{equation}
Thus the second-kind comparison follows whenever
$1\le\Norm{s-tx_{n+1}}$.
\end{lemma}

\begin{proof}
Subtract the two sides of the stored reconstruction
\eqref{eq:stage-represents}; the result is the first identity in
\eqref{eq:stage-error-recurrence}.  Multiply it by
$x_{n+1}=\rho_n^{-1}$ to obtain the second.  Substitution into
\eqref{eq:competitor-error-decomposition} gives
\eqref{eq:lakein-factor}.  Taking quadrances, using multiplicativity, and
multiplying the factor bound by the nonnegative $\Norm{E_z(c_n)}$ proves the
comparison.
\end{proof}

\leanref{RFAP.HurwitzTrace.Stage.error_add_residual_mul_previous_eq_zero;
RFAP.HurwitzTrace.Stage.previous_error_eq_neg_next_mul_current;
RFAP.HurwitzTrace.Stage.competitor_error_eq_factor_mul_current;
RFAP.HurwitzTrace.Stage.atLeastAsGoodAs_of_one_le_factor_normSq}

The algebra has now isolated the only difficult quantity: the Lakein factor
$s-tx_{n+1}$.  The denominator budget constrains the integral coefficients
$s,t$, while the Hurwitz trace constrains $x_{n+1}$ and the backward
denominator ratio.  The next three sections turn those constraints into the
quadrance-one exterior inequality without introducing a limiting value or an
infinite tail.

\section{Fundamental-square and succession geometry}
\label{sec:succession-geometry}

This section supplies the finite geometric constraints used by the last
factor cases.  The fundamental square gives exterior-disc laws for every
noninitial complete quotient.  A two-part stage invariant then proves strict
growth of adjacent denominator quadrances.  Normalising those displayed
denominators produces the backward ratio $r_n=q_{n-1}/q_n$, whose transition
is an exact finite fractional transformation.  Finally, cleared exterior
laws exclude a displayed list of small adjacent digit pairs; these are the
forbidden successors consumed by the Hermitian-form induction in the next
section.

\begin{definition}[Gaussian units, diagonals, and corner turn]
\label{def:small-gaussian-classes}
The four Gaussian units and four diagonal Gaussian integers are the finite
classes
\[
 \mathcal U=\{1,-1,i,-i\},
 \qquad
 \mathcal D=\{1+i,1-i,-1+i,-1-i\}.
\]
Thus $u\in\mathcal U$ exactly in the cases explicitly listed by the
predicate \emph{Gaussian unit}, and $c\in\mathcal D$ exactly in the cases
listed by \emph{diagonal}.  For a diagonal $c$, define its corner turn by
\[
                              c^\sharp=-\overline c.
\]
The transform permutes $\mathcal D$ and is the direction in which a failed
corner configuration propagates one step backward.
\end{definition}

\leanref{RFAP.HurwitzGeometry.imaginaryUnit;
RFAP.HurwitzGeometry.IsGaussianUnit;
RFAP.HurwitzGeometry.IsDiagonal;
RFAP.HurwitzGeometry.conjugate;
RFAP.HurwitzGeometry.cornerTurn;
RFAP.HurwitzGeometry.cornerTurn_isDiagonal}

\begin{lemma}[small Gaussian-integral quadrances]
\label{lem:small-gaussian-quadrances}
For $w\in\Zi$:
\begin{enumerate}[label=(\roman*)]
\item if $\Norm w<4$, then $\Norm w\in\{0,1,2\}$;
\item if $\Norm w<6$, then $\Norm w\in\{0,1,2,4,5\}$;
\item quadrance one means $w\in\mathcal U$, quadrance two means
$w\in\mathcal D$, quadrance four means $w\in\{2,-2,2i,-2i\}$, and
quadrance five gives the eight values $\pm1\pm2i,\ \pm2\pm i$.
\end{enumerate}
\end{lemma}

\begin{proof}
Write $w=a+ib$.  The assumptions bound each integral square $a^2,b^2$ by
the same displayed natural threshold.  For quadrance below four, each
coordinate is therefore $-1,0$, or $1$; for quadrance below six it is one of
$-2,-1,0,1,2$.  The finite Cartesian products give (i) and (ii), and sorting
the same checked pairs by $a^2+b^2$ gives (iii).
\end{proof}

\leanref{RFAP.HurwitzGeometry.normNat_eq_zero_or_one_or_two_of_lt_four;
RFAP.HurwitzGeometry.normNat_eq_zero_or_one_or_two_or_four_or_five_of_lt_six;
RFAP.HurwitzGeometry.isGaussianUnit_of_normNat_eq_one;
RFAP.HurwitzGeometry.isDiagonal_of_normNat_eq_two;
RFAP.HurwitzGeometry.axisDouble_cases_of_normNat_eq_four;
RFAP.HurwitzGeometry.normFive_cases_of_normNat_eq_five}

\begin{lemma}[complete-quotient exterior laws]
\label{lem:complete-quotient-exterior}
If a displayed fraction $x$ is nonterminal and $x^+=T(x)=\rho(x)^{-1}$,
then
\begin{equation}
\label{eq:complete-quotient-size}
                     2\le\Norm{x^+},
 \qquad
                     2\le\Norm{[x^+]_{\Zi}}.
\end{equation}
For every $u\in\mathcal U$,
\begin{equation}
\label{eq:unit-exterior}
                       1\le\Norm{x^+-u}.
\end{equation}
At the positive real unit the half-open tie convention strengthens this to
$1<\Norm{x^+-1}$.
\end{lemma}

\begin{proof}
Since $x^+\rho(x)=1$, multiplicativity and
$\Norm{\rho(x)}\le1/2$ give $2\le\Norm{x^+}$.  For a unit $u$, the
fundamental-square inequalities give
\[
                 \Norm{\rho(x)}
          \le \Norm{u\rho(x)-1}.
\]
For $u=1$ the upper real inequality
$\operatorname{Re}\rho(x)<1/2$ makes this strict.  Multiplying by the
positive $\Norm{x^+}$ and using $x^+\rho(x)=1$ gives
\eqref{eq:unit-exterior} and its strict case.  If the nearest digit of $x^+$
had quadrance zero or one, subtracting that digit would simultaneously have
quadrance at most $1/2$ by the fundamental square and at least one by the
exterior law (or $x^+$ itself would contradict
\eqref{eq:complete-quotient-size}).  Hence the digit quadrance is at least
two.
\end{proof}

\leanref{RFAP.HurwitzGeometry.residual_normSq_le_unit_mul_sub_one;
RFAP.HurwitzGeometry.one_le_next_sub_unit_normSq;
RFAP.HurwitzGeometry.next_normSq_ge_two;
RFAP.HurwitzGeometry.two_le_next_digit_normNat;
RFAP.HurwitzGeometry.one_lt_next_sub_one_normSq}

\begin{definition}[stage geometry]
\label{def:stage-geometry}
For a Hurwitz stage with adjacent denominators $q_{n-1},q_n$ and next
complete quotient $x_{n+1}$, its \emph{stage geometry} consists of:
\begin{enumerate}[label=(\roman*)]
\item strict denominator growth
\begin{equation}
\label{eq:denominator-growth}
                         \Norm{q_{n-1}}<\Norm{q_n};
\end{equation}
\item for every diagonal $c\in\mathcal D$, exclusion of the joint corner
failure
\begin{equation}
\label{eq:corner-failure}
 \Norm{cq_n+q_{n-1}}\le\Norm{q_n}
 \quad\hbox{and}\quad
 \Norm{x_{n+1}-c}<1.
\end{equation}
\end{enumerate}
The second clause is a finite natural-extension state: it retains precisely
the information needed to prove the first clause again after advancing.
\end{definition}

\leanref{RFAP.HurwitzGeometry.StageGeometry}

\begin{lemma}[generated stages have stage geometry]
\label{lem:generated-stage-geometry}
The first Hurwitz stage has stage geometry, and every nonterminal advance
preserves it.  Consequently every generated stage satisfies
\eqref{eq:denominator-growth}.
\end{lemma}

\begin{proof}
At the first stage the denominators are $0,1$, so growth is immediate, and a
diagonal multiplier has quadrance two and cannot make
$\Norm{c\cdot1+0}\le\Norm1$.

Assume the invariant at stage $n$.  If the next denominator
\[
                         q_{n+1}=a_{n+1}q_n+q_{n-1}
\]
were no longer than $q_n$, the small-coefficient geometry forced by
\eqref{eq:denominator-growth} classifies $a_{n+1}$ as a diagonal.  Its own
fundamental-square residual is strictly inside the unit disc, so these two
facts are exactly the corner failure \eqref{eq:corner-failure} at stage $n$.
This proves the new strict growth.

For corner safety at stage $n+1$, expand a proposed failure using the two
successive recurrences and the identities
$x_{n+1}=a_{n+1}+x_{n+2}^{-1}$ and
$c^\sharp=-\overline c$.  The cleared quadrance comparison transfers the
failure to stage $n$ at $c^\sharp$.  Since
$c^\sharp\in\mathcal D$, the old corner-safety clause rejects it.  Induction
on the displayed generation certificate proves the result for every stage.
\end{proof}

\leanref{RFAP.HurwitzGeometry.next_digit_isDiagonal_of_combination_norm_le;
RFAP.HurwitzGeometry.corner_failure_transfers_backward;
RFAP.HurwitzGeometry.firstStage_geometry;
RFAP.HurwitzGeometry.advanceStage_geometry;
RFAP.HurwitzGeometry.geometry_of_generated;
RFAP.HurwitzGeometry.stage_denominatorGrowth;
RFAP.HurwitzGeometry.generated_of_mem_stages}

This is the point at which the familiar picture of expanding convergent
denominators becomes an actual finite proof.  The fundamental-square bound
alone does not make the chronological recurrence monotone term by term;
the corner-safety memory is what rules out the diagonal cancellation that
could otherwise spoil the next step.

\begin{definition}[backward denominator ratio]
\label{def:backward-ratio}
At a generated stage define the displayed Gaussian rational
\begin{equation}
\label{eq:backward-ratio}
                          r_n=\frac{q_{n-1}}{q_n}.
\end{equation}
The denominator is nonzero by strict growth.  On advancing,
\begin{equation}
\label{eq:backward-ratio-recurrence}
                         r_{n+1}=(a_{n+1}+r_n)^{-1}.
\end{equation}
\end{definition}

\leanref{RFAP.HurwitzGeometry.Stage.denominatorRatio;
RFAP.HurwitzGeometry.current_denominator_ne_zero_of_generated;
RFAP.HurwitzGeometry.denominatorRatio_mul_current;
RFAP.HurwitzGeometry.advanceStage_denominatorRatio}

\begin{lemma}[the backward ratio lies in the open unit disc]
\label{lem:backward-ratio-unit-disc}
Every generated stage satisfies
\begin{equation}
\label{eq:backward-ratio-unit-disc}
                              \Norm{r_n}<1.
\end{equation}
\end{lemma}

\begin{proof}
Multiply \eqref{eq:backward-ratio} by the nonzero displayed $q_n$ and take
quadrances:
\[
                       \Norm{r_n}\Norm{q_n}=\Norm{q_{n-1}}.
\]
The strict denominator growth of
\cref{lem:generated-stage-geometry} and positivity of $\Norm{q_n}$ give
\eqref{eq:backward-ratio-unit-disc} by exact rational division.
\end{proof}

\leanref{RFAP.HurwitzGeometry.denominatorRatio_normSq_lt_one_of_generated}

\begin{lemma}[cleared successor exterior law]
\label{lem:successor-exterior}
Let $a_{j}$ and $a_{j+1}$ be two displayed noninitial adjacent digits, and
let $x_{j+1}$ be the complete quotient beginning with $a_{j+1}$.  For every
$u\in\mathcal U$,
\begin{equation}
\label{eq:successor-exterior}
 \Norm{x_{j+1}}
 \le
 \Norm{(a_j-u)x_{j+1}+1}.
\end{equation}
For $u=1$ the inequality is strict.
\end{lemma}

\begin{proof}
The preceding complete quotient is
$x_j=a_j+x_{j+1}^{-1}$.  Apply
\eqref{eq:unit-exterior} to $x_j-u$ and multiply by the positive
$\Norm{x_{j+1}}$.  The cleared identity
\[
                 (x_j-u)x_{j+1}=(a_j-u)x_{j+1}+1
\]
gives \eqref{eq:successor-exterior}.  The strict exterior law at $u=1$
preserves strictness under the same positive multiplication.
\end{proof}

\leanref{RFAP.HurwitzGeometry.successor_exterior_cleared;
RFAP.HurwitzGeometry.successor_strict_exterior_cleared}

\begin{definition}[forbidden small successors]
\label{def:forbidden-successors}
The relation $\mathcal F(a_j,a_{j+1})$ consists of the following twenty-eight
ordered pairs.  To keep the finite table readable, the left column fixes the
following digit $a_{j+1}$ and the right column lists its forbidden preceding
digits $a_j$.
\[
\begin{array}{c|l}
a_{j+1} & \multicolumn{1}{c}{\text{forbidden }a_j}\\ \hline
-2      & 1-i,\ 1+i\\
-1-i    & -2i,\ 1-2i,\ 1-i,\ 2-i,\ 2\\
-1+i    & 2i,\ 1+i,\ 1+2i,\ 2,\ 2+i\\
-2i     & -1-i,\ 1-i\\
 2i     & -1+i,\ 1+i\\
 1-i    & -2-i,\ -2,\ -1-2i,\ -1-i,\ -2i\\
 1+i    & -2,\ -2+i,\ -1+i,\ -1+2i,\ 2i\\
 2      & -1-i,\ -1+i
\end{array}
\]
It is an inductive finite relation with one constructor for each displayed
entry, not a predicate obtained from an unbounded search.
\end{definition}

\leanref{RFAP.HurwitzGeometry.ForbiddenSmallSuccessor}

\begin{lemma}[forbidden successors cannot occur]
\label{lem:no-forbidden-successor}
No two adjacent noninitial digits of a generated finite Hurwitz trace belong
to $\mathcal F$.
\end{lemma}

\begin{proof}
Write the following quotient as $x_{j+1}=a_{j+1}+\sigma$ with
$\sigma\in\mathfrak F$.  For a proposed table entry, choose the displayed
unit $u$ attached to that entry and subtract the two sides of
\eqref{eq:successor-exterior}.  Expansion gives one of a small collection of
signed gaps, including
\[
 \Norm\sigma-1,\qquad -1\pm2\operatorname{Re}\sigma,
 \qquad -1\pm2\operatorname{Im}\sigma,
 \qquad -3\pm2\operatorname{Re}\sigma,
 \qquad -3\pm2\operatorname{Im}\sigma.
\]
The fundamental-square inequalities make the appropriate gap negative.
Where a lower boundary could make a gap zero, the strict $u=1$ exterior law
or the preceding inversion supplies the needed strict inequality.  Thus the
proposed pair makes the cleared exterior inequality point the wrong way.
Elimination on the twenty-eight constructors performs these exact polynomial
calculations one by one.
\end{proof}

\leanref{RFAP.HurwitzGeometry.successorGap;
RFAP.HurwitzGeometry.successor_signs_of_one_add_i;
RFAP.HurwitzGeometry.residual_eq_next_value_inv;
RFAP.HurwitzGeometry.not_forbiddenSmallSuccessor}

The table is less elegant than the phrase ``and similarly in the other
quadrants'', but it records what that phrase omits: which boundary is strict,
which unit exterior disc is used, and which ordered digit pair is actually
excluded.  Section~\ref{sec:hermitian-invariant} uses precisely this finite
relation to close the second pullback step.

\section{Hermitian root forms and the exceptional invariant}
\label{sec:hermitian-invariant}

This section proves the glossary's root-form invariant.  The remaining
comparison between quadrance-five and diagonal coefficients is encoded by eight integral
Hermitian forms evaluated at the backward ratio.  Pulling a form through the
exact recurrence $r_{n+1}=(a_{n+1}+r_n)^{-1}$ is a coefficient calculation.
Most pullbacks have an immediate sum-of-squares certificate; twenty-four
first-step exceptions require one preceding digit.  Their second pullbacks
are either certified nonnegative or demand one of the forbidden successors
already excluded in \cref{lem:no-forbidden-successor}.  A two-stage induction
then establishes all eight inequalities at every generated stage.

\begin{definition}[integral Hermitian form]
\label{def:integral-hermitian-form}
For displayed integers $A,B,C,D$ define
\begin{equation}
\label{eq:integral-hermitian-form}
 \mathcal H_{A,B,C,D}(z)
   =A+B\Norm z+2\bigl(C\operatorname{Re}z+D\operatorname{Im}z\bigr).
\end{equation}
The associated integral coefficient record is written $F=(A,B,C,D)$, with
evaluation $F[z]=\mathcal H_{A,B,C,D}(z)$ and determinant
\begin{equation}
\label{eq:hermitian-determinant}
                             \det F=AB-(C^2+D^2).
\end{equation}
\end{definition}

\leanref{RFAP.HurwitzGeometry.finiteHermitianForm;
RFAP.HurwitzGeometry.IntegralHermitianForm;
RFAP.HurwitzGeometry.IntegralHermitianForm.evaluate;
RFAP.HurwitzGeometry.IntegralHermitianForm.determinant}

\begin{lemma}[finite unit-disc certificate]
\label{lem:hermitian-slack-certificate}
Let $F=(A,B,C,D)$ and suppose $\Norm z\le1$.  If
\begin{equation}
\label{eq:slack-one-certificate}
                  B+1>0,
 \qquad (A-1)(B+1)-(C^2+D^2)\ge0,
\end{equation}
then $F[z]\ge0$.
\end{lemma}

\begin{proof}
Put $z=x+iy$.  Direct expansion gives the displayed identity
\begin{align*}
 (B+1)F[z]
   ={}&((B+1)x+C)^2+((B+1)y+D)^2\\
     &+(B+1)(1-\Norm z)
       +\bigl((A-1)(B+1)-(C^2+D^2)\bigr).
\end{align*}
Every term on the right is nonnegative by the supplied rational
inequalities.  Exact division by the positive $B+1$ proves the result.  The
formal lemma also permits a displayed nonnegative slack other than one; the
form above is the only instance needed by the finite classification.
\end{proof}

\leanref{RFAP.HurwitzGeometry.finiteHermitianForm_nonnegative_of_certificate;
RFAP.HurwitzGeometry.IntegralHermitianForm.evaluate_nonnegative_of_slack_one}

\begin{definition}[integral pullback]
\label{def:hermitian-pullback}
For $F=(A,B,C,D)$ and $a=u+iv\in\Zi$, its pullback through
$z\mapsto(a+z)^{-1}$ is the integral form
\begin{equation}
\label{eq:hermitian-pullback-coefficients}
 F\mathbin{\triangleleft}a=
 \bigl(A\Norm a+B+2(Cu-Dv),\ A,\ Au+C,\ Av-D\bigr).
\end{equation}
\end{definition}

\leanref{RFAP.HurwitzGeometry.IntegralHermitianForm.pullback}

\begin{lemma}[pullback identity and determinant]
\label{lem:hermitian-pullback}
If $a+z\ne0$, then
\begin{equation}
\label{eq:hermitian-pullback}
 F[(a+z)^{-1}]\Norm{a+z}=(F\mathbin{\triangleleft}a)[z].
\end{equation}
Moreover
\begin{equation}
\label{eq:pullback-determinant}
                 \det(F\mathbin{\triangleleft}a)=\det F.
\end{equation}
\end{lemma}

\begin{proof}
The inverse has coordinates
\[
 (a+z)^{-1}
 =\frac{\operatorname{Re}(a+z)-i\operatorname{Im}(a+z)}{\Norm{a+z}}.
\]
Substitute these two displayed rational coordinates into
\eqref{eq:integral-hermitian-form} and multiply by $\Norm{a+z}$.  Collecting
the constant, radial, real-cross, and imaginary-cross coefficients gives
exactly \eqref{eq:hermitian-pullback-coefficients}.  Expanding
\eqref{eq:hermitian-determinant} for those four new coefficients cancels all
$u,v$ terms and leaves $AB-(C^2+D^2)$.
\end{proof}

\leanref{RFAP.HurwitzGeometry.finiteHermitianForm_inverse_add_mul_normSq;
RFAP.HurwitzGeometry.IntegralHermitianForm.evaluate_pullback_mul_normSq;
RFAP.HurwitzGeometry.IntegralHermitianForm.pullback_determinant}

\begin{lemma}[transport through one generated advance]
\label{lem:hermitian-transport}
Let $r_n$ be the backward ratio at a generated nonterminal stage and let
$a_{n+1}$ be its next digit.  If
\[
               (F\mathbin{\triangleleft}a_{n+1})[r_n]\ge0,
\]
then $F[r_{n+1}]\ge0$.
\end{lemma}

\begin{proof}
The affine denominator $a_{n+1}+r_n$ cannot vanish.  Otherwise
$r_n=-a_{n+1}$, but \cref{lem:backward-ratio-unit-disc} gives
$\Norm{r_n}<1$, whereas \eqref{eq:complete-quotient-size} gives
$\Norm{a_{n+1}}\ge2$.  Now use
$r_{n+1}=(a_{n+1}+r_n)^{-1}$ in
\eqref{eq:hermitian-pullback}.  Its other factor
$\Norm{a_{n+1}+r_n}$ is positive, so the pullback's nonnegativity divides
through to $F[r_{n+1}]\ge0$.
\end{proof}

\leanref{RFAP.HurwitzGeometry.denominatorRatio_step_sum_ne_zero;
RFAP.HurwitzGeometry.IntegralHermitianForm.evaluate_advance_nonnegative_of_pullback}

\begin{definition}[Lakein root forms]
\label{def:lakein-root-forms}
Let
\[
 \mathcal X=\{(-3,-1),(-3,1),(-1,-3),(-1,3),
                 (1,-3),(1,3),(3,-1),(3,1)\}.
\]
For $(C,D)\in\mathcal X$, the corresponding \emph{Lakein root form} is
\begin{equation}
\label{eq:lakein-root-form}
             L_{C,D}=(4,2,C,D),
 \qquad
             L_{C,D}[z]
               =4+2\Norm z+2(C\operatorname{Re}z+D\operatorname{Im}z).
\end{equation}
All eight have determinant $-2$.  The \emph{root-form condition} at $z$
means $L_{C,D}[z]\ge0$ for every displayed pair $(C,D)\in\mathcal X$.
\end{definition}

\leanref{RFAP.HurwitzGeometry.IsLakeinRootCross;
RFAP.HurwitzGeometry.lakeinRootForm;
RFAP.HurwitzGeometry.LakeinRootFormsNonnegative}

\begin{lemma}[root forms are the quadrance-five exterior test]
\label{lem:root-form-exterior}
Let $s,t\in\Zi$ satisfy $\Norm s=5$ and $\Norm t=2$.  Then the two integral
cross terms
\[
 C=s_{\rm re}t_{\rm re}+s_{\rm im}t_{\rm im},
 \qquad
 D=-s_{\rm re}t_{\rm im}+s_{\rm im}t_{\rm re}
\]
form one of the eight pairs in $\mathcal X$, and
\begin{equation}
\label{eq:root-form-factor}
                       \Norm{s+tr}=1+L_{C,D}[r].
\end{equation}
Consequently the root-form condition at $r$ implies
$1\le\Norm{s+tr}$ for every such $s,t$.
\end{lemma}

\begin{proof}
Lemma~\ref{lem:small-gaussian-quadrances} lists the eight possible $s$ and
four possible $t$.  Substitution gives exactly the eight cross pairs in
$\mathcal X$.  Expanding the real and imaginary coordinates of $s+tr$, and
then using $\Norm s=5$ and $\Norm t=2$, gives
\[
 \Norm{s+tr}=5+2\Norm r
      +2(C\operatorname{Re}r+D\operatorname{Im}r),
\]
which is \eqref{eq:root-form-factor}.  Add the certified nonnegative root
form to one.
\end{proof}

\leanref{RFAP.HurwitzGeometry.one_le_normSq_normFive_add_diagonal_mul_of_rootForms}

\begin{definition}[first exceptional pullbacks]
\label{def:first-exceptional-hermitian}
For a noninitial digit $a$ with $\Norm a\ge2$, pull back each of the eight
root forms once.  A constant-coefficient gap certifies every result except
the following twenty-four displayed form--digit pairs:
\[
\begin{array}{c|l}
a & \multicolumn{1}{c}{\text{exceptional }F=L_{C,D}\mathbin{\triangleleft}a}
\\ \hline
 1+i &(6,4,1,5),\ (2,4,1,3),\ (2,4,3,1),\ (6,4,5,1)\\
 1-i &(2,4,1,-3),\ (6,4,1,-5),\ (2,4,3,-1),\ (6,4,5,-1)\\
 2   &(6,4,5,1),\ (6,4,5,-1)\\
-1-i &(6,4,-5,-1),\ (2,4,-3,-1),\ (2,4,-1,-3),\ (6,4,-1,-5)\\
-2i  &(6,4,-1,-5),\ (6,4,1,-5)\\
-1+i &(6,4,-5,1),\ (2,4,-3,1),\ (6,4,-1,5),\ (2,4,-1,3)\\
 2i  &(6,4,-1,5),\ (6,4,1,5)\\
-2   &(6,4,-5,1),\ (6,4,-5,-1)
\end{array}
\]
The inductive relation \emph{first exceptional Hermitian} records these
pairs with one constructor per row entry.
\end{definition}

\leanref{RFAP.HurwitzGeometry.FirstExceptionalHermitian}

\begin{lemma}[first-pullback classification]
\label{lem:first-pullback-classification}
For every root form $L$ and every Gaussian integer $a$ with
$\Norm a\ge2$, either
\begin{equation}
\label{eq:first-pullback-gap}
                 7\le (L\mathbin{\triangleleft}a)_{\rm constant},
\end{equation}
or $(L\mathbin{\triangleleft}a,a)$ is one of the twenty-four exceptions in
\cref{def:first-exceptional-hermitian}.  In the first case the pulled-back
form is nonnegative throughout the closed unit disc.
\end{lemma}

\begin{proof}
Every root form has constant four and determinant $-2$.  The pullback
determinant is still $-2$.  If its new constant is at least seven, the
determinant identity rearranges the second inequality of
\eqref{eq:slack-one-certificate}; its radial coefficient is four, so the
first inequality is positive.  Lemma~\ref{lem:hermitian-slack-certificate}
applies.

It remains to classify failure of \eqref{eq:first-pullback-gap}.  The
explicit pullback constant has the lower estimate
\[
 (A-1)\Norm a+B-(C^2+D^2)
       \le (F\mathbin{\triangleleft}a)_{\rm constant}.
\]
For $F=L$ this settles $\Norm a\ge5$.  Under
$2\le\Norm a<5$, Lemma~\ref{lem:small-gaussian-quadrances} leaves only the
four diagonal and four axial-double digits.  Substituting those eight digits
into the eight root forms gives either the gap or exactly the displayed
twenty-four exceptions.
\end{proof}

\leanref{RFAP.HurwitzGeometry.IntegralHermitianForm.pullback_constant_lower;
RFAP.HurwitzGeometry.IntegralHermitianForm.pullback_evaluate_nonnegative_of_constant_gap;
RFAP.HurwitzGeometry.FirstExceptionalHermitian}

\begin{lemma}[exceptional second pullback]
\label{lem:exceptional-second-pullback}
Let $(F,a)$ be one of the twenty-four first-step exceptions.  Let $b$ be the
preceding noninitial digit, so $\Norm b\ge2$, and suppose
$\mathcal F(b,a)$ does not hold.  If $\Norm r\le1$ and the eight root forms
are nonnegative at $r$, then
\begin{equation}
\label{eq:exceptional-second-pullback}
                       (F\mathbin{\triangleleft}b)[r]\ge0.
\end{equation}
\end{lemma}

\begin{proof}
If $\Norm b\ge11$, the pullback-constant lower estimate gives the same
slack-one certificate directly.  If $2\le\Norm b<11$, the two coordinate
squares show that $b$ belongs to an explicit list of thirty-two Gaussian
integers with coordinates between $-3$ and $3$.  Substitution of those
digits into the twenty-four exceptional forms classifies each of the 768
displayed cases as one of:
\begin{enumerate}[label=(\roman*)]
\item a forbidden pair $\mathcal F(b,a)$;
\item one of the original eight root forms;
\item one of the eight augmented forms
$(8,6,\pm7,\pm1)$ or $(8,6,\pm1,\pm7)$; or
\item a form satisfying the slack-one certificate
\eqref{eq:slack-one-certificate}.
\end{enumerate}
Case (i) is excluded by hypothesis, (ii) by the root-form condition, and
(iv) by \cref{lem:hermitian-slack-certificate}.  Each augmented form is one
root form plus $4\Norm{r+u}$ for a displayed Gaussian unit $u$, so case (iii)
is also nonnegative.  These finite alternatives prove
\eqref{eq:exceptional-second-pullback}.
\end{proof}

\leanref{RFAP.HurwitzGeometry.FirstExceptionalHermitian.pullback_nonnegative}

The 768 substitutions are intentionally kernel-checked coefficient
calculations, but reproducing them as a paper table would hide rather than
clarify the argument.  The meaningful finite interface is the four-way
classification above: only its first branch depends on actual Hurwitz
succession, and Section~\ref{sec:succession-geometry} has already ruled that
branch out.

\begin{theorem}[root-form invariant]
\label{thm:root-form-invariant}
For every generated Hurwitz stage and every $(C,D)\in\mathcal X$,
\begin{equation}
\label{eq:root-form-invariant}
                              L_{C,D}[r_n]\ge0.
\end{equation}
\end{theorem}

\begin{proof}
The first ratio is $r_0=0$, where all eight forms have value four.  At the
first advance, apply the first-pullback classification at $r_0$.  A regular
pullback is nonnegative on the unit disc; an exceptional pullback is still
nonnegative at zero because its displayed constant is two or six.

For the induction step, retain the root-form condition both at the current
stage and after its next possible advance.  To prove it two advances later,
pull a target root form back through the following digit.  A regular first
pullback is nonnegative at the middle ratio by
\cref{lem:hermitian-slack-certificate}.  For an exceptional first pullback,
pull once more through the preceding digit.  The previous ratio has
quadrance below one by \cref{lem:backward-ratio-unit-disc}, its root forms are
nonnegative by the induction hypothesis, and the digit pair is not forbidden
by \cref{lem:no-forbidden-successor}; hence
\cref{lem:exceptional-second-pullback} applies.  Transport forward once and
then once more using \cref{lem:hermitian-transport}.  The two-component
invariant aligns this two-step calculation with ordinary induction on the
one-step generated-stage certificate.
\end{proof}

\leanref{RFAP.HurwitzGeometry.lakeinRootFormsNonnegative_of_generated}

\begin{corollary}[quadrance-five denominator safety]
\label{cor:norm-five-denominator-safe}
At every generated stage, for all $s,t\in\Zi$ with
$\Norm s=5$ and $\Norm t=2$,
\begin{equation}
\label{eq:norm-five-denominator-safe}
        \Norm{q_n}\le\Norm{s q_n+t q_{n-1}}.
\end{equation}
\end{corollary}

\begin{proof}
By \cref{thm:root-form-invariant,lem:root-form-exterior},
$1\le\Norm{s+t r_n}$.  Multiply by $\Norm{q_n}>0$ and use
\[
                    (s+t r_n)q_n=sq_n+tq_{n-1}.
\]
The resulting quadrance inequality is
\eqref{eq:norm-five-denominator-safe}.
\end{proof}

\leanref{RFAP.HurwitzGeometry.NormFiveDenominatorSafe;
RFAP.HurwitzGeometry.normFiveDenominatorSafe_of_rootForms;
RFAP.HurwitzGeometry.normFiveDenominatorSafe_of_generated}

This corollary is the sole fact exported from the long exceptional
classification into the final Lakein comparison.  Its form is deliberately
denominator-only: if a competitor coefficient pair ever has the exceptional
quadrances $(5,2)$, the required smaller-denominator hypothesis contradicts
\eqref{eq:norm-five-denominator-safe} before any appeal to the source value.

\section{The finite Lakein inequality}
\label{sec:lakein-proof}

This section assembles the glossary's final comparison.  The exact complete
denominator makes the current error a reciprocal in quadrance.  A large
previous-basis coefficient is then excluded by a determinant estimate.  If
that test fails, the coefficient has quadrance below four and hence is zero,
a unit, or diagonal.  The first two cases use elementary exterior laws; the
diagonal case leaves a current coefficient of quadrance below six, where the
root-form invariant removes the sole exceptional value five.  The resulting
Lakein-factor bound completes every nonterminal stage, while a terminal stage
has exact zero error.

\begin{definition}[complete denominator]
\label{def:complete-denominator}
At a nonterminal stage put
\begin{equation}
\label{eq:complete-denominator}
                    \Lambda_n=x_{n+1}q_n+q_{n-1}\in\Qi.
\end{equation}
This is an unnormalised displayed Gaussian rational, not a new continuant.
The stage reconstruction gives
\begin{equation}
\label{eq:complete-denominator-reconstruction}
                    z\Lambda_n=x_{n+1}p_n+p_{n-1}.
\end{equation}
\end{definition}

\leanref{RFAP.HurwitzTrace.Stage.completeDenominator;
RFAP.HurwitzTrace.Stage.scaled_represents}

\begin{lemma}[exact current-error product]
\label{lem:current-error-product}
At every nonterminal stage,
\begin{equation}
\label{eq:current-error-product}
 E_z(c_n)\Lambda_n=-\Delta(c_n,c_{n-1}),
 \qquad
 \Norm{E_z(c_n)}\Norm{\Lambda_n}=1.
\end{equation}
Both factors in the second identity are positive.
\end{lemma}

\begin{proof}
Multiply $E_z(c_n)=zq_n-p_n$ by \eqref{eq:complete-denominator}, substitute
\eqref{eq:complete-denominator-reconstruction}, and expand:
\begin{align*}
 (zq_n-p_n)\Lambda_n
   &=q_n(x_{n+1}p_n+p_{n-1})
       -p_n(x_{n+1}q_n+q_{n-1})\\
   &=q_np_{n-1}-p_nq_{n-1}
    =-\Delta(c_n,c_{n-1}).
\end{align*}
The adjacent determinant is explicitly $1$ or $-1$, so taking quadrances
gives the second identity.  A product equal to one cannot contain a
zero-quadrance factor.
\end{proof}

\leanref{RFAP.HurwitzTrace.Stage.current_error_mul_completeDenominator;
RFAP.HurwitzTrace.Stage.current_error_norm_product_eq_one;
RFAP.HurwitzTrace.Stage.current_error_normSq_pos;
RFAP.HurwitzTrace.Stage.completeDenominator_normSq_pos}

\begin{lemma}[complete-denominator growth]
\label{lem:complete-denominator-growth}
Every generated nonterminal stage satisfies
\begin{equation}
\label{eq:complete-denominator-growth}
                         \Norm{q_n}<\Norm{\Lambda_n}.
\end{equation}
\end{lemma}

\begin{proof}
At the first stage, $\Lambda_0=x_1$ and
$\Norm{x_1}\ge2>1=\Norm{q_0}$.  For the induction step, the exact recurrence
\[
                  \Lambda_{n+1}=x_{n+2}\Lambda_n
\]
transfers the preceding strict inequality after multiplication by the
positive $\Norm{x_{n+2}}$.  Suppose nevertheless that
$\Norm{\Lambda_{n+1}}\le\Norm{q_{n+1}}$.  The resulting two rational
quadrance estimates force $\Norm{x_{n+2}}<3$.  The four unit exterior laws
then classify its nearest digit as diagonal.  Expanding the corresponding
fundamental-square lune shows
$\Norm{q_{n+2}}\le\Norm{q_{n+1}}$, contradicting the strict denominator
growth of \cref{lem:generated-stage-geometry}.  Hence the proposed failure is
rejected at each generated step.
\end{proof}

\leanref{RFAP.HurwitzGeometry.advanceStage_completeDenominator;
RFAP.HurwitzGeometry.quotient_normSq_lt_three_of_complete_failure;
RFAP.HurwitzGeometry.next_digit_isDiagonal_of_normSq_lt_three;
RFAP.HurwitzGeometry.diagonal_combination_norm_le_of_complete_failure;
RFAP.HurwitzGeometry.completeDenominatorGrowth_of_generated}

\begin{lemma}[large previous coefficient]
\label{lem:large-previous-coefficient}
Let a displayed competitor $c=(p,q)$ have $q\ne0$ and
$\Norm q<\Norm{q_n}$.  Let $s,t$ be its adjacent-basis coefficients from
\eqref{eq:competitor-coefficients}.  If
\begin{equation}
\label{eq:large-previous-coefficient}
       4\Norm{q_n}<\Norm t\,\Norm{\Lambda_n},
\end{equation}
then
\[
                         \Norm{E_z(c_n)}\le\Norm{E_z(c)}.
\]
\end{lemma}

\begin{proof}
For any current pair and competitor, direct expansion of the cross
determinant gives
\[
 \Delta(c_n,c)=q_nE_z(c)-qE_z(c_n).
\]
The two-term quadrance estimate therefore gives
\begin{equation}
\label{eq:determinant-error-bound}
 \Norm{\Delta(c_n,c)}
 \le2\bigl(\Norm{q_n}\Norm{E_z(c)}
           +\Norm q\Norm{E_z(c_n)}\bigr).
\end{equation}
If the claimed comparison failed strictly, the smaller-error and
smaller-denominator hypotheses would make the right side strictly less than
$4\Norm{q_n}\Norm{E_z(c_n)}$.  Since the adjacent determinant is a unit,
$\Norm t=\Norm{\Delta(c_n,c)}$.  Multiply the strict estimate by the positive
$\Norm{\Lambda_n}$ and use
\eqref{eq:current-error-product}; the result is
\[
                 \Norm t\,\Norm{\Lambda_n}<4\Norm{q_n},
\]
contrary to \eqref{eq:large-previous-coefficient}.
\end{proof}

\leanref{RFAP.HurwitzTrace.determinant_eq_denominator_mul_error_sub;
RFAP.HurwitzTrace.determinant_normSq_le;
RFAP.HurwitzTrace.Stage.previousCoefficient_normSq;
RFAP.HurwitzTrace.Stage.atLeastAsGoodAs_of_large_previousCoefficient}

\begin{lemma}[small previous-coefficient reduction]
\label{lem:small-previous-coefficient}
Under the hypotheses of \cref{lem:large-previous-coefficient}, if
\eqref{eq:large-previous-coefficient} does not hold, then
\begin{equation}
\label{eq:previous-coefficient-small}
                              \Norm t<4.
\end{equation}
Hence $\Norm t$ is $0$, $1$, or $2$.
\end{lemma}

\begin{proof}
If $\Norm t\ge4$, the complete-denominator growth
\eqref{eq:complete-denominator-growth} gives
\[
 4\Norm{q_n}<4\Norm{\Lambda_n}
             \le\Norm t\,\Norm{\Lambda_n},
\]
which is exactly the supposedly false large-coefficient test.  Therefore
$\Norm t<4$, and the three displayed possibilities follow from
\cref{lem:small-gaussian-quadrances}.
\end{proof}

\leanref{RFAP.HurwitzGeometry.previousCoefficient_normNat_lt_four_of_not_large;
RFAP.HurwitzGeometry.normNat_eq_zero_or_one_or_two_of_lt_four}

\begin{lemma}[small-coefficient Lakein factors]
\label{lem:small-coefficient-factors}
Assume
\begin{equation}
\label{eq:competitor-denominator-combination}
        q=sq_n+tq_{n-1},\qquad q\ne0,\qquad \Norm q<\Norm{q_n}.
\end{equation}
Then every case in \eqref{eq:previous-coefficient-small} either is
impossible or satisfies
\begin{equation}
\label{eq:small-factor-exterior}
                         1\le\Norm{s-tx_{n+1}}.
\end{equation}
\end{lemma}

\begin{proof}
If $\Norm t=0$, then $t=0$ and $q=sq_n$.  A nonzero $s$ makes this
denominator at least as large as $q_n$, while $s=0$ contradicts $q\ne0$.

Suppose $\Norm t=1$.  Conjugating by the unit $t$ rewrites
$s=t c$ with $c=\overline t s$, without changing quadrances.  The strict
denominator inequality becomes
$\Norm{cq_n+q_{n-1}}<\Norm{q_n}$ and forces $\Norm c<4$.  Thus $c$ is zero,
a unit, or diagonal.  The first two cases follow from
$\Norm{x_{n+1}}\ge2$ and the unit exterior laws.  In the diagonal case,
failure of \eqref{eq:small-factor-exterior} would put $x_{n+1}$ strictly
within the corresponding corner disc while the displayed denominator
combination does not grow.  This is exactly the corner failure excluded by
stage geometry.

Finally suppose $\Norm t=2$, so $t$ is diagonal.  The strict denominator
combination and \eqref{eq:denominator-growth} give the finite estimate
$\Norm s<6$.  Lemma~\ref{lem:small-gaussian-quadrances} leaves
$\Norm s=0,1,2,4$, or $5$.  The first three cases follow respectively from
the complete-quotient size, a unit exterior estimate, and association of two
diagonals.  Quadrance four factors through a displayed diagonal corner and
is again excluded by corner safety if its Lakein factor enters the unit
disc.  In the remaining quadrance-five case,
\cref{cor:norm-five-denominator-safe} says
\[
                 \Norm{q_n}\le\Norm{sq_n+tq_{n-1}}=\Norm q,
\]
contradicting \eqref{eq:competitor-denominator-combination}.  All cases are
therefore impossible or satisfy \eqref{eq:small-factor-exterior}.
\end{proof}

\leanref{RFAP.HurwitzGeometry.coefficient_normNat_lt_four_of_combination_norm_le;
RFAP.HurwitzGeometry.one_le_factor_normSq_of_unit_previousCoefficient;
RFAP.HurwitzGeometry.currentCoefficient_normNat_lt_six_of_diagonal_previous;
RFAP.HurwitzGeometry.one_le_factor_normSq_of_unit_and_diagonal;
RFAP.HurwitzGeometry.one_le_factor_normSq_of_diagonal_coefficients;
RFAP.HurwitzGeometry.one_le_factor_normSq_of_four_and_diagonal;
RFAP.HurwitzGeometry.one_le_factor_normSq_of_diagonal_previousCoefficient_of_not_five}

\begin{theorem}[finite Lakein comparison at a nonterminal stage]
\label{thm:finite-lakein-nonterminal}
Let $c_n=(p_n,q_n)$ be the current pair of a generated nonterminal stage.
For every displayed $c=(p,q)\in\Zi^2$ satisfying
\[
                         q\ne0,
             \qquad      \Norm q<\Norm{q_n},
\]
one has
\begin{equation}
\label{eq:finite-lakein-nonterminal}
                         \Norm{zq_n-p_n}\le\Norm{zq-p}.
\end{equation}
\end{theorem}

\begin{proof}
Use the explicit adjacent-basis coefficients $s,t$ of
\cref{lem:adjacent-basis-decomposition}; its denominator identity is
\eqref{eq:competitor-denominator-combination}.  Decide the rational
large-coefficient comparison \eqref{eq:large-previous-coefficient}.  If it
holds, \cref{lem:large-previous-coefficient} proves the result directly.  If
it does not, \cref{lem:small-previous-coefficient} and
\cref{lem:small-coefficient-factors} prove
$1\le\Norm{s-tx_{n+1}}$.  The error factorisation
\eqref{eq:lakein-factor} then gives
\[
 \Norm{zq-p}
   =\Norm{s-tx_{n+1}}\Norm{zq_n-p_n}
   \ge\Norm{zq_n-p_n}.
\]
\end{proof}

\leanref{RFAP.HurwitzGeometry.Stage.atLeastAsGoodAs_of_generated_of_normFiveDenominatorSafe;
RFAP.HurwitzGeometry.Stage.atLeastAsGoodAs_of_generated_nonterminal}

\begin{theorem}[finite Hurwitz--Lakein inequality]
\label{thm:finite-lakein-full}
At every generated stage of the finite Hurwitz trace, every displayed
Gaussian-integral competitor $(p,q)$ with
$q\ne0$ and $\Norm q<\Norm{q_n}$ satisfies
\begin{equation}
\label{eq:finite-lakein-full}
                         \Norm{zq_n-p_n}\le\Norm{zq-p}.
\end{equation}
\end{theorem}

\begin{proof}
If the stage is nonterminal, apply
\cref{thm:finite-lakein-nonterminal}.  If its integral remainder is zero,
then \eqref{eq:hurwitz-remainder-norm} and $D\ne0$ give $\rho_n=0$.  The
stage error recurrence \eqref{eq:stage-error-recurrence} reduces to
$E_z(c_n)=0$.  Its quadrance is therefore zero and hence no larger than the
nonnegative quadrance of any displayed competitor error.
\end{proof}

\leanref{RFAP.HurwitzGeometry.Stage.atLeastAsGoodAs_of_generated;
RFAP.ApproximationPaper.hurwitzConvergent_atLeastAsGoodAs}

The theorem compares division-free residuals.  If both displayed
denominators are nonzero, dividing by their quadrances converts each to the
ordinary squared approximation error, but that changes the comparison and
is not the theorem proved here.  The second-kind form is exactly what the
adjacent determinant and smaller-denominator budget control; it also remains
meaningful without choosing reduced representatives for the competitor.

\section{Finite all-best approximations}
\label{sec:all-best}

This section proves the stronger conclusion available for a displayed
rational source.  It first identifies the glossary's best pairs with minimal
source-lattice vectors.  The exact source representation then supplies a
global denominator cutoff, while nearest-integer geometry and denominator
clearance supply a global numerator cutoff.  Those bounds turn the universal
definition into an executable finite test.  The same rectangle then decides
whether two best pairs occupy consecutive denominator levels and returns the
complete next level above a supplied pair.  These are rectangular exhaustive
procedures rather than Chevallier's local structural minimal-vector update;
the proof nevertheless recovers his local empty-cylinder characterization and
the sharp index-one/index-two determinant dichotomy.  It also gives the
centered-coset classification in the index-two case, the resulting
four-candidate error-ratio bound, a finite twelve-value normal form for
two putative successive index-two transitions, and the exact cylinder
calculation excluding every value in that normal form.  Finally it uses these
structural results to generate the next level from bounded coefficients of one
adjacent pair, with a proved no-skipped-level specification and candidate-count
bound.  A local denominator-and-error comparison pass then returns exactly the
next occupied level without scanning the global source rectangle.

\begin{lemma}[best approximation equals cylinder minimality]
\label{lem:best-equals-minimal}
For a displayed source $z$ and pair $c=(p,q)$ with $q\ne0$, the two-tier
conditions \eqref{eq:best-strict}--\eqref{eq:best-weak} hold if and only if
the vector $(E_z(c),q)$ is minimal in its closed error--denominator cylinder.
\end{lemma}

\begin{proof}
First restrict attention to competitors with nonzero denominator.  Suppose
$c$ is best and $d=(a,b)$ lies in its closed cylinder.  If
$\Norm b<\Norm q$, the cylinder inequality
$\mathcal E_z(d)\le\mathcal E_z(c)$ contradicts the strict best comparison.
Hence $\Norm b=\Norm q$; the weak best comparison supplies the reverse error
inequality, so both cylinder coordinates are equal.

Conversely, suppose the cylinder condition holds.  For a lower-height
competitor, failure of the strict error comparison would put that competitor
in the closed cylinder, whose boundary conclusion would falsely equate the
two denominator quadrances.  For a no-higher competitor, failure of the weak
comparison would again put it in the cylinder and would falsely equate its
strictly smaller error with the reference error.  Thus both tiers follow.

It remains to admit the lattice vectors with $b=0$.  At the fixed denominator
$q$, take the coordinatewise nearest Gaussian integer to $zq$.  The weak
equal-height comparison and the fundamental-square estimate give
\begin{equation}
\label{eq:best-error-half}
                         \mathcal E_z(c)\le\frac12.
\end{equation}
A nonzero pair $(a,0)$ has error $-a$, whose integral quadrance is at least
one.  It therefore cannot lie in the cylinder of $c$.  This proves the full
source-lattice statement without postulating a carrier beyond the displayed
coefficient pairs.
\end{proof}

\leanref{RFAP.BestApproximation.isBest_iff_isDenominatorMinimal;
RFAP.BestApproximation.IsBest.errorQuadrance_le_half;
RFAP.BestApproximation.isBest_iff_isGaussLatticeMinimal}

\begin{lemma}[global finite rectangle]
\label{lem:all-best-rectangle}
Let $z=N/D$ be the supplied Gaussian fraction and put
\begin{equation}
\label{eq:all-best-bounds}
                 B=\Norm D,
        \qquad   P=2(\Norm N\,B+B).
\end{equation}
Every best pair $(p,q)$ satisfies
\begin{equation}
\label{eq:all-best-norm-bounds}
                         \Norm q\le B,
             \qquad     \Norm p\le P.
\end{equation}
Consequently each real and imaginary component of $q$ has magnitude at most
$B$, and each component of $p$ has magnitude at most $P$.
\end{lemma}

\begin{proof}
The exact pair $(N,D)$ is a displayed competitor with zero error.  If
$B<\Norm q$, the strict lower-height clause for a best pair would say
$\mathcal E_z(p,q)<0$, contrary to nonnegativity.  Hence $\Norm q\le B$.

Equation~\eqref{eq:best-error-half} supplies the other ingredient.  Define
the cleared Gaussian-integer error
\begin{equation}
\label{eq:all-best-cleared-error}
                         W=Nq-pD\in\Zi.
\end{equation}
Direct multiplication by $D$ gives
\begin{equation}
\label{eq:all-best-cleared-norm}
             \Norm W=\mathcal E_z(p,q)\Norm D\le B.
\end{equation}
Since $pD=Nq-W$, multiplicativity and the two-term quadrance estimate give
\[
 \Norm p\,B
   =\Norm{pD}
   \le2\bigl(\Norm N\Norm q+\Norm W\bigr)
   \le2(\Norm N B+B)=P.
\]
The nonzero integral $D$ has $B\ge1$, so $\Norm p\le\Norm p,B\le P$.
Finally, for a Gaussian integer $u+iv$, each of $|u|$ and $|v|$ is bounded
by $u^2+v^2$; applying this elementary displayed inequality gives the
component bounds.
\end{proof}

\leanref{RFAP.BestApproximation.IsBest.denominatorQuadrance_le_source;
RFAP.BestApproximation.clearedError_norm_identity;
RFAP.BestApproximation.clearedError_normNat_le_denominator;
RFAP.BestApproximation.numeratorQuadrance_le_bound;
RFAP.BestApproximation.IsBest.mem_candidates}

\begin{lemma}[finite descent and the mixed-cylinder criterion]
\label{lem:mixed-cylinder-criterion}
Let $d$ be a displayed pair with nonzero denominator,
$\Norm{q_d}\le B$, and $\mathcal E_z(d)\le1/2$.  There is a best pair $c$
such that
\begin{equation}
\label{eq:finite-best-descent}
       \Norm{q_c}\le\Norm{q_d},
       \qquad \mathcal E_z(c)\le\mathcal E_z(d).
\end{equation}
Consequently, two best pairs $c_-,c_+$ with
$\Norm{q_-}<\Norm{q_+}$ occupy consecutive best levels if and only if no
nonzero displayed pair lies in the mixed open cylinder
\begin{equation}
\label{eq:mixed-open-cylinder}
 \mathring{C}_z(c_-,c_+)=
 \{d:\mathcal E_z(d)<\mathcal E_z(c_-),\
          \Norm{q_d}<\Norm{q_+}\}.
\end{equation}
In particular, for Gaussian integers $g,h$ not both zero,
\begin{equation}
\label{eq:linear-combination-exterior}
 \mathcal E_z(c_-)\le\mathcal E_z(gc_-+hc_+)
 \quad\text{or}\quad
 \Norm{q_+}\le\Norm{gq_-+hq_+}.
\end{equation}
\end{lemma}

\begin{proof}
Restrict the global candidate rectangle to nonzero denominators of quadrance
at most $\Norm{q_d}$.  The supplied pair $d$ belongs to that finite list.
Recursively select the lexicographically least pair, first by error quadrance
and then, among exact error ties, by denominator quadrance.  The selected pair
has no greater error or denominator than $d$.

To verify that it is best, consider a competitor below the selected height.
If its error exceeds $1/2$, the selected error is already strictly smaller.
Otherwise the same global bounds put the competitor in the checked list.
Lexicographic minimality gives the weak error inequality, and equality would
force the selected denominator to be no greater, contradicting strict lower
height.  The no-higher weak clause follows by the same two cases without the
last contradiction.  This proves \eqref{eq:finite-best-descent} by an
executable finite selection.

Now suppose $c_-,c_+$ are consecutive and that a nonzero $d$ lies in
\eqref{eq:mixed-open-cylinder}.  A vertical pair has integral error quadrance
at least one, whereas $\mathcal E_z(c_-)\le1/2$, so $q_d\ne0$.  Bestness of
$c_-$ forces $\Norm{q_-}<\Norm{q_d}$.  Apply the descent just proved to obtain
a best $c$ no higher and no worse than $d$.  Bestness of $c_-$ again gives
$\Norm{q_-}<\Norm{q_c}$, while
$\Norm{q_c}\le\Norm{q_d}<\Norm{q_+}$, contradicting consecutiveness.

Conversely, a best pair at an intermediate level has strictly smaller error
than $c_-$ and therefore lies in the mixed open cylinder.  This proves the
equivalence.  Finally, coefficient independence of $c_-,c_+$ makes
$gc_-+hc_+$ nonzero whenever $(g,h)\ne(0,0)$; applying cylinder emptiness and
deciding which strict inequality fails gives
\eqref{eq:linear-combination-exterior}.
\end{proof}

\leanref{RFAP.BestApproximation.RankLE;
RFAP.BestApproximation.selectRankMinimum;
RFAP.BestApproximation.exists_best_below_of_small;
RFAP.BestApproximation.InOpenMixedCylinder;
RFAP.BestApproximation.MixedCylinderEmpty;
RFAP.BestApproximation.areConsecutiveBest_iff_mixedCylinderEmpty;
RFAP.BestApproximation.AreConsecutiveBest.linearCombination_pairNonzero;
RFAP.BestApproximation.AreConsecutiveBest.linearCombination_outside_mixedCylinder;
RFAP.ApproximationPaper.consecutiveBest_iff_mixedCylinderEmpty;
RFAP.ApproximationPaper.exists_bestBelow_of_small}

\begin{theorem}[executable finite all-best selector]
\label{thm:all-best-selector}
For a displayed source $z=N/D$, form the component rectangles from
\cref{lem:all-best-rectangle}.  For every candidate reference pair in that
rectangle, check \eqref{eq:best-error-half} and both comparisons
\eqref{eq:best-strict}--\eqref{eq:best-weak} against every pair in the same
rectangle.  Retain precisely the successful references.  The resulting
finite list $\operatorname{AllBest}(z)$ satisfies
\begin{equation}
\label{eq:all-best-membership}
 c\in\operatorname{AllBest}(z)
       \quad\Longleftrightarrow\quad
 c\text{ is a best approximation to }z.
\end{equation}
Its unfiltered candidate list has exactly
$16(B+1)^2(P+1)^2$ entries, and the output has no more entries than this.
\end{theorem}

\begin{proof}
Soundness is the only point at which a candidate outside the rectangle could
appear to matter.  If the strict comparison failed, the offending competitor
would have no greater error than the reference; if the weak comparison
failed, it would have strictly smaller error.  In either case its error is at
most $1/2$, and its denominator is no higher than the reference denominator,
which is at most $B$.  The proof of
\cref{lem:all-best-rectangle} therefore places its numerator and denominator
inside the same checked rectangle.  The finite test has already rejected the
proposed failure.

For completeness, every genuinely best reference satisfies
\eqref{eq:best-error-half} and the bounds of
\cref{lem:all-best-rectangle}, so it occurs in the rectangle and passes every
checked comparison by definition.  Each signed component list contains both
signs of the $B+1$ or $P+1$ possible magnitudes, including a harmless second
copy of zero.  The two denominator coordinates and two numerator coordinates
therefore give
\[
       \bigl(2(B+1)\bigr)^2\bigl(2(P+1)\bigr)^2
        =16(B+1)^2(P+1)^2
\]
displayed entries.  Structural filtering cannot increase their number.
\end{proof}

\leanref{RFAP.ConstructiveGaussianHeight.GaussianInt.componentSquare;
RFAP.BestApproximation.PassesBestTest;
RFAP.BestApproximation.passesBestTest_iff_isBest;
RFAP.BestApproximation.candidates_length;
RFAP.BestApproximation.allBest;
RFAP.BestApproximation.allBest_length_le;
RFAP.BestApproximation.mem_allBest_iff_isBest}

\begin{theorem}[executable consecutive-best successor]
\label{thm:consecutive-best-selector}
Fix a displayed source $z$ and a displayed lower pair $c_-$.  Filter the same
candidate rectangle by retaining a candidate $c_+$ exactly when both endpoint
pairs pass the best test, their denominator quadrances satisfy
$\Norm{q_-}<\Norm{q_+}$, and every candidate at an intermediate denominator
quadrance fails the best test.  The resulting finite list
$\operatorname{NextBest}(z,c_-)$ satisfies
\begin{equation}
\label{eq:next-best-selector-membership}
 c_+\in\operatorname{NextBest}(z,c_-)
 \quad\Longleftrightarrow\quad
 c_-,c_+\text{ occupy consecutive best levels}.
\end{equation}
Its length is at most $16(B+1)^2(P+1)^2$.
\end{theorem}

\begin{proof}
Suppose first that the finite test accepts $c_+$.  Exactness of the endpoint
best tests proves that $c_-$ and $c_+$ are best.  If a genuinely best pair
lay at an intermediate denominator quadrance, the global rectangle lemma
would place it among the checked candidates, contradicting the final clause
of the test.  The two endpoint levels are therefore consecutive.

Conversely, suppose the endpoints occupy consecutive best levels.  Both
endpoint tests pass by \eqref{eq:all-best-membership}.  No candidate passing
the best test can have intermediate denominator quadrance, because exactness
of that test would turn it into the best pair excluded by consecutiveness.
Thus the finite conjunction passes and the filter retains $c_+$.  Filtering
the same candidate rectangle also gives the stated length bound.
\end{proof}

\leanref{RFAP.BestApproximation.AreConsecutiveBest;
RFAP.BestApproximation.PassesConsecutiveTest;
RFAP.BestApproximation.passesConsecutiveTest_iff;
RFAP.BestApproximation.consecutiveSuccessors;
RFAP.BestApproximation.consecutiveSuccessors_length_le;
RFAP.BestApproximation.mem_consecutiveSuccessors_iff;
RFAP.ApproximationPaper.mem_consecutiveBestSuccessors_iff}

\begin{theorem}[total next-best step for a finite RAT response]
\label{thm:rat-response-next-best}
Let $T$ be any RFAP-1 RAT, let $\sigma$ be one of its displayed current
states, and execute a positive-accuracy request
$\sigma'=T(\sigma,\eps)$.  Let $z_R$ be the canonical common-denominator
presentation of the rational anchor of $\sigma'$.  Then $z_R$ belongs to the
returned winding-one box,
\[
 \diamq(B')\le\eps^2,
 \qquad
 \Norm{f(z_R)}\le L_1(f,A)\diamq(B'),
\]
and $B'$ refines the current box.  The executable list
$\operatorname{AllBest}(z_R)$ contains all and only best approximants to
$z_R$ and is nonempty.

If $c$ is any member of this list, a terminating optional calculation returns
one $c_+$ satisfying
\[
       c,c_+\text{ occupy consecutive best levels},
\]
or returns empty.  It returns empty if and only if
$\mathcal E_{z_R}(c)=0$.  Consequently every nonterminal call produces the
next occupied best level, and every emitted pair is a genuine best
approximant to the displayed RAT response.
\end{theorem}

\begin{proof}
The RFAP-1 track call supplies the refined winding-one state, its squared
diagonal bound, and its rational anchor.  Canonical common-denominator
conversion preserves that anchor exactly.  Nonzero winding gives the stated
anchor residual bound, while the track record gives refinement and
accuracy.  Conversion from the positive natural common denominator to a real
Gaussian-integer denominator again preserves the value exactly, so
\cref{thm:all-best-selector,thm:consecutive-best-selector} apply.  Finite
descent from the exact pair representing $z_R$ proves that the all-best table
is nonempty.

For totality, the same descent produces a best pair of zero error.  If $c$
has nonzero error, its best property forces
its denominator quadrance to be strictly below that terminal level.  Taking
the least denominator among the finite best table strictly above $c$ produces
the consecutive successor.  If $c$ has zero error, a later best pair would
have to improve strictly below zero, contradicting nonnegativity of
quadrance.  These two finite arguments prove both directions of the empty
option criterion.
\end{proof}

\leanref{RFAP.RATBestApproximation.trackResponseState_refines;
RFAP.RATBestApproximation.trackResponseState_accurate;
RFAP.RATBestApproximation.trackResponseSample_inside;
RFAP.RATBestApproximation.trackResponseSample_residual_le;
RFAP.RATBestApproximation.mem_trackBestApproximations_iff;
RFAP.RATBestApproximation.trackBestApproximations_nonempty;
RFAP.RATBestApproximation.trackNextBestApproximation?_sound;
RFAP.RATBestApproximation.trackNextBestApproximation_or_exactTerminal;
RFAP.RATBestApproximation.respondToTrack;
RFAP.ApproximationPaper.respondToRAT}

\begin{corollary}[RFAP-II best-response strengthening of FTA and RFAP]
\label{cor:fta-rfap-best-response}
Every output branch of the RFAP-1 finite constructive FTA carries the
following additional RFAP-II arithmetic capability.
\begin{enumerate}[label=(\roman*)]
\item A degree-one exact terminal $z$ has a nonempty executable table
      containing all and only its best Gaussian-rational approximants.
\item For every one of the degree-indexed RATs in a higher-degree output,
      every displayed current state and every positive rational accuracy
      request produce the complete proof-bearing response of
      \cref{thm:rat-response-next-best}.
\end{enumerate}
The same alternatives apply to every signed numerator or denominator site in
the RFAP-1 rational-function action principle: an exact linear site uses its
displayed terminal point, while every finite call to every higher-degree site
RAT returns the refined box, residual-bound anchor, all-and-only-best table,
and total next-best-or-terminal step.  Existing multiplicity signs, winding
certificates, separation, ownership, and maximality data are retained
unchanged.
\end{corollary}

\begin{proof}
Case-split the RFAP-1 positive-degree FTA result.  In the linear branch,
canonical common-denominator conversion of the exact displayed point followed
by the finite all-best calculation gives the first alternative.  In the
higher branch, each output is already a degree-indexed
\texttt{WindingOneTrack}; apply \cref{thm:rat-response-next-best} to the
selected track, current state, and request.  This is the finite construction
packaged by \texttt{strengthenFTAResult}.

Each signed rational-function site has the identical exact-or-higher shape.
The exact branch uses the point construction, and the higher branch applies
the generic track construction to each track of the focused squarefree
factor.  This is \texttt{strengthenRationalSite}.  Both constructions only
append RFAP-II arithmetic records, so none of the Paper-I certificates is
altered and no dependency from RFAP-1 back to RFAP-2 is introduced.
\end{proof}

\leanref{RFAP.RATBestApproximation.PointBestResponse;
RFAP.RATBestApproximation.respondToPoint;
RFAP.RATBestApproximation.TrackBestResponse;
RFAP.RATBestApproximation.respondToTrack;
RFAP.ApproximationPaper.FTAApproximationCapability;
RFAP.ApproximationPaper.strengthenFTAResult;
RFAP.ApproximationPaper.RationalSiteApproximationCapability;
RFAP.ApproximationPaper.strengthenRationalSite}

\begin{lemma}[coefficient primitivity and Cramer certificate]
\label{lem:coefficient-cramer}
A best pair has no nonunit common Gaussian-integer scale: if
$c=\lambda d$ coefficientwise and $q_d\ne0$, then $\Norm\lambda=1$.

For arbitrary displayed pairs $u,v,x$, put
$\delta=\det(u,v)$.  The following coefficientwise identity is an executable
Cramer certificate:
\begin{equation}
\label{eq:coefficient-cramer}
 \det(x,v)u+\det(u,x)v=\delta x.
\end{equation}
Multiplying the two displayed coefficients on the left by $\bar\delta$ gives
\begin{equation}
\label{eq:coefficient-cramer-norm}
 \bar\delta\det(x,v)u+\bar\delta\det(u,x)v
      =\Norm\delta\,x.
\end{equation}
In particular, if $\Norm\delta=1$, the displayed coefficients in
\eqref{eq:coefficient-cramer-norm} express every $x$ as an integral
combination of $u,v$.  If $\Norm\delta=2$, the same coefficients express
$2x$ as such a combination, giving the finite annihilation certificate behind
the centered index-two alternative.
\end{lemma}

\begin{proof}
If $c=\lambda d$, multiplication of Gaussian quadrances gives
\[
 \Norm{q_c}=\Norm\lambda\Norm{q_d},
 \qquad
 \mathcal E_z(c)=\Norm\lambda\mathcal E_z(d).
\]
The nonzero denominator makes $\Norm\lambda$ positive.  If it exceeded one,
$d$ would have strictly smaller denominator while its error was no greater,
contradicting the strict lower-height clause for $c$.  Hence
$\Norm\lambda=1$.

Both coordinates of \eqref{eq:coefficient-cramer} are obtained by expanding
two products and cancelling the cross terms.  Multiplication by
$\bar\delta$ and the integral identity
$\bar\delta\delta=\Norm\delta$ give
\eqref{eq:coefficient-cramer-norm}.  When the norm is one, coefficientwise
multiplication by the real Gaussian integer $1$ leaves $x$ unchanged.
\end{proof}

\leanref{RFAP.BestApproximation.scalePair;
RFAP.BestApproximation.errorQuadrance_scalePair;
RFAP.BestApproximation.denominatorQuadrance_scalePair;
RFAP.BestApproximation.IsBest.scalingCoefficient_norm_eq_one;
RFAP.BestApproximation.linearCombination_cramer;
RFAP.BestApproximation.linearCombination_cramer_norm;
RFAP.BestApproximation.linearCombination_cramer_of_index_one;
RFAP.BestApproximation.linearCombination_cramer_of_index_two}

\begin{theorem}[constructive index-one/index-two dichotomy]
\label{lem:finite-coefficient-index}
For consecutive best pairs $c_-=(p_-,q_-)$ and $c_+=(p_+,q_+)$, put
\begin{equation}
\label{eq:coefficient-index-quadrance}
 I(c_-,c_+)=\Norm{p_-q_+-p_+q_-}.
\end{equation}
Then
\begin{equation}
\label{eq:coefficient-index-bound}
             I(c_-,c_+)\in\{1,2\}.
\end{equation}
\end{theorem}

\begin{proof}
The determinant-error identity is
\[
 p_-q_+-p_+q_-
   =q_-E_z(c_+)-q_+E_z(c_-).
\]
If the determinant vanished, multiplicativity of quadrance would give
\[
 \Norm{q_-}\mathcal E_z(c_+)
   =\Norm{q_+}\mathcal E_z(c_-).
\]
But $c_+$ is best and $\Norm{q_-}<\Norm{q_+}$, so
$\mathcal E_z(c_+)<\mathcal E_z(c_-)$.  Multiplying this strict inequality
by the positive $\Norm{q_-}$, and the denominator inequality by the
nonnegative $\mathcal E_z(c_-)$, makes the left side strictly smaller than
the right.  Thus the displayed Gaussian-integer determinant is nonzero and
its quadrance is positive.

It remains to prove the sharp upper bound without importing a lattice
index theorem.  Run the finite extended Gaussian Euclidean algorithm on
$p_-,q_-$.  It returns displayed $g,s,t,p_0,q_0$ satisfying
\[
       g=sp_-+tq_- ,\qquad (p_-,q_-)=g(p_0,q_0).
\]
Primitivity of $c_-$ forces $\Norm g=1$.  Multiplying the B\'ezout
coefficients by $\bar g$ therefore gives $a,b\in\Zi$ with
$ap_-+bq_-=1$.  The displayed complement $V=(-b,a)$ satisfies
$\det(c_-,V)=1$, so the Cramer certificate gives an exact normal form
\begin{equation}
\label{eq:index-basis-normal-form}
                    c_+=\beta c_-+\delta V,
       \qquad \delta=\det(c_-,c_+).
\end{equation}

Apply one nearest-Gaussian-integer division to $\beta/\delta$.  Writing
$\beta=k\delta+r$ and $V'=kc_-+V$ turns
\eqref{eq:index-basis-normal-form} into
\begin{equation}
\label{eq:index-reduced-normal-form}
                    c_+=r c_-+\delta V',
       \qquad 2\Norm r\le\Norm\delta=I(c_-,c_+).
\end{equation}
The half-open square supplies a little more than this usual half-quadrance
bound.  If $\Norm\delta\ge4$, direct integral calculation of the dot and
cross coordinates of the remainder gives the alternative
\begin{equation}
\label{eq:index-remainder-clearance}
  2(\Norm r+1)\le\Norm\delta
  \quad\text{or}\quad
  \delta=r(-1+i).
\end{equation}
Indeed, failure of the first inequality leaves only
$\Norm\delta=2\Norm r$ or $2\Norm r+1$.  The odd case is incompatible with
the two scaled half-open coordinate bounds and Lagrange's identity.  In the
even case equality forces both scaled coordinates to be the included lower
endpoint, so multiplication out gives the displayed diagonal factorization.
This is a finite integer argument; no volume, compactness, or quotient
cardinality occurs.

In the first case of \eqref{eq:index-remainder-clearance}, $r\ne0$ by
primitivity of $c_+$.  The two-term quadrance estimate and the strict endpoint
inequalities give
\[
 \begin{split}
 \Norm\delta\,\mathcal E_z(V')
 &\le2\bigl(\mathcal E_z(c_+)+\Norm r\,\mathcal E_z(c_-)\bigr)
  <2(\Norm r+1)\mathcal E_z(c_-),\\
 \Norm\delta\,\Norm{q_{V'}}
 &\le2\bigl(\Norm{q_+}+\Norm r\,\Norm{q_-}\bigr)
  <2(\Norm r+1)\Norm{q_+}.
 \end{split}
\]
Thus $V'$ lies strictly inside the mixed cylinder, contradicting
\cref{lem:mixed-cylinder-criterion}.  In the second case,
\eqref{eq:index-reduced-normal-form} reads
$c_+=r(c_-+(-1+i)V')$.  Primitivity of $c_+$ forces $\Norm r=1$, and hence
$\Norm\delta=2$, contrary to $\Norm\delta\ge4$.

Therefore $0<I(c_-,c_+)<4$.  An integral Gaussian quadrance below four is
$0$, $1$, or $2$; positivity excludes zero and proves
\eqref{eq:coefficient-index-bound}.
\end{proof}

\leanref{RFAP.GaussianEuclidean.extendedGCD;
RFAP.GaussianEuclidean.remainder_dot_cross_bounds;
RFAP.GaussianEuclidean.remainder_clearance_or_diagonal_factor;
RFAP.BestApproximation.coefficientDeterminant;
RFAP.BestApproximation.coefficientIndexQuadrance;
RFAP.BestApproximation.toGaussian_coefficientDeterminant;
RFAP.BestApproximation.determinant_ne_zero_of_ordered_best;
RFAP.BestApproximation.AreConsecutiveBest.coefficientIndexQuadrance_pos;
RFAP.BestApproximation.IsBest.euclideanGCD_norm_eq_one;
RFAP.BestApproximation.IsBest.coefficientDeterminant_bezoutComplement_eq_one;
RFAP.BestApproximation.IsBest.exists_reducedCoefficientBasis;
RFAP.BestApproximation.AreConsecutiveBest.coefficientIndexQuadrance_le_two;
RFAP.BestApproximation.AreConsecutiveBest.coefficientIndexQuadrance_eq_one_or_two;
RFAP.ApproximationPaper.consecutiveBest_coefficientIndexQuadrance_pos;
RFAP.ApproximationPaper.consecutiveBest_coefficientIndexQuadrance_eq_one_or_two}

\begin{lemma}[finite index-two error chamber]
\label{lem:index-two-error-chamber}
Let $u,v$ occupy consecutive best levels for the displayed source $z$.  If
$\Norm{\det(u,v)}=2$, then
\begin{equation}
\label{eq:index-two-error-ratio}
                  \mathcal E_z(u)<4\mathcal E_z(v).
\end{equation}
\end{lemma}

\begin{proof}
For each Gaussian unit $\alpha\in\{1,-1,i,-i\}$, the parity construction
produces a displayed nonzero pair $c_\alpha$ satisfying
\[
       \det(u,v)c_\alpha=v-\alpha u.
\]
The mixed cylinder between $u$ and $v$ is empty.  Consequently each of these
four identities supplies one of two exact exterior inequalities: either the
error quadrance of $v-\alpha u$ is at least
$2\mathcal E_z(u)$, or its denominator quadrance is at least
$2\Norm{q_v}$.

The four alternatives are incompatible with
$4\mathcal E_z(v)\le\mathcal E_z(u)$.  Two error alternatives attached to
opposite units contradict the parallelogram identity.  Two attached to
adjacent units contradict the corresponding pair of orthogonal-projection
identities under the assumed factor-four bound.  If at most one alternative
is an error alternative, at least three are denominator alternatives, hence
two use opposite units and again contradict the parallelogram identity.
This finite sixteen-case calculation proves
\eqref{eq:index-two-error-ratio}; it contains no square root or angular
choice.
\end{proof}

\leanref{RFAP.BestApproximation.AreConsecutiveBest.exists_centeredUnitCandidate_outsideCylinder;
RFAP.BestApproximation.AreConsecutiveBest.centeredUnit_selfObstruction;
RFAP.BestApproximation.four_unit_obstructions_force_error_lt_four_mul;
RFAP.BestApproximation.AreConsecutiveBest.errorQuadrance_lower_lt_four_mul_upper_of_index_two;
RFAP.ApproximationPaper.consecutiveBest_indexTwoErrorRatio}

For Gaussian integers $a$ and $\delta$, write $\delta\mid a$ only when a
displayed Gaussian integer $k$ satisfies $a=\delta k$.  This witness-bearing
notation is important below: no quotient ring is being introduced.

\begin{theorem}[finite centered-coset and repeated-index classification]
\label{thm:finite-centered-coset}
Let $u,v$ be consecutive best pairs and put $\delta=\det(u,v)$.  Every
displayed coefficient pair $x\in\Zi^2$ satisfies one of
\begin{align}
 x&=a u+b v &&(a,b\in\Zi),
 \label{eq:centered-integral-class}\\
 \delta x&=a u+b v &&(a,b\in\Zi, \delta\nmid a, \delta\nmid b).
 \label{eq:centered-odd-class}
\end{align}
If $\Norm\delta=1$, the first alternative always holds.  If
$\Norm\delta=2$, the two alternatives give the ordinary span and its
centered coset.  Conversely, for every $a,b$ with
$\delta\nmid a$ and $\delta\nmid b$, the equation
\eqref{eq:centered-odd-class} has a displayed pair $x$ as a solution.

Moreover, suppose $u,v$ and $v,w$ are consecutive best pairs and both
determinants have quadrance two.  There are displayed
$\varepsilon,k\in\Zi$ such that
\begin{equation}
\label{eq:repeated-index-normal-form}
          w=\varepsilon u+\delta k v,
          \qquad \Norm\varepsilon=1,
          \qquad \Norm k\in\{1,2,4\}.
\end{equation}
\end{theorem}

\begin{proof}
Every quadrance-two Gaussian integer is one of the four diagonal elements
$\pm1\pm i$.  Direct multiplication therefore proves the exact test
\begin{equation}
\label{eq:diagonal-divisibility-parity}
 \delta\mid a
 \quad\Longleftrightarrow\quad
 \Ree(a)+\Imm(a)\text{ is even}.
\end{equation}
The integer quotient and remainder on division by two give the parity split
constructively.  In particular, the nondivisible class is the odd class; it
is closed under multiplication, while the sum of two members of that class
is divisible by $\delta$.

Complete $u$ to the explicit B\'ezout basis $(u,V)$ used in the proof of
\cref{lem:finite-coefficient-index}.  Write
\[
       v=\beta u+\delta V,
       \qquad x=\gamma u+\eta V.
\]
If $\delta\mid\beta$, the first identity gives $v$ a common factor of
quadrance two, contrary to primitivity; hence $\delta\nmid\beta$.  If
$\delta\mid\eta$, substitution gives integral coefficients for $x$ in
$u,v$.  Otherwise
\[
 \delta x=(\delta\gamma-\beta\eta)u+\eta v.
\]
Both displayed coefficients are nondivisible: if the first were divisible,
then so would be $\beta\eta$, contradicting closure of the odd class under
multiplication.  This proves the forward classification.  Conversely, when
$a,b$ are nondivisible, so is $\beta b$, and therefore
$a+\beta b=\delta q$ for a computed $q\in\Zi$.  The pair
$x=q u+bV$ then satisfies $\delta x=a u+bv$.

For the last assertion, first apply the classification to $w$.  Its centered
alternative would force a nondivisible cleared coefficient to have
quadrance two after taking determinants.  But every quadrance-two Gaussian
integer is divisible by $\delta$, a contradiction.  Hence
$w=\varepsilon u+bv$; taking the determinant with $v$ and using both
index-two hypotheses gives $\Norm\varepsilon=1$.

It remains to show $\delta\mid b$.  The centered construction for the second
pair supplies the displayed half-sum $h$ with
$\det(v,w)h=v+w$.  Classifying $h$ relative to $u,v$, its integral
alternative is ruled out immediately by multiplicativity of quadrance: it
would give $2(2\Norm r)=2$ for an integral coefficient $r$.  Comparing the
two coefficient expressions in the centered alternative shows that
$1+b$ is nondivisible by $\delta$.  Since $-1$ is also nondivisible, their
sum $b=(1+b)+(-1)$ is divisible by $\delta$.  The resulting quotient is the
displayed $k$ in \eqref{eq:repeated-index-normal-form}.

The strict increase from the denominator level of $u$ to that of $w$ shows
$k\ne0$.  It remains to bound its quadrance.  Apply
\cref{lem:index-two-error-chamber} to $u,v$.  In the error identity obtained
from \eqref{eq:repeated-index-normal-form}, the first summand has quadrance
$\mathcal E_z(u)<4\mathcal E_z(v)$, whereas the second has quadrance
$2\Norm k\,\mathcal E_z(v)$.  If $\Norm k\ge5$, the latter is therefore at
least $10\mathcal E_z(v)$.  The exact two-coordinate Cauchy identity
\[
 (x_1y_1+x_2y_2)^2+(x_1y_2-x_2y_1)^2
   =(x_1^2+x_2^2)(y_1^2+y_2^2)
\]
and the expansion of the quadrance of their sum then show that the error of
$w$ cannot be smaller than $\mathcal E_z(v)$.  This contradicts the strict
best comparison between the consecutive levels $v,w$.  Hence
$0<\Norm k\le4$; the displayed coordinate-square classification leaves
exactly $\Norm k=1,2,$ or $4$.

This proves the finite counterpart of the coefficient bound in Chevallier's
repeated-index lemma and supplies the bounded input for the exclusion below.
\end{proof}

\leanref{RFAP.GaussianCoset.coordinateEven_or_coordinateOdd;
RFAP.GaussianCoset.divides_iff_coordinateEven_of_norm_two;
RFAP.GaussianCoset.divides_of_norm_two;
RFAP.GaussianCoset.not_divides_mul_of_norm_two;
RFAP.GaussianCoset.divides_add_of_not_divides_of_norm_two;
RFAP.BestApproximation.InCoefficientSpan;
RFAP.BestApproximation.InCenteredCoefficientCoset;
RFAP.BestApproximation.IsBest.inCoefficientSpan_or_inCenteredCoefficientCoset_of_index_two;
RFAP.BestApproximation.IsBest.exists_centeredCoefficientPair_of_index_two;
RFAP.BestApproximation.AreConsecutiveBest.coefficientSpan_or_centeredCoset;
RFAP.BestApproximation.AreConsecutiveBest.integralTransition_with_divisible_middleCoefficient;
RFAP.BestApproximation.AreConsecutiveBest.repeatedIndexTwo_boundedNormalForm;
RFAP.BestApproximation.AreConsecutiveBest.repeatedIndexTwo_finiteNormalForm;
RFAP.ApproximationPaper.consecutiveBest_coefficientSpan_or_centeredCoset;
RFAP.ApproximationPaper.consecutiveBest_indexTwoTransition;
RFAP.ApproximationPaper.consecutiveBest_indexTwoFiniteTransition}

\begin{theorem}[no successive index-two transitions]
\label{thm:no-successive-index-two}
Let $u,v$ and $v,w$ occupy consecutive best levels for the displayed source
$z$.  If
\[
                       \Norm{\det(u,v)}=2,
\]
then
\[
                       \Norm{\det(v,w)}=1.
\]
\end{theorem}

\begin{proof}
Suppose instead that both determinant quadrances are two.  By
\cref{thm:finite-centered-coset}, write
\[
        w=\varepsilon u+\delta k v,
        \qquad \Norm\varepsilon=1,
        \qquad \Norm k\in\{1,2,4\},
        \qquad \delta=\det(u,v).
\]
Put $a=\delta k$.  Rotate the lower error and denominator by the displayed
unit $\varepsilon$, and divide by the two displayed nonzero reference
coordinates.  This constructs exact Gaussian-rational coordinates $x,y$
for which
\begin{align}
 \Norm x&<1,& 1&<4\Norm x,&
 \Norm{1+ax}&<\Norm x,\label{eq:repeated-normalized-error}\\
 \Norm y&<1,& 1&<\Norm{y+a}.\label{eq:repeated-normalized-denominator}
\end{align}
The first and last strict inequalities are just the error decrease and
denominator increase through the two consecutive levels; the middle bound in
\eqref{eq:repeated-normalized-error} is
\cref{lem:index-two-error-chamber}.  The two affine expressions are the
coefficientwise error and denominator identities obtained from the displayed
normal form.

For every displayed Gaussian unit $\zeta$, the centered candidate outside
the first mixed cylinder gives
\begin{equation}
  2\le\Norm{x-\zeta}
  \quad\text{or}\quad
  2\le\Norm{1-\zeta y},
  \label{eq:repeated-self-obstruction}
\end{equation}
while the same candidate outside the next mixed cylinder gives
\begin{equation}
  2\Norm x\le\Norm{x-\zeta}
  \quad\text{or}\quad
  2\Norm{y+a}\le\Norm{1-\zeta y}.
  \label{eq:repeated-next-obstruction}
\end{equation}
These are homogeneous consequences of the displayed cylinder certificates;
normalization cancels only positive rational quadrances.

Multiplicativity gives
$\Norm a=2\Norm k\in\{2,4,8\}$.  The bounded Gaussian-integer
classification makes $a$ a displayed unit associate of, respectively,
$1+i$, $2$, or $2+2i$.  Rotating $x$ by that unit and $y$ by its displayed
conjugate inverse preserves every quadrance above and merely permutes the
four choices of $\zeta$.  It therefore suffices to exclude those three
canonical coefficients.

For each coefficient, substitute $\zeta\in\{1,-1,i,-i\}$ into
\eqref{eq:repeated-self-obstruction}--\eqref{eq:repeated-next-obstruction}
and expand the quadrances coefficientwise.  For $a=2$, the opposite real
unit tests force the boundary equality for the real projection; together
with \eqref{eq:repeated-normalized-error} and $1<4\Norm x$ this gives both
$5\Norm x<1$ and $1<4\Norm x$.  For $a=1+i$ and $a=2+2i$, the real and
imaginary tests leave only the corresponding diagonal crossing patterns.
Their strict transition alternatives force a coordinate of $x$ or $y$
outside its open unit disc; the sole weak boundary pattern would give a
rational square whose value is three, which is impossible by the reduced
numerator--denominator parity argument.  Thus all three canonical
coefficients are impossible.

The index-one/index-two dichotomy
\cref{lem:finite-coefficient-index} leaves quadrance one as the only value
for $\det(v,w)$.
\end{proof}

\leanref{RFAP.RepeatedIndex.NormalizedCylinderData;
RFAP.RepeatedIndex.normalizeCylinderData;
RFAP.RepeatedIndex.transportNormalizedData;
RFAP.RepeatedIndex.diagonalStep_obstructions_impossible;
RFAP.RepeatedIndex.doubledAxisStep_obstructions_impossible;
RFAP.RepeatedIndex.doubledDiagonalStep_obstructions_impossible;
RFAP.RepeatedIndex.repeatedIndexTwo_impossible;
RFAP.RepeatedIndex.nextCoefficientIndexQuadrance_eq_one_of_eq_two;
RFAP.ApproximationPaper.consecutiveBest_repeatedIndexTwo_impossible;
RFAP.ApproximationPaper.consecutiveBest_nextIndex_eq_one_of_indexTwo}

\begin{theorem}[bounded coefficient-local successor generation]
\label{thm:local-successor-generation}
Let $u,v$ occupy consecutive best levels for $z=N/D$, and put
\[
 C=4\bigl(\Norm D+\Norm{q_u}\bigr),
 \qquad \delta=\det(u,v).
\]
Form a finite candidate list by taking each of the eight
$a\in\Zi$ with $\Norm a\in\{1,2\}$, every displayed $b\in\Zi$ whose two
coordinate magnitudes are at most $C$, and both expressions
\begin{equation}
 au+bv,
 \qquad
 \operatorname{half}\bigl(\bar\delta(au+bv)\bigr).
 \label{eq:local-successor-forms}
\end{equation}
Here $\operatorname{half}$ divides both displayed integral coordinates by
two; the second expression is a total calculation, and is used as exact
division only when the centered identity proves those coordinates even.

Call a candidate \emph{improving} when its denominator quadrance is strictly
greater than that of $v$ but no greater than $\Norm D$, and its error
quadrance is strictly smaller than that of $v$.  Retain an improving candidate
$c$ when every improving candidate in the same local list has denominator
quadrance at least $\Norm{q_c}$ and every one at the same denominator level has
error quadrance at least $\mathcal E_z(c)$.  This coefficient-local test
returns exactly the pairs $w$ for which $v,w$ occupy consecutive best levels.
In particular, the local generator skips no next best level.  Before
filtering, its length is at most
\begin{equation}
                         64(C+1)^2.
 \label{eq:local-successor-count}
\end{equation}
\end{theorem}

\begin{proof}
Let $w$ occupy the next best level above $v$.  If $\Norm\delta=1$, the
coefficientwise Cramer identity gives
\[
                              w=au+bv.
\]
Taking the determinant with $v$ shows
$\Norm a=\Norm{\det(v,w)}\in\{1,2\}$.

If $\Norm\delta=2$, \cref{thm:no-successive-index-two} gives
$\Norm{\det(v,w)}=1$.  The integral-span branch would instead make that
quadrance equal to $2\Norm a$, so it is impossible.  The centered-coset
branch is therefore forced:
\[
                              \delta w=au+bv.
\]
Taking determinants now gives $\Norm a=1$.  Since
$\bar\delta\delta=2$, multiplying this identity by $\bar\delta$ proves that
the second expression in \eqref{eq:local-successor-forms} computes $w$
coefficientwise; no quotient is selected.

Both branches have the common denominator identity
\[
     bq_v=sq_w-aq_u,
     \qquad \Norm s\le2,
     \qquad \Norm a\le2,
\]
where $s=1$ in the integral case and $s=\delta$ in the centered case.
The elementary difference bound
$\Norm{x-y}\le2(\Norm x+\Norm y)$, the source bound
$\Norm{q_w}\le\Norm D$, and $q_v\ne0$ give
\[
 \Norm b
 \le \Norm b\Norm{q_v}
 \le 2\bigl(2\Norm D+2\Norm{q_u}\bigr)=C.
\]
Thus $a$ and $b$ occur in the displayed coefficient lists, and $w$ occurs in
the local candidate list.

For completeness, let $w$ be an actual next pair.  Emptying the mixed cylinder
between $v$ and $w$ excludes every improving local contender of smaller
denominator quadrance.  Bestness of $w$ excludes a smaller error at the same
level.  Hence $w$ passes the local test.

For soundness, let $c$ pass that test.  Finite descent below $c$ produces a
best pair $r$ with no greater denominator or error.  Since $c$ improves on the
best pair $v$, the denominator of $r$ is strictly greater than that of $v$.
Select the least denominator quadrance among the finite complete list of best
pairs strictly above $v$; call a selected pair $w$.  Then $v,w$ are
consecutive and
\[
                  \Norm{q_w}\le\Norm{q_r}\le\Norm{q_c}.
\]
The completeness argument already places $w$ in the local candidate list, so
the local test for $c$ gives the reverse inequality.  Thus $c$ and $w$ have
the same denominator quadrance.  The same-level clause of the local test and
bestness of $w$ give equality of their error quadrances.  Bestness is unchanged
under these two exact ties, so $c$ occupies the same next level as $w$.

The finite all-best list is used here only to prove soundness of the local
calculation.  Evaluating the successor update itself enumerates and compares
only the displayed coefficient candidates.  This proves soundness and
completeness.
There are eight leading coefficients, $4(C+1)^2$ records in the signed
coordinate square for $b$, and two forms in
\eqref{eq:local-successor-forms}; multiplication gives
\eqref{eq:local-successor-count}.
\end{proof}

\leanref{RFAP.LocalSuccessor.halfGaussianInt;
RFAP.LocalSuccessor.divideByNormTwo;
RFAP.LocalSuccessor.divideByNormTwo_mul;
RFAP.LocalSuccessor.trailingCoefficientBound;
RFAP.LocalSuccessor.candidates;
RFAP.LocalSuccessor.next_mem_candidates;
RFAP.LocalSuccessor.selectDenominatorMinimum;
RFAP.LocalSuccessor.exists_consecutiveBest_below;
RFAP.LocalSuccessor.PassesLocalSuccessorTest;
RFAP.LocalSuccessor.passesLocalSuccessorTest_of_consecutive;
RFAP.LocalSuccessor.consecutive_of_passesLocalSuccessorTest;
RFAP.LocalSuccessor.passesLocalSuccessorTest_iff;
RFAP.LocalSuccessor.successors;
RFAP.LocalSuccessor.mem_successors_iff;
RFAP.LocalSuccessor.candidates_length_le;
RFAP.ApproximationPaper.localSuccessorCoefficientBound;
RFAP.ApproximationPaper.localConsecutiveBestSuccessors;
RFAP.ApproximationPaper.passesLocalSuccessorTest;
RFAP.ApproximationPaper.passesLocalSuccessorTest_iff;
RFAP.ApproximationPaper.mem_localConsecutiveBestSuccessors_iff;
RFAP.ApproximationPaper.localSuccessorCandidates_length_le}

\begin{theorem}[finite Dirichlet-window certificate]
\label{thm:finite-dirichlet-window}
For the supplied source $z=N/D$, put $B=\Norm D$.  Let $c_-$ and $c_+$
occupy consecutive best levels and define their exact
squared Dirichlet product by
\begin{equation}
\label{eq:dirichlet-window-product}
 \mathcal D_z(c_-,c_+)
   =\Norm{q_+}\mathcal E_z(c_-).
\end{equation}
For every natural quadrance request $H$ satisfying
\begin{equation}
\label{eq:dirichlet-window-request}
                 \Norm{q_-}<H\le\Norm{q_+},
\end{equation}
the lower pair has $q_-\ne0$ and supplies the certificate
\begin{equation}
\label{eq:dirichlet-window-bound}
 H\mathcal E_z(c_-)
   \le\mathcal D_z(c_-,c_+)
   \le\frac{B}{2},
 \qquad
 \mathcal D_z(c_-,c_+)\le\frac{1189}{102}I(c_-,c_+)
    \le\frac{1189}{51}.
\end{equation}
The exact product also satisfies the universal strict lower bounds
\begin{equation}
\label{eq:dirichlet-window-lower}
 \frac14<\mathcal D_z(c_-,c_+),
 \qquad
 I(c_-,c_+)=2\ \Longrightarrow\
 \frac12<\mathcal D_z(c_-,c_+).
\end{equation}
Writing
\[
 x=\frac{E_z(c_+)}{E_z(c_-)},\qquad
 y=\frac{q_-}{q_+},\qquad G=\Norm{1-yx},
\]
the finite coordinate-octant calculation gives
\begin{equation}
\label{eq:finite-octant-window-upper}
 \frac{205}{256}\le G,
 \qquad
 I(c_-,c_+)=1\quad\Longrightarrow\quad
 \mathcal D_z(c_-,c_+)\le\frac{256}{205}.
\end{equation}
The index-two centered-coset calculation strengthens the exact product to
\begin{equation}
\label{eq:index-two-window-upper}
 I(c_-,c_+)=2\quad\Longrightarrow\quad
 \frac{15}{56}\le\Norm x,\qquad
 \frac{15}{56}\le\Norm y,\qquad
 \frac{205}{128}\le G,\qquad
 \mathcal D_z(c_-,c_+)\le\frac{256}{205}.
\end{equation}
Consequently the index split gives the improved universal bound
\begin{equation}
\label{eq:index-aware-window-bound}
 \mathcal D_z(c_-,c_+)\le\frac{256}{205}
   <\frac{112}{71}<\frac{28}{15}<2<\frac{1189}{102}
   <\frac{1189}{51}<\frac{816}{35}.
\end{equation}
More exactly, every displayed rational $\tau$ satisfying
\begin{equation}
\label{eq:optimal-rational-gap-threshold-test}
 \frac34\le\tau<1,\qquad \tau^2-12\tau+9\ge0
\end{equation}
gives the index-specific gap bounds
\begin{equation}
\label{eq:optimal-rational-gap-bounds}
 \begin{aligned}
 I(c_-,c_+)=1&\quad\Longrightarrow\quad \tau\le G,\\
 I(c_-,c_+)=2&\quad\Longrightarrow\quad 2\tau\le G,
 \end{aligned}
\end{equation}
and hence the universal division-free certificate
\begin{equation}
\label{eq:optimal-rational-window-bound}
             \tau\mathcal D_z(c_-,c_+)\le1.
\end{equation}
The concrete value $\tau=205/256$ recovers
\eqref{eq:index-aware-window-bound}.  The cut itself, rather than a limiting
choice of $\tau$, is the stronger theorem.
The earlier one-coefficient calculation is retained as the separately
checkable root-free algebraic boundary
\begin{equation}
\label{eq:normalized-gap-quadratic}
 \frac12\le G
 \quad\hbox{or}\quad
G^2-3G+\frac14<0.
\end{equation}
Consequently every displayed rational $t$ satisfying
\begin{equation}
\label{eq:rational-gap-threshold-test}
 t<\frac12,\qquad t^2-3t+\frac14\ge0
\end{equation}
is a certified lower bound $t\le G$ and gives the division-free window bound
\begin{equation}
\label{eq:rational-threshold-window-bound}
       t\mathcal D_z(c_-,c_+)\le I(c_-,c_+)\le2.
\end{equation}
The displayed threshold $102/1189$ satisfies
\eqref{eq:rational-gap-threshold-test}; its integer certificate is
\begin{equation}
\label{eq:quadratic-gap-integer-certificate}
 4\mathbin{\cdot}102^2-12\mathbin{\cdot}102\mathbin{\cdot}1189+1189^2=1.
\end{equation}
Thus one exact finite output pair answers every integer squared-denominator
request in its whole consecutive window.
\end{theorem}

\begin{proof}
The nonzero denominator and the strict lower inequality in
\eqref{eq:dirichlet-window-request} are already part of the best-pair and
request data.  Multiplying $H\le\Norm{q_+}$ by the nonnegative
$\mathcal E_z(c_-)$ gives the first inequality in
\eqref{eq:dirichlet-window-bound}.  For the second, the global rectangle gives
$\Norm{q_+}\le B$, while the equal-height nearest-integer comparison in
\eqref{eq:best-error-half} gives $\mathcal E_z(c_-)\le1/2$.
Multiplying the two nonnegative inequalities proves the claim.

For the universal upper bound, use the displayed $x,y,G$ from the statement.
The lower error and upper denominator are nonzero: otherwise the strict error
or denominator ordering of consecutive best levels would fail.  Thus these
are displayed Gaussian-rational calculations, and denominator ordering also
gives $\Norm y<1$.  The empty-mixed-cylinder criterion, applied to the one
coefficient candidate $c_+-ac_-$, says for every $a\in\Zi$ that
\begin{equation}
\label{eq:normalized-coefficient-exterior}
 1\le\Norm{x-a}
 \quad\hbox{or}\quad
 1\le\Norm{1-ay}.
\end{equation}
More generally, the same empty-cylinder calculation applied to
$gc_-+hc_+$ gives, for every displayed $g,h\in\Zi$ not both zero,
\begin{equation}
\label{eq:normalized-linear-combination-exterior}
 1\le\Norm{g+hx}
 \quad\hbox{or}\quad
 1\le\Norm{gy+h}.
\end{equation}
The one-coefficient alternative above is the special case $(g,h)=(-a,1)$.

Choose $a$ by coordinatewise nearest-Gaussian-integer rounding and put
$r=x-a$, so $\Norm r\le1/2$.  The elementary weighted quadrance estimate
\begin{equation}
\label{eq:weighted-two-leg}
 60\Norm{u+v}\le204\Norm u+85\Norm v
\end{equation}
is just the nonnegativity of $\Norm{12u-5v}$.  The weights were chosen so that
\[
             204\frac{35}{408}+85\frac12=60.
\]
Suppose instead that $G<35/408$.  If $a=0$, then $x=r$ and
$\Norm{yx}<1/2$; applying
\eqref{eq:weighted-two-leg} to $1=(1-yx)+yx$ would give $\Norm1<1$.
If $a\ne0$, the first alternative in
\eqref{eq:normalized-coefficient-exterior} is impossible because
$\Norm r\le1/2$.  Hence the second holds.  But $\Norm{yr}<1/2$ and
\[
                  1-ay=(1-yx)+yr,
\]
so \eqref{eq:weighted-two-leg} again gives $\Norm{1-ay}<1$, a contradiction.
This proves $G\ge35/408$ using one explicitly computed coefficient.

The exact Cauchy calculation gives the stronger structural statement
\eqref{eq:normalized-gap-quadratic}.  If $G\ge1/2$ there is nothing to prove.
Otherwise put $R=\Norm{yr}<1/2$.  In the $a=0$ branch
$1=(1-yx)+yr$; in the other branch
\eqref{eq:normalized-coefficient-exterior} gives
$1\le\Norm{(1-yx)+yr}$.  If $d$ is the displayed coordinate dot product of
these two summands, expansion and the two-coordinate Cauchy identity give
\[
 1\le G+R+2d,
 \qquad d^2\le GR.
\]
Since $G+R<1$, squaring the positive margin gives
$(1-G-R)^2\le4GR$.  With $\delta=1/2-R>0$, direct expansion yields
\[
 (1-G-R)^2-4GR
   =\left(\frac12-G\right)^2-2G
      +\delta(1+2G+\delta).
\]
The left side is nonpositive and the last term is positive, so
$(1/2-G)^2-2G<0$, which is exactly the second alternative in
\eqref{eq:normalized-gap-quadratic}.  This records the boundary by rational
coefficient arithmetic rather than by naming either root of its quadratic.
To obtain the last assertion of the theorem, suppose a displayed $t$
satisfies \eqref{eq:rational-gap-threshold-test} but $G<t$.  The first branch
of \eqref{eq:normalized-gap-quadratic} is then impossible.  In the second
branch, $G+t<1<3$ and direct factorization gives
\[
 G^2-3G+\frac14
 =t^2-3t+\frac14+(t-G)(3-G-t)>0,
\]
contradicting that branch.  Thus sharper rational thresholds can be certified
by two finite comparisons, without changing the proof route.

Direct coefficient arithmetic factors the determinant identity as
\begin{equation}
\label{eq:normalized-window-index}
\mathcal D_z(c_-,c_+)G=I(c_-,c_+).
\end{equation}
Multiplying $t\le G$ by the nonnegative exact window gives
\eqref{eq:rational-threshold-window-bound}.
Combining this exact identity with $G\ge35/408$ gives
$\mathcal D_z(c_-,c_+)\le(408/35)I(c_-,c_+)$, the direct output of the
weighted estimate.  For the sharper displayed value $t=102/1189$, multiplying
$t^2-3t+1/4$ by the positive rational $4\mathbin{\cdot}1189^2$ gives the
integer certificate in \eqref{eq:quadratic-gap-integer-certificate}, hence
$t^2-3t+1/4>0$.  The quadratic boundary therefore gives $t\le G$, and exact
reciprocal cancellation gives
\[
 \mathcal D_z(c_-,c_+)\le\frac{1189}{102}I(c_-,c_+)
   \le\frac{1189}{51}.
\]

The additional coefficient tests improve that one-coefficient estimate.
Multiplication by a Gaussian unit, followed if necessary by coordinate
conjugation, places $x$ in one of the eight closed coordinate octants.
Multiply $y$ by the inverse unit and conjugate it whenever $x$ is conjugated.
This preserves $yx$ up to conjugation, hence preserves $G$.  Multiplying
$g$ by the matching unit, or conjugating both $g$ and $h$, shows directly
that \eqref{eq:normalized-linear-combination-exterior} is preserved as well.
We may therefore write
\[
\begin{gathered}
 x=a+ib,\qquad y=c+id,\qquad 0\le b\le a,\\
 X=\Norm x,\qquad Y=\Norm y,\qquad P=ac-bd.
\end{gathered}
\]
Here $X,Y<1$ and
\begin{equation}
\label{eq:first-octant-gap-expansion}
                 G=1+XY-2P.
\end{equation}
The three instances of
\eqref{eq:normalized-coefficient-exterior} with coefficients $1,i,1+i$
expand, coefficient by coefficient, to
\begin{align}
 2a\le X &\quad\hbox{or}\quad 2c\le Y,
       \label{eq:first-octant-one-obstruction}\\
 2b\le X &\quad\hbox{or}\quad -2d\le Y,
       \label{eq:first-octant-i-obstruction}\\
 2(a+b)\le X+1 &\quad\hbox{or}\quad c-d\le Y.
       \label{eq:first-octant-diagonal-obstruction}
\end{align}
If the first alternative of
\eqref{eq:first-octant-one-obstruction} holds, then
$2a\le X=a^2+b^2\le2a^2$.  Since $0\le a<1$, this forces $a=0$, and then
$b=0$; thus $x=0$ and $G=1$.  In every other case we have
$2c\le Y$.  Three elementary square expansions will also be used:
\begin{equation}
\label{eq:first-octant-square-bounds}
 2a\le1+X,\qquad 2(a+b)\le1+2X,\qquad -2d\le1+Y.
\end{equation}

It remains only to inspect the two other displayed alternatives.  The
elementary square identity
\[
 0\le Y\bigl((u-\tfrac12)^2+v^2\bigr)
       +(1-Y)\frac14
\]
shows, whenever $u^2+v^2=U$ and $0\le Y\le1$, that
$uY\le UY+1/4$.  If $-2d\le Y$, the two denominator-axis obstructions and
$Y<1$ force $c\le0$ or $d\ge0$.  In either sign case the preceding component
bound gives $2P\le XY+1/4$.  The same conclusion follows when
$2b\le X$ and $c-d\le Y$: the diagonal denominator obstruction forces the
corresponding sign alternative, and one applies the component bound to
either $a$ or $-d$.  Hence both of these chambers satisfy
\[
                 G=1+XY-2P\ge\frac34.
\]

The sole remaining chamber has both non-real exclusions on the error side:
\[
             2b\le X,\qquad 2(a+b)\le X+1,
             \qquad 2c\le Y.
\]
The first two inequalities imply
\[
 2b-b^2\le a^2,qquad 2b-b^2\le(1-a)^2.
\]
As $0\le a\le1$, one of $a$ and $1-a$ is at most $1/2$; consequently
\[
 b^2-2b+\frac14\ge0,qquad 0\le b\le\frac12.
\]
At the displayed rational $r=13/97$, exact arithmetic gives
$r^2-2r+1/4<0$.  Factoring the difference of these two quadratic values
therefore forces $b\le r$.

The extra strength comes from retaining both coefficients in
\eqref{eq:normalized-linear-combination-exterior}.  Taking
$(g,h)=(1,-(1-i))$ gives
\begin{equation}
\label{eq:first-octant-norm-two-obstruction}
 1\le\Norm{1-(1-i)x}
 \quad\hbox{or}\quad
 1\le\Norm{y-(1-i)}.
\end{equation}
If $x\ne0$, then $a>0$.  Since $X<1$ and $0\le b\le a$, we have
$a^2<a$ and $b^2\le b$, hence $X<a+b$.  But the error alternative in
\eqref{eq:first-octant-norm-two-obstruction} expands to $a+b\le X$ and is
therefore impossible.  If $x=0$, then $G=1$ already.  Thus in the remaining
case the denominator alternative supplies the second projection
\begin{equation}
\label{eq:first-octant-second-denominator-projection}
                 2(c-d)\le Y+1.
\end{equation}
Together with $2c\le Y$, it gives
\[
 2P=2(a-b)c+2b(c-d)
    \le(a-b)Y+b(Y+1)=aY+b,
\]
and consequently
\begin{equation}
\label{eq:first-octant-projected-gap}
                 G\ge1+XY-aY-b.
\end{equation}

As a preliminary, only small rational comparisons remain.  Exact
cross-multiplication gives
\[
 \frac{71}{112}
 \le\frac34+r^2-r,
 \qquad
 \left(\frac34+r^2-r\right)-\frac{71}{112}
   =\frac{13}{1053808}>0.
\]
Since $b\le r\le1/2$, the nonnegative product
$(r-b)(1-r-b)$ gives $r^2-r\le b^2-b$.  Therefore
\begin{equation}
\label{eq:first-octant-scalar-bound}
 \frac{71}{112}
 \le (a-\tfrac12)^2+\frac34+b^2-b
 =1+X-a-b.
\end{equation}
If $X\le a$, subtracting the last expression from the right side of
\eqref{eq:first-octant-projected-gap} leaves
$(1-Y)(a-X)\ge0$.  If $a\le X$, that right side is instead at least
$1-b\ge1-r\ge71/112$.  Thus this first four-coefficient calculation gives
$G\ge71/112$.

Retaining the other norm-two leading combination sharpens the same finite
octant calculation.  In addition to the inequalities above, the coefficient
$1+i$ supplies
\[
       a-b\le X\quad\hbox{or}\quad 2(c+d)\le Y+1.
\]
The five alternatives divide the displayed $(c,d)$ data into four regions:
\[
\begin{array}{ll}
 \text{green:}&Y<c-d,\\
 \text{red:}&c-d\le Y<-2d,\\
 \text{exterior:}&-2d\le Y\ \hbox{and}\ 2(c+d)\le Y+1,\\
 \text{blue:}&Y+1<2(c+d).
\end{array}
\]
This is only a decidable split of rational inequalities.  In the red and
blue regions, the displayed signs and projections give $2P\le XY$, hence
$G\ge1$.  In the exterior region they give either $8P^2\le XY$ or
$7P^2\le XY$; both imply the common dyadic certificate
\[
                         256P^2\le51XY,
\]
which yields $G\ge205/256$ by completing a square.

The green region is the extremal one.  It forces $d<0$.  If $c\le0$, a
direct expansion writes the required remainder as a sum of the nonnegative
terms $-2ac$, $Xc^2$, $(X-2b)d^2$, and
$\tfrac b2(2d+1)^2$.  If $0\le c$, put $e=-d$.  The error and denominator
rows then have the symmetric form
\[
\begin{gathered}
 0\le b\le a,\quad 2b\le X,\quad2(a+b)\le X+1,\\
 0\le c\le e,\quad 2c\le Y,\quad2(e+c)\le Y+1.
\end{gathered}
\]
Each row implies its slope quadratic, and an exact homogeneous identity
combining the two quadratics with $(be-ac)^2$ gives
\[
                         4(ac+be)^2\le XY.
\]
The remaining radial estimate is the following one-vector certificate:
\begin{equation}
\label{eq:dyadic-one-vector-projection}
 b+X(a-X)\le\frac{51}{256}.
\end{equation}
First the two error projections give $b\le35/256$.  When
$256b\le13$, the identity
$4-27u(1-u)^2=(4-3u)(3u-1)^2$ for $u=1-a$ gives
\[
 b+X(a-X)\le\frac{13}{256}+\frac4{27}\le\frac{51}{256}.
\]
For $13\le t=256b\le35$, put $L=(5t+82)/512$.  Splitting once at
$t=21$, direct polynomial products prove
\[
 L^2\le2b-b^2,qquad
 b+\frac14(1-L)^2\le\frac{51}{256};
\]
the diagonal projection gives $L\le1-a$, and hence
\eqref{eq:dyadic-one-vector-projection}.  Applying this certificate to the
row with smaller norm and using the two linear projections of $ac+be$ gives
\[
             G=1+XY-2(ac+be)\ge1-\frac{51}{256}
               =\frac{205}{256}.
\]
Thus all four regions, and therefore all eight finite symmetry choices,
satisfy the first inequality in
\eqref{eq:finite-octant-window-upper}; no limiting value or completed complex
input enters the calculation.  When $I(c_-,c_+)=1$, the exact identity
\eqref{eq:normalized-window-index} immediately gives
$\mathcal D_z(c_-,c_+)\le256/205$.

The same chamber calculation admits an exact rational parameter.  Let
$\tau$ satisfy
\eqref{eq:optimal-rational-gap-threshold-test} and put $r=1-\tau$.
Then $r\ge0$ and direct expansion gives
\begin{equation}
\label{eq:optimal-projection-polynomial}
 r^2+10r-2=\tau^2-12\tau+9\ge0.
\end{equation}
The one-vector projection lemma used in the green region has the exact form
\[
 0\le a,b,\quad X=a^2+b^2\le1,\quad
 2b\le X,\quad 2(a+b)\le X+1
 \quad\Longrightarrow\quad
 b+X(a-X)\le r
\]
for every displayed $r\ge0$ satisfying
$r^2+10r-2\ge0$.  Its proof first derives
$b^2-2b+1/4\ge0$ and $b\le1/7$.  A split at the displayed rational $3/64$
then uses, respectively, the cubic identity already displayed above and a
rational linear support to the same coordinate quadratic; after clearing
denominators, each remainder is an explicit product of nonnegative rational
factors.  Thus the lemma identifies the whole polynomial cut, rather than
only the point $51/256$.

Apply it to the error row and to the transformed denominator row.  The same
two projection inequalities used above now give $G\ge1-r=\tau$ in the green
region.  The red and blue regions already give $G\ge1$, the exterior region
gives $G\ge6/7$, and the nonpositive-real branch of the green region gives
$G\ge13/14$; direct polynomial comparisons put every $\tau$ satisfying
\eqref{eq:optimal-rational-gap-threshold-test} below those three rational
floors.  Hence the first line of
\eqref{eq:optimal-rational-gap-bounds} holds in every one of the eight
displayed octants.

It remains to treat coefficient index two.
Suppose $I(c_-,c_+)=2$.  For each Gaussian unit $\zeta$, the centered-coset
certificate constructed in \cref{lem:index-two-error-chamber}, normalized by
the positive lower-error and upper-denominator quadrances, gives
\begin{equation}
\label{eq:index-two-unit-obstruction}
 2\le\Norm{x-\zeta}
 \quad\hbox{or}\quad
 2\le\Norm{1-\zeta y}.
\end{equation}
Put $X=\Norm x$, $Y=\Norm y$, $\alpha=1-X>0$, and $\beta=1-Y>0$.
The alternatives for opposite units must cross: the two error alternatives
for $\zeta,-\zeta$ would have sum $2(1+X)<4$, and the two denominator
alternatives would have sum $2(1+Y)<4$.  Expanding the crossed alternatives
for $\zeta=1,-1$ therefore gives
\[
       4\Ree(x)\Ree(y)\le-\alpha\beta,
\]
while the alternatives for $\zeta=i,-i$ give
\[
       \alpha\beta\le4\Imm(x)\Imm(y).
\]
Adding these inequalities and cancelling the positive rational factor two
shows
\[
 2\Ree(yx)=2\bigl(\Ree(x)\Ree(y)-\Imm(x)\Imm(y)\bigr)
   \le-\alpha\beta.
\]
Hence, entirely coefficientwise,
\begin{equation}
\label{eq:index-two-balanced-gap}
 G-1=X Y-2\Ree(yx)\ge X Y+\alpha\beta.
\end{equation}

The individual crossed alternatives give a second useful consequence.  In
the real direction they are exactly one of
\[
 \alpha\le-2\Ree(x),\quad \beta\le2\Ree(y),
 \qquad\hbox{or}\qquad
 \beta\le-2\Ree(y),\quad \alpha\le2\Ree(x).
\]
In the imaginary direction they are exactly one of
\[
 \alpha\le-2\Imm(x),\quad \beta\le-2\Imm(y),
 \qquad\hbox{or}\qquad
 \beta\le2\Imm(y),\quad \alpha\le2\Imm(x).
\]
The right sides selected above are nonnegative because $\alpha,\beta\ge0$.
Squaring the two bounds which concern $\alpha$, adding them, and cancelling
two gives
\[
             (1-X)^2=\alpha^2\le2X.
\]
The two complementary bounds give the identical result
$(1-Y)^2\le2Y$.  Thus each of $X,Y$ satisfies
\begin{equation}
\label{eq:index-two-coordinate-quadratic}
                 U^2-4U+1\le0.
\end{equation}

No root of this quadratic is needed.  Put $s=15/56$.  Exact arithmetic gives
\[
                  s^2-4s+1=\frac1{3136}>0.
\]
If a displayed $U$ satisfying
\eqref{eq:index-two-coordinate-quadratic} had $U<s$, then
\[
 U^2-4U+1
   =s^2-4s+1+(s-U)(4-U-s)>0,
\]
a contradiction.  Hence $s\le X,Y$.

It remains to bound the balanced term in
\eqref{eq:index-two-balanced-gap}.  Write
\[
                  F=XY+(1-X)(1-Y).
\]
If $Y\le1/2$, then
$F-Y=(1-X)(1-2Y)\ge0$, so $F\ge Y\ge s$.  If $1/2\le Y$, then
\[
 F-s=(X-s)(2Y-1)+(1-Y)(1-2s)\ge0.
\]
Thus $F\ge15/56$, giving the preliminary bound $G\ge71/56$.

The orientation already chosen for the index-one calculation extracts more
from the same four centered alternatives.  Write
\[
 x=a+ib,\qquad y=-u+iv,\qquad
 \ell(U)=\frac{2(U+1)}7,\qquad
 h(U)=\frac{1-U}{2}.
\]
The first-octant inequalities and the forced orientation of the unit
alternatives give
\[
 0\le b\le a,\qquad h(X)\le b,\qquad
 0\le h(Y)\le u,v.
\]
For every displayed $15/56\le U\le1$, the exact endpoint factorization
proves
\begin{equation}
\label{eq:centered-affine-minorant}
                    \ell(U)^2\le\frac U2.
\end{equation}
Thus $\ell(X)\le a$ and
$\ell(Y)\le\max(u,v)$, with no square root selected.

There are now only two rational orientation cases.  If $u\le v$, the
coordinate products dominate the cross floor
\[
 XY+\ell(X)(1-Y)+(1-X)\ell(Y).
\]
If $v\le u$, they dominate the aligned floor
\[
 XY+2\ell(X)\ell(Y)+\frac{(1-X)(1-Y)}2.
\]
Both displayed expressions are increasing from the rational corner
$X=Y=15/56$ in the required directions; clearing denominators at that corner
gives the same exact bound $77/128$.  Consequently
\[
 G\ge1+\frac{77}{128}=\frac{205}{128}.
\]
The exact identity \eqref{eq:normalized-window-index} now gives
\[
 \mathcal D_z(c_-,c_+)\le
 \frac{2}{205/128}=\frac{256}{205}
\]
in index two.  The index-one and index-two cases therefore have the same
displayed bound and prove \eqref{eq:index-aware-window-bound}.  Every
comparison in this calculation is rational and exact.

To obtain the second line of
\eqref{eq:optimal-rational-gap-bounds}, retain the parameter $\tau$ and put
\[
 s=\frac{\tau}{3},\qquad
 m=\frac{s}{1-s^2},\qquad
 \ell_s(U)=m(U+1).
\]
The threshold test gives $1/4\le s<1/3$ and
\begin{equation}
\label{eq:optimal-centered-norm-cut}
 9(s^2-4s+1)=\tau^2-12\tau+9\ge0.
\end{equation}
If a displayed $U<1$ satisfying
$U^2-4U+1\le0$ had $U<s$, then
\[
 U^2-4U+1
  =s^2-4s+1+(s-U)(4-U-s)>0,
\]
a contradiction.  Hence $s\le X,Y$.

The denominator $1-s^2$ is positive.  Moreover
\[
 (1-s^2)^2-12s^2
   =(s^2+4s+1)(s^2-4s+1)\ge0,
\]
so $12m^2\le1$.  At the lower endpoint the defining cancellation is
\begin{equation}
\label{eq:optimal-centered-affine-endpoint}
              \ell_s(s)(1-s)=s.
\end{equation}
Together with $(1-s)^2\ge2s$, this gives
$\ell_s(s)^2\le s/2$.  For $s\le U\le1$, the exact difference
factorization
\begin{equation}
\label{eq:optimal-centered-affine-minorant}
 \frac U2-\ell_s(U)^2
 =\left(\frac s2-\ell_s(s)^2\right)
  +(U-s)\left(\frac12-m^2(U+s+2)\right)
\end{equation}
is nonnegative: $U+s+2\le4$ and $4m^2\le1/2$.
Thus $\ell_s(X)\le a$ and
$\ell_s(Y)\le\max(u,v)$ by the same square comparison as before.

In the cross case $u\le v$, the resulting floor simplifies to
\begin{align*}
 &XY+\ell_s(X)(1-Y)+(1-X)\ell_s(Y)\\
 &\qquad=2m+(1-2m)XY
 \ge2m+(1-2m)s^2
 =s^2+2s
 \ge6s-1=2\tau-1.
\end{align*}
Here $m\le1/2$, and the penultimate inequality is precisely
$s^2-4s+1\ge0$.  In the aligned case $v\le u$, the floor is increasing from
$(s,s)$: its coordinate slope is
\[
 U+2m^2(U+1)-\frac{1-U}{2}\ge0
\]
for $U\ge s$, using $s\ge1/4$ and $m\ge s$.  At the lower corner the aligned
floor exceeds the cross floor by the explicit square
\[
 2\left(\ell_s(s)-\frac{1-s}{2}\right)^2.
\]
It therefore has the same lower bound $2\tau-1$.  Adding the leading one in
the determinant-gap expansion proves $G\ge2\tau$.  With
$I(c_-,c_+)=2$, cancellation in
\eqref{eq:normalized-window-index} gives
$\tau\mathcal D_z(c_-,c_+)\le1$.  Together with the index-one calculation,
this proves \eqref{eq:optimal-rational-window-bound}.

For the lower bound, put
$A=\Norm{q_+}\mathcal E_z(c_-)$ and
$A'=\Norm{q_-}\mathcal E_z(c_+)$.  Strict denominator and error ordering give
$A'<A$.  The coefficient determinant identity and the elementary quadrance
difference estimate therefore give
\[
\begin{split}
 I(c_-,c_+)
  &=\Norm{q_-E_z(c_+)-q_+E_z(c_-)}\\
  &\le2(A'+A)<4A.
\end{split}
\]
The index is one or two, which yields
\eqref{eq:dirichlet-window-lower}.  The Lean statements retain the
division-free forms $1<4A$ and, in the index-two case, $2<4A$.

The use of a natural quadrance request is intentional: if a conventional
modulus threshold is denoted by $Q$, then $H$ is its squared counterpart, but
the theorem never has to introduce $Q=\sqrt H$.  The exact product
\eqref{eq:dirichlet-window-product} is normally much sharper than $B/2$ and is
computed from the returned pairs.  The source-independent rational
$256/205$ is a convenient fixed instance, while
\eqref{eq:optimal-rational-window-bound} records the whole finite quadratic
cut.  Conventionally naming its smaller endpoint as $6-3\sqrt3$ makes the
reciprocal endpoint $(2+\sqrt3)/3$, Chevallier's optimal squared constant;
the theorem itself needs neither name.
\end{proof}

\begin{theorem}[finite sharpness beyond every scale]
\label{thm:finite-dirichlet-window-sharpness}
Let $\tau\in\Q$ satisfy
\begin{equation}
\label{eq:sharper-rational-window-test}
 \frac34\le\tau<1,
 \qquad \tau^2-12\tau+9<0,
\end{equation}
and let $T\in\N$ be a supplied finite scale.  Put
\[
 \delta=-(\tau^2-12\tau+9),
 \qquad
 k=\left\lceil\frac{24}{\delta}\right\rceil+T+1.
\]
Starting from $(m_0,d_0)=(1,3)$, execute the subtraction-free recurrence
\begin{equation}
\label{eq:critical-pell-recurrence}
 (m_{j+1},d_{j+1})
   =(m_j+d_j,\,2m_j+3d_j),
 \qquad n_j=m_j+d_j,
\end{equation}
through stage $k$.  Define
\[
 s=\frac{m_k}{n_k},
 \qquad
 x=\frac12+i\frac{s}{2},
 \qquad
 y=\frac{s}{2}-i\frac12,
\]
and the two displayed Gaussian-integral pairs
\begin{equation}
\label{eq:critical-pell-pairs}
 c_-=(-2m_k,\,m_k-n_ki),
 \qquad
 c_+=(-m_k-n_ki,\,2n_k).
\end{equation}
Finally put
\begin{equation}
\label{eq:critical-pell-source}
 e_-=-\frac{1}{q_+(1-yx)},
 \qquad
 z=\frac{p_-+e_-}{q_-},
\end{equation}
where $c_\pm=(p_\pm,q_\pm)$.  These are executable displayed rational
data.  They satisfy
\begin{equation}
\label{eq:finite-window-sharpness-conclusion}
 \begin{aligned}
  T&<n_k<\Norm{q_+}=4n_k^2,\\
  c_-,c_+&\text{ occupy consecutive best levels for }z,\\
  1&<\tau\mathcal D_z(c_-,c_+).
 \end{aligned}
\end{equation}
Thus no displayed rational threshold satisfying
\eqref{eq:sharper-rational-window-test} can replace the quadratic cut as a
universal finite-window bound.
\end{theorem}

\begin{proof}
The recurrence preserves the natural-number identity
\begin{equation}
\label{eq:critical-pell-identity}
 d_j^2=2m_j^2+2m_jd_j+1,
 \qquad\text{equivalently}\qquad
 m_j^2+n_j^2=4m_jn_j+1.
\end{equation}
It also preserves $2m_j<d_j\le3m_j$.  Hence
\[
 \frac14\le s<\frac13,
 \qquad
 n_k^2(s^2-4s+1)=1,
\]
and $n_k>k>T$.  The second identity in
\eqref{eq:critical-pell-identity} gives
$p_-q_+-p_+q_-=1$.  Direct substitution into
\eqref{eq:critical-pell-source} gives
\[
 E_z(c_-)=e_-,
 \qquad E_z(c_+)=e_-x,
 \qquad \frac{q_-}{q_+}=y.
\]
The finite critical-cylinder calculation checks every nonzero pair of
Gaussian-integral coefficients after an explicit coefficient bound; at
$1/4\le s<1/3$ with $s^2-4s+1>0$, one of its normalized error or denominator
coordinates is exterior.  The displayed Cramer identity then proves that
both pairs are best and that their mixed open cylinder is empty.  They
therefore occupy consecutive best levels.

Put $G_s=\Norm{1-yx}$.  The definitions of $e_-$ and $\mathcal D$ cancel
exactly to
\begin{equation}
\label{eq:critical-gap-window-reciprocal}
             G_s\mathcal D_z(c_-,c_+)=1.
\end{equation}
Writing $P(s)=s^2-4s+1=1/n_k^2$, the exact factorization
\[
 16G_s=48s+P(s)(s^2+4s+17)
\]
gives the convenient strict estimate $G_s<3s+2P(s)$.  On the sharper side
of the threshold boundary, exact difference factorization gives
\[
 -(\tau^2-12\tau+9)<12(\tau-3s).
\]
The definition of $k$, the growth $n_k>k$, and $P(s)n_k^2=1$ give
\[
 24P(s)<-(\tau^2-12\tau+9).
\]
Consequently $2P(s)<\tau-3s$, and hence $G_s<\tau$.  The window product is
positive, so multiplying this strict inequality by it and applying
\eqref{eq:critical-gap-window-reciprocal} proves
$1<\tau\mathcal D_z(c_-,c_+)$.
\end{proof}

\leanref{RFAP.FiniteWindowSharpness.pellState;
RFAP.FiniteWindowSharpness.pellDifference_sq;
RFAP.FiniteWindowSharpness.pellDifference_bounds;
RFAP.FiniteWindowSharpness.pellEquation;
RFAP.FiniteWindowSharpness.pellParameter_isFiniteCritical;
RFAP.FiniteWindowSharpness.criticalPairs_areConsecutiveBest;
RFAP.FiniteWindowSharpness.criticalGap_mul_dirichletWindow;
RFAP.FiniteWindowSharpness.criticalGap_lt_sharperThreshold;
RFAP.FiniteWindowSharpness.requestedScale_lt_criticalUpperDenominatorQuadrance;
RFAP.FiniteWindowSharpness.finiteCriticalWindowSharpness;
RFAP.ApproximationPaper.finiteSharpnessStage;
RFAP.ApproximationPaper.finiteSharpnessSource;
RFAP.ApproximationPaper.finiteSharpnessLowerPair;
RFAP.ApproximationPaper.finiteSharpnessUpperPair;
RFAP.ApproximationPaper.finiteSharpnessDenominatorScale;
RFAP.ApproximationPaper.sharperGapThreshold_has_finiteCounterexample}

The preceding source depends on the requested scale.  A stronger finite
quantifier order is available when the whole desired run length is supplied
at once.

\begin{theorem}[finite same-source robust sharpness run]
\label{thm:finite-same-source-sharpness-run}
Let $\tau\in\Q$ satisfy
\eqref{eq:sharper-rational-window-test}, and supply $T,N\in\N$.  A terminating
calculation returns one displayed Gaussian-rational source $z$, one proper
rational box $B$ whose lower-left anchor is $z$, and a list
\[
  (c^-_0,c^+_0),\ldots,(c^-_{N-1},c^+_{N-1})
\]
of exactly $N$ pairs.  Their upper denominator quadrances obey
\[
 T<\Norm{q^+_0}<\Norm{q^+_1}<\cdots<\Norm{q^+_{N-1}}
\]
when $N>0$.  More importantly, for every displayed $w\in B$ and every $j<N$,
the pair
$(c^-_j,c^+_j)$ occupies consecutive best levels for $w$ and satisfies
\[
       1<\tau\mathcal D_w(c^-_j,c^+_j).
\]
Thus all data in one supplied finite run share one source description and one
uniform rational neighbourhood.  The source and box may still depend on
$T$ and $N$.
\end{theorem}

\begin{proof}
Put $b$ for the Pell stage selected by $\tau$ and $T$, and choose the common
source at stage $b+N$.  For $j<N$, transport the critical determinant-one
basis from the backward parameter $b+N-j$ to the forward parameter
$b+N+j$.  The subtraction-free Pell identities keep every parameter in the
critical chamber.  The finite mixed-cylinder calculation and its Cramer
certificate therefore prove that the transported pair is consecutive best
for the one chosen source.  The forward stages increase strictly with $j$,
so the corresponding upper denominator quadrances are strictly increasing;
their first value already lies beyond $T$.  The exact gap identity and the
same threshold estimate used in
Theorem~\ref{thm:finite-dirichlet-window-sharpness} give the strict violation
for every member of the list.

It remains to make these strict conclusions uniform.  For one window, first
bound every possible mixed-cylinder competitor by the explicit Gaussian
coefficient square with coefficient quadrance below $32$.  This is a finite
$20736$-entry list.  Each strict residual comparison in that list computes a
positive rational perturbation budget; the basis residuals and the window
inequality compute two further positive budgets.  Taking their finite minimum
gives a proper rational box on which that window remains consecutive and
violating.  Taking the finite minimum of these budgets over the $N$ displayed
windows gives $B$.  For $N=0$ the list is empty and the construction uses its
explicit default positive radius.  In either case all minima and all
membership tests are bounded rational calculations.
\end{proof}

\leanref{RFAP.CriticalBadWindowRun.criticalBadWindowRun;
RFAP.RobustCriticalRun.criticalFullyRobustRun;
RFAP.ApproximationPaper.robustFiniteSharpnessRun}

\begin{corollary}[two-window refinement]
\label{cor:two-dirichlet-windows}
Let $c_0,c_1,c_2$ occupy three successive best levels, and abbreviate
\[
 I_0=I(c_0,c_1),\quad I_1=I(c_1,c_2),\qquad
 \mathcal D_0=\mathcal D_z(c_0,c_1),\quad
 \mathcal D_1=\mathcal D_z(c_1,c_2).
\]
For every displayed rational $t$ satisfying
\eqref{eq:rational-gap-threshold-test},
\begin{equation}
\label{eq:two-dirichlet-window-threshold}
 t\mathcal D_0\le1\quad\hbox{or}\quad t\mathcal D_1\le1,
 \qquad
 t(\mathcal D_0+\mathcal D_1)\le3.
\end{equation}
For every displayed $\tau$ satisfying
\eqref{eq:optimal-rational-gap-threshold-test}, the stronger individual
window theorem gives
\begin{equation}
\label{eq:two-dirichlet-window-optimal-threshold}
 \tau\mathcal D_0\le1,\qquad
 \tau\mathcal D_1\le1,\qquad
 \tau(\mathcal D_0+\mathcal D_1)\le2.
\end{equation}
In particular,
\begin{equation}
\label{eq:two-dirichlet-window-bound}
 \begin{aligned}
 I_0+I_1&\le3, &
 \mathcal D_0&\le\frac{256}{205}, &
 \mathcal D_1&\le\frac{256}{205},\\
 I_j=2&\ \Longrightarrow\
   \mathcal D_j\le\frac{256}{205}, &
 \mathcal D_0+\mathcal D_1&\le\frac{512}{205}<4
   <\frac{1189}{34}.
 \end{aligned}
\end{equation}
Thus the finite octant and centered index-two calculations bound both
individual windows before they are added.
\end{corollary}

\begin{proof}
The index dichotomy gives $I_0=1$ or $I_0=2$.  In the first case the first
index is one.  In the second, the no-successive-index-two theorem gives
$I_1=1$.  Hence one index is one and in either case $I_0+I_1\le3$.
The threshold form of \cref{thm:finite-dirichlet-window} gives
$t\mathcal D_j\le I_j$ for each window.  Selecting the index-one inequality
and adding the two inequalities proves
\eqref{eq:two-dirichlet-window-threshold}.  The certified choice
$t=102/1189$ gives the weaker sum bound $1189/34$.
Applying \eqref{eq:optimal-rational-window-bound} separately to both
transitions and adding proves
\eqref{eq:two-dirichlet-window-optimal-threshold}.
The index-specific bounds in \cref{thm:finite-dirichlet-window} give the
universal estimate $\mathcal D_j\le256/205$ in either case.  Adding the two
universal estimates gives $512/205$; the remaining comparisons in
\eqref{eq:two-dirichlet-window-bound} are exact rational arithmetic.
\end{proof}

\begin{corollary}[finite-chain window bound]
\label{cor:finite-dirichlet-chain}
Let $c_0,\ldots,c_n$ be a supplied finite list in which every adjacent pair
occupies consecutive best levels.  For $0\le j<n$, put
\[
 I_j=I(c_j,c_{j+1}),\qquad
 \mathcal D_j=\mathcal D_z(c_j,c_{j+1}).
\]
Then
\begin{equation}
\label{eq:finite-chain-index-bound}
 2\sum_{j<n} I_j\le3n+(n\bmod2)\le3n+1.
\end{equation}
For every displayed rational $t$ satisfying
\eqref{eq:rational-gap-threshold-test}, the same chain satisfies
\begin{equation}
\label{eq:finite-chain-window-bound}
 2t\sum_{j<n}\mathcal D_j\le3n+(n\bmod2)\le3n+1.
\end{equation}
For every displayed $\tau$ satisfying
\eqref{eq:optimal-rational-gap-threshold-test}, the direct universal family
instead gives
\begin{equation}
\label{eq:finite-chain-optimal-threshold-bound}
       \tau\sum_{j<n}\mathcal D_j\le n.
\end{equation}
For the explicit threshold used above this becomes
\begin{equation}
\label{eq:finite-chain-refined-window-bound}
 2\sum_{j<n}\mathcal D_j
 \le\frac{1189}{102}\bigl(3n+(n\bmod2)\bigr).
\end{equation}
The finite-octant and centered-coset estimates give the stronger direct
accumulation
\begin{equation}
\label{eq:finite-chain-index-aware-window-bound}
 \sum_{j<n}\mathcal D_j
 \le\frac{256}{205}n.
\end{equation}
\end{corollary}

\begin{proof}
Group the transitions into adjacent blocks of two.  Each complete block has
index sum at most three by \cref{cor:two-dirichlet-windows}; precisely when
$n$ is odd, one transition remains and its index is at most two.  This finite
induction gives
\eqref{eq:finite-chain-index-bound}.  Summing the exact inequalities
$t\mathcal D_j\le I_j$ and multiplying by two gives
\eqref{eq:finite-chain-window-bound}.  Taking $t=102/1189$ and cancelling its
exact reciprocal $1189/102$ gives
\eqref{eq:finite-chain-refined-window-bound}.  The Lean proof performs
precisely this recursion on the supplied list and also checks that its
computed index and window lists have the same length.  Applying
\eqref{eq:optimal-rational-window-bound} to each displayed transition and
recursing over the supplied list proves
\eqref{eq:finite-chain-optimal-threshold-bound}.  For the final fixed bound,
each transition costs at most $256/205$ by
\eqref{eq:index-aware-window-bound}; summing those $n$ displayed inequalities
gives the result directly.
\end{proof}

\leanref{RFAP.BestApproximation.dirichletWindowQuadrance;
RFAP.DirichletWindow.AreConsecutiveBest.normalizedCoefficientExterior;
RFAP.DirichletWindow.AreConsecutiveBest.normalizedLinearCombinationExterior;
RFAP.DirichletWindow.quadraticGap_neg_of_sumExterior;
RFAP.DirichletWindow.AreConsecutiveBest.normalizedDeterminantGap_quadratic_neg_of_lt_half;
RFAP.DirichletWindow.AreConsecutiveBest.normalizedDeterminantGap_boundary;
RFAP.DirichletWindow.threshold_le_of_quadraticBoundary;
RFAP.DirichletWindow.AreConsecutiveBest.threshold_le_normalizedDeterminantGap_of_quadratic;
RFAP.DirichletWindow.AreConsecutiveBest.quadraticThreshold_mul_dirichletWindow_le_index;
RFAP.DirichletWindow.AreConsecutiveBest.quadraticThreshold_mul_dirichletWindow_le_two;
RFAP.DirichletWindow.AreConsecutiveBest.one_twelfth_le_normalizedDeterminantGap;
RFAP.DirichletWindow.refinedGapThreshold_lt_half;
RFAP.DirichletWindow.refinedGapThreshold_quadratic_nonnegative;
RFAP.DirichletWindow.quadraticGapThreshold_lt_half;
RFAP.DirichletWindow.quadraticGapThreshold_quadratic_nonnegative;
RFAP.DirichletWindow.quadraticGapThreshold_integerCertificate;
RFAP.DirichletWindow.quadraticGapThreshold_mul_quadraticIndexFactor;
RFAP.DirichletWindow.quadraticUniversalWindowBound_mul_denominator;
RFAP.DirichletWindow.quadraticUniversalWindowBound_lt_refinedUniversalWindowBound;
RFAP.DirichletWindow.refinedGapThreshold_mul_denominator;
RFAP.DirichletWindow.refinedGapThreshold_budget;
RFAP.DirichletWindow.refinedIndexFactor_mul_denominator;
RFAP.DirichletWindow.refinedUniversalWindowBound_mul_denominator;
RFAP.DirichletWindow.AreConsecutiveBest.refinedGap_le_normalizedDeterminantGap;
RFAP.DirichletWindow.AreConsecutiveBest.dirichletWindow_mul_normalizedGap_eq_index;
RFAP.DirichletWindow.AreConsecutiveBest.dirichletWindowQuadrance_le_factor_mul_index_of_quadraticThreshold;
RFAP.DirichletWindow.AreConsecutiveBest.dirichletWindowQuadrance_le_factor_mul_two_of_quadraticThreshold;
RFAP.DirichletWindow.AreConsecutiveBest.dirichletWindowQuadrance_le_refinedIndexFactor;
RFAP.DirichletWindow.AreConsecutiveBest.dirichletWindowQuadrance_le_refinedUniversalBound;
RFAP.DirichletWindow.AreConsecutiveBest.dirichletWindowQuadrance_le_quadraticIndexFactor;
RFAP.DirichletWindow.AreConsecutiveBest.dirichletWindowQuadrance_le_quadraticUniversalBound;
RFAP.DirichletWindow.AreConsecutiveBest.dirichletWindowQuadrance_le_twelve_mul_index;
RFAP.DirichletWindow.AreConsecutiveBest.dirichletWindowQuadrance_le_twentyFour;
RFAP.IndexOneWindow.half_le_determinantGap_of_firstOctantObstructions;
RFAP.IndexOneWindow.fifteen_twenty_eighths_le_determinantGap_of_firstOctantObstructions;
RFAP.IndexOneWindow.fifteen_twenty_eighths_le_determinantGap_of_coefficientObstructions;
RFAP.IndexOneWindow.half_le_determinantGap_of_coefficientObstructions;
RFAP.IndexOneWindow.AreConsecutiveBest.fifteen_twenty_eighths_le_normalizedDeterminantGap;
RFAP.IndexOneWindow.normTwoIndexOneGapThreshold_mul_denominator;
RFAP.IndexOneWindow.normTwoIndexOneWindowBound_mul_denominator;
RFAP.IndexOneWindow.normTwoIndexOneWindowBound_mul_gapThreshold;
RFAP.IndexOneWindow.AreConsecutiveBest.normTwoIndexOneGapThreshold_le_normalizedDeterminantGap;
RFAP.SharpIndexOneWindow.fourRegionSplit_of_fiveCoefficientObstructions;
RFAP.SharpIndexOneWindow.second_add_norm_mul_first_sub_norm_le_projectionBound;
RFAP.SharpIndexOneWindow.four_crossSquare_le_normProduct_of_slopeQuadratics;
RFAP.SharpIndexOneWindow.firstRegionGapBound_le_gap_of_greenRegion;
RFAP.SharpIndexOneWindow.firstRegionGapBound_le_gap_of_redRegion;
RFAP.SharpIndexOneWindow.firstRegionGapBound_le_gap_of_exteriorRegion;
RFAP.SharpIndexOneWindow.firstRegionGapBound_le_gap_of_blueRegion;
RFAP.SharpIndexOneWindow.firstRegionGapBound_le_gap_of_firstOctantFiveObstructions;
RFAP.OptimalIndexOneWindow.projection_le_of_firstRegionObstructions;
RFAP.OptimalIndexOneWindow.gapThreshold_le_gap_of_firstOctantFiveObstructions;
RFAP.IndexOneWindow.sharpIndexOneGapBound_mul_denominator;
RFAP.IndexOneWindow.sharpIndexOneWindowBound_mul_denominator;
RFAP.IndexOneWindow.AreConsecutiveBest.sharpIndexOneGapBound_le_normalizedDeterminantGap;
RFAP.IndexOneWindow.AreConsecutiveBest.dirichletWindowQuadrance_le_sharpIndexOneWindowBound_of_indexOne;
RFAP.OptimalIndexTwoWindow.componentLower_sq_le_half;
RFAP.OptimalIndexTwoWindow.twice_gapThreshold_le_determinantGap_of_firstOctantObstructions;
RFAP.IndexOneWindow.twice_gapThreshold_le_of_centeredUnitObstructions;
RFAP.IndexOneWindow.AreConsecutiveBest.twice_gapThreshold_le_normalizedDeterminantGap_of_indexTwo;
RFAP.IndexOneWindow.AreConsecutiveBest.gapThreshold_mul_dirichletWindow_le_one;
RFAP.IndexOneWindow.AreConsecutiveBest.half_le_normalizedDeterminantGap;
RFAP.IndexOneWindow.AreConsecutiveBest.dirichletWindowQuadrance_le_twentyEight_fifteenths_of_indexOne;
RFAP.IndexOneWindow.AreConsecutiveBest.dirichletWindowQuadrance_le_normTwoIndexOneWindowBound_of_indexOne;
RFAP.IndexOneWindow.AreConsecutiveBest.dirichletWindowQuadrance_le_two_of_indexOne;
RFAP.IndexOneWindow.AreConsecutiveBest.dirichletWindowQuadrance_le_twentyEight_fifteenths;
RFAP.IndexOneWindow.AreConsecutiveBest.dirichletWindowQuadrance_le_normTwoUniversalWindowBound;
RFAP.IndexOneWindow.AreConsecutiveBest.dirichletWindowQuadrance_le_two;
RFAP.IndexTwoWindow.one_add_balancedNormProduct_le_determinantGap_of_centeredUnitObstructions;
RFAP.IndexTwoWindow.complement_sq_le_two_norms_of_centeredUnitObstructions;
RFAP.IndexTwoWindow.fifteen_fiftySixths_le_norms_of_centeredUnitObstructions;
RFAP.IndexTwoWindow.seventyOne_fiftySixths_le_determinantGap_of_centeredUnitObstructions;
RFAP.IndexTwoWindow.firstOctant_projectionBounds_of_centeredUnitObstructions;
RFAP.IndexTwoWindow.centeredComponentLower_sq_le_half;
RFAP.IndexTwoWindow.dyadicCenteredGapThreshold_le_determinantGap_of_firstOctantObstructions;
RFAP.IndexTwoWindow.one_le_determinantGap_of_centeredUnitObstructions;
RFAP.IndexTwoWindow.AreConsecutiveBest.errorRatio_normSq_lt_one;
RFAP.IndexTwoWindow.AreConsecutiveBest.normalizedCenteredUnitObstruction;
RFAP.IndexTwoWindow.AreConsecutiveBest.seventyOne_fiftySixths_le_normalizedDeterminantGap_of_indexTwo;
RFAP.IndexTwoWindow.AreConsecutiveBest.one_le_normalizedDeterminantGap_of_indexTwo;
RFAP.IndexTwoWindow.AreConsecutiveBest.dirichletWindowQuadrance_le_oneHundredTwelve_seventyFirsts_of_indexTwo;
RFAP.IndexTwoWindow.AreConsecutiveBest.dirichletWindowQuadrance_le_two_of_indexTwo;
RFAP.IndexTwoWindow.AreConsecutiveBest.dirichletWindowQuadrance_le_quadraticIndexFactorUniversal;
RFAP.IndexOneWindow.AreConsecutiveBest.dyadicCenteredGapThreshold_le_normalizedDeterminantGap_of_indexTwo;
RFAP.IndexOneWindow.AreConsecutiveBest.dirichletWindowQuadrance_le_dyadicCenteredWindowBound_of_indexTwo;
RFAP.IndexOneWindow.AreConsecutiveBest.dirichletWindowQuadrance_le_sharpDyadicUniversalWindowBound;
RFAP.IndexTwoWindow.quadraticIndexFactor_mul_denominator;
RFAP.DirichletBlock.refinedIndexFactor_mul_three;
RFAP.DirichletBlock.refinedTwoWindowSumBound_mul_denominator;
RFAP.DirichletBlock.quadraticIndexFactor_mul_three;
RFAP.DirichletBlock.quadraticTwoWindowSumBound_mul_denominator;
RFAP.DirichletBlock.finiteOctantTwoWindowSumBound_mul_denominator;
RFAP.DirichletBlock.dyadicTwoWindowSumBound_mul_denominator;
RFAP.DirichletBlock.coefficientIndexQuadrance_add_le_three;
RFAP.DirichletBlock.quadraticThreshold_mul_one_dirichletWindow_le_one;
RFAP.DirichletBlock.quadraticThreshold_mul_dirichletWindowSum_le_three;
RFAP.DirichletBlock.one_dirichletWindowQuadrance_le_refinedIndexFactor;
RFAP.DirichletBlock.one_dirichletWindowQuadrance_le_quadraticIndexFactor;
RFAP.DirichletBlock.dirichletWindowQuadrance_add_le_refinedTwoWindowSumBound;
RFAP.DirichletBlock.dirichletWindowQuadrance_add_le_quadraticTwoWindowSumBound;
RFAP.DirichletBlock.dirichletWindowQuadrance_add_le_quadraticUniversalWindowBound;
RFAP.DirichletBlock.dirichletWindowQuadrance_add_le_finiteOctantTwoWindowSumBound;
RFAP.DirichletBlock.dirichletWindowQuadrance_add_le_dyadicTwoWindowSumBound;
RFAP.DirichletBlock.gapThreshold_mul_dirichletWindowQuadrance_add_le_two;
RFAP.DirichletBlock.dirichletWindowQuadrance_add_le_four;
RFAP.DirichletChain.dirichletWindowQuadrances_length_eq_coefficientIndexQuadrances_length;
RFAP.DirichletChain.two_mul_coefficientIndexQuadrances_sum_le_with_parity;
RFAP.DirichletChain.two_mul_coefficientIndexQuadrances_sum_le;
RFAP.DirichletChain.quadraticThreshold_mul_dirichletWindowQuadrances_sum_le_indexSum;
RFAP.DirichletChain.two_mul_quadraticThreshold_mul_dirichletWindowQuadrances_sum_le;
RFAP.DirichletChain.two_mul_quadraticThreshold_mul_dirichletWindowQuadrances_sum_le_with_parity;
RFAP.DirichletChain.two_mul_dirichletWindowQuadrances_sum_le_factor_mul_parityBound_of_quadraticThreshold;
RFAP.DirichletChain.two_mul_dirichletWindowQuadrances_sum_le_refinedIndexFactor_mul_parityBound;
RFAP.DirichletChain.two_mul_dirichletWindowQuadrances_sum_le_quadraticIndexFactor_mul_parityBound;
RFAP.DirichletChain.dirichletWindowQuadrances_sum_le_quadraticIndexFactor_mul_length;
RFAP.DirichletChain.dirichletWindowQuadrances_sum_le_twentyEight_fifteenths_mul_length;
RFAP.DirichletChain.dirichletWindowQuadrances_sum_le_normTwoUniversalWindowBound_mul_length;
RFAP.DirichletChain.dirichletWindowQuadrances_sum_le_sharpDyadicUniversalWindowBound_mul_length;
RFAP.DirichletChain.gapThreshold_mul_dirichletWindowQuadrances_sum_le_length;
RFAP.DirichletChain.fifteen_mul_dirichletWindowQuadrances_sum_le_twentyEight_mul_length;
RFAP.DirichletChain.fifteen_mul_dirichletWindowQuadrances_sum_le_with_parity;
RFAP.DirichletChain.dirichletWindowQuadrances_sum_le_two_mul_length;
RFAP.BestApproximation.AreConsecutiveBest.coefficientIndexQuadrance_lt_four_mul_dirichletWindow;
RFAP.BestApproximation.AreConsecutiveBest.one_lt_four_mul_dirichletWindow;
RFAP.BestApproximation.AreConsecutiveBest.two_lt_four_mul_dirichletWindow_of_indexTwo;
RFAP.BestApproximation.AreConsecutiveBest.lower_witnesses_dirichletWindow;
RFAP.BestApproximation.AreConsecutiveBest.dirichletWindowQuadrance_le_sourceBound;
RFAP.ApproximationPaper.consecutiveDirichletWindowQuadrance;
RFAP.ApproximationPaper.consecutiveBest_index_lt_four_mul_dirichletWindow;
RFAP.ApproximationPaper.one_lt_four_mul_consecutiveDirichletWindow;
RFAP.ApproximationPaper.two_lt_four_mul_consecutiveDirichletWindow_of_indexTwo;
RFAP.ApproximationPaper.one_twelfth_le_consecutiveNormalizedDeterminantGap;
RFAP.ApproximationPaper.refinedNormalizedDeterminantGapThreshold_lt_half;
RFAP.ApproximationPaper.refinedNormalizedDeterminantGapThreshold_quadratic_nonnegative;
RFAP.ApproximationPaper.refinedNormalizedDeterminantGapThreshold_mul_denominator;
RFAP.ApproximationPaper.refinedIndexDirichletWindowFactor_mul_denominator;
RFAP.ApproximationPaper.refinedUniversalDirichletWindowBound_mul_denominator;
RFAP.ApproximationPaper.refinedTwoWindowDirichletSumBound_mul_denominator;
RFAP.ApproximationPaper.quadraticNormalizedDeterminantGapThreshold_lt_half;
RFAP.ApproximationPaper.quadraticNormalizedDeterminantGapThreshold_quadratic_nonnegative;
RFAP.ApproximationPaper.quadraticNormalizedDeterminantGapThreshold_integerCertificate;
RFAP.ApproximationPaper.quadraticUniversalDirichletWindowBound_mul_denominator;
RFAP.ApproximationPaper.indexAwareUniversalDirichletWindowBound_mul_denominator;
RFAP.ApproximationPaper.quadraticUniversalDirichletWindowBound_lt_refined;
RFAP.ApproximationPaper.quadraticTwoWindowDirichletSumBound_mul_denominator;
RFAP.ApproximationPaper.finiteOctantNormalizedDeterminantGap_mul_denominator;
RFAP.ApproximationPaper.finiteOctantIndexOneDirichletWindowBound_mul_denominator;
RFAP.ApproximationPaper.centeredIndexTwoNormalizedDeterminantGap_mul_denominator;
RFAP.ApproximationPaper.centeredIndexTwoDirichletWindowBound_mul_denominator;
RFAP.ApproximationPaper.finiteOctantUniversalDirichletWindowBound_mul_denominator;
RFAP.ApproximationPaper.finiteOctantTwoWindowDirichletSumBound_mul_denominator;
RFAP.ApproximationPaper.dyadicNormalizedDeterminantGap_mul_denominator;
RFAP.ApproximationPaper.dyadicUniversalDirichletWindowBound_mul_denominator;
RFAP.ApproximationPaper.dyadicTwoWindowDirichletSumBound_mul_denominator;
RFAP.ApproximationPaper.refinedThreshold_le_consecutiveNormalizedDeterminantGap;
RFAP.ApproximationPaper.quadraticThreshold_le_consecutiveNormalizedDeterminantGap;
RFAP.ApproximationPaper.consecutiveNormalizedDeterminantGap_boundary;
RFAP.ApproximationPaper.threshold_le_consecutiveNormalizedDeterminantGap_of_quadratic;
RFAP.ApproximationPaper.quadraticThreshold_mul_consecutiveDirichletWindow_le_index;
RFAP.ApproximationPaper.quadraticThreshold_mul_consecutiveDirichletWindow_le_two;
RFAP.ApproximationPaper.consecutiveDirichletWindow_mul_normalizedGap_eq_index;
RFAP.ApproximationPaper.consecutiveBest_dirichletWindow_le_twelve_mul_index;
RFAP.ApproximationPaper.consecutiveBest_dirichletWindow_le_twentyFour;
RFAP.ApproximationPaper.consecutiveBest_dirichletWindow_le_refinedUniversalBound;
RFAP.ApproximationPaper.consecutiveBest_dirichletWindow_le_quadraticIndexFactor;
RFAP.ApproximationPaper.half_le_consecutiveNormalizedDeterminantGap;
RFAP.ApproximationPaper.finiteOctantGap_le_consecutiveNormalizedDeterminantGap;
RFAP.ApproximationPaper.dyadicGap_le_consecutiveNormalizedDeterminantGap;
RFAP.ApproximationPaper.dyadicCenteredIndexTwoGap_le_consecutiveNormalizedDeterminantGap;
RFAP.ApproximationPaper.consecutiveBest_dirichletWindow_le_finiteOctantIndexOneBound;
RFAP.ApproximationPaper.consecutiveBest_dirichletWindow_le_two_of_indexOne;
RFAP.ApproximationPaper.centeredIndexTwoNormalizedDeterminantGap_le_consecutive;
RFAP.ApproximationPaper.consecutiveBest_dirichletWindow_le_centeredIndexTwoBound;
RFAP.ApproximationPaper.one_le_consecutiveNormalizedDeterminantGap_of_indexTwo;
RFAP.ApproximationPaper.consecutiveBest_dirichletWindow_le_two_of_indexTwo;
RFAP.ApproximationPaper.consecutiveBest_dirichletWindow_le_finiteOctantUniversalBound;
RFAP.ApproximationPaper.consecutiveBest_dirichletWindow_le_dyadicUniversalBound;
RFAP.ApproximationPaper.consecutiveBest_dirichletWindow_le_indexAwareUniversalBound;
RFAP.ApproximationPaper.consecutiveBest_dirichletWindow_le_quadraticUniversalBound;
RFAP.ApproximationPaper.successiveBest_coefficientIndexQuadrance_add_le_three;
RFAP.ApproximationPaper.quadraticThreshold_mul_one_successiveDirichletWindow_le_one;
RFAP.ApproximationPaper.quadraticThreshold_mul_successiveDirichletWindowSum_le_three;
RFAP.ApproximationPaper.successiveBest_one_dirichletWindow_le_refinedIndexFactor;
RFAP.ApproximationPaper.successiveBest_dirichletWindow_add_le_refinedTwoWindowSumBound;
RFAP.ApproximationPaper.successiveBest_one_dirichletWindow_le_quadraticIndexFactor;
RFAP.ApproximationPaper.successiveBest_dirichletWindow_add_le_quadraticTwoWindowSumBound;
RFAP.ApproximationPaper.successiveBest_dirichletWindow_add_le_indexAwareTwoWindowSumBound;
RFAP.ApproximationPaper.successiveBest_dirichletWindow_add_le_finiteOctantTwoWindowSumBound;
RFAP.ApproximationPaper.successiveBest_dirichletWindow_add_le_dyadicTwoWindowSumBound;
RFAP.ApproximationPaper.successiveBest_dirichletWindow_add_le_four;
RFAP.ApproximationPaper.consecutiveDirichletWindowQuadrances_length_eq_indexQuadrances_length;
RFAP.ApproximationPaper.finiteBestChain_two_mul_coefficientIndexQuadrances_sum_le_with_parity;
RFAP.ApproximationPaper.finiteBestChain_two_mul_coefficientIndexQuadrances_sum_le;
RFAP.ApproximationPaper.finiteBestChain_quadraticThreshold_mul_windowSum_le_indexSum;
RFAP.ApproximationPaper.finiteBestChain_two_mul_quadraticThreshold_mul_windowSum_le;
RFAP.ApproximationPaper.finiteBestChain_two_mul_quadraticThreshold_mul_windowSum_le_with_parity;
RFAP.ApproximationPaper.finiteBestChain_two_mul_windowSum_le_refinedIndexFactor_mul_parityBound;
RFAP.ApproximationPaper.finiteBestChain_two_mul_windowSum_le_quadraticIndexFactor_mul_parityBound;
RFAP.ApproximationPaper.finiteBestChain_windowSum_le_indexAwareUniversalBound_mul_length;
RFAP.ApproximationPaper.finiteBestChain_windowSum_le_finiteOctantUniversalBound_mul_length;
RFAP.ApproximationPaper.finiteBestChain_windowSum_le_dyadicUniversalBound_mul_length;
RFAP.ApproximationPaper.finiteBestChain_fifteen_mul_windowSum_le_twentyEight_mul_length;
RFAP.ApproximationPaper.finiteBestChain_fifteen_mul_windowSum_le_with_parity;
RFAP.ApproximationPaper.finiteBestChain_windowSum_le_two_mul_length;
RFAP.ApproximationPaper.mem_consecutiveBestSuccessors_witnesses_dirichletWindow;
RFAP.ApproximationPaper.mem_consecutiveBestSuccessors_dirichletWindow_le_sourceBound;
RFAP.ApproximationPaper.mem_consecutiveBestSuccessors_dirichletWindow_le_twentyFour;
RFAP.ApproximationPaper.mem_consecutiveBestSuccessors_dirichletWindow_le_refinedUniversalBound;
RFAP.ApproximationPaper.mem_consecutiveBestSuccessors_dirichletWindow_le_quadraticUniversalBound;
RFAP.ApproximationPaper.mem_consecutiveBestSuccessors_dirichletWindow_le_indexAwareUniversalBound;
RFAP.ApproximationPaper.mem_consecutiveBestSuccessors_dirichletWindow_le_finiteOctantUniversalBound;
RFAP.ApproximationPaper.mem_consecutiveBestSuccessors_dirichletWindow_le_dyadicUniversalBound;
RFAP.CriticalBadWindowRun.criticalBadWindowRun;
RFAP.RobustCriticalRun.criticalFullyRobustRun;
RFAP.SuccessorDynamics.regeneratedTransition_atom;
RFAP.SuccessorDynamics.stoppedBlockHitCertificate;
RFAP.RobustSuccessorSharpness.robustSuccessorReturnProbability_eq_one;
RFAP.RobustSuccessorSharpness.finiteChevallierAnalogue;
RFAP.FiniteDirichletGlobal.globalDirichletRequest;
RFAP.FiniteDirichletGlobal.finiteChevallierBoundary}

\begin{lemma}[the terminal Hurwitz pair is best]
\label{lem:terminal-hurwitz-best}
The terminal current pair of every generated finite Hurwitz trace satisfies
the full best-approximation predicate.  Every displayed rational trace has
such a terminal stage, so $\operatorname{AllBest}(z)$ is nonempty.
\end{lemma}

\begin{proof}
At a terminal stage, the current error is zero, so the weak comparison with
every competitor is immediate.  For strictness, suppose a competitor
$c=(p,q)$ with $0<\Norm q<\Norm{q_n}$ also had zero error.  The cross-error
identity
\[
 \Delta(c_n,c)=q_nE_z(c)-qE_z(c_n)
\]
would give $\Delta(c_n,c)=0$.  In the adjacent-continuant basis of
\cref{lem:adjacent-basis-decomposition}, its previous coefficient is
$t=\Delta(c_n,c_{n-1})\Delta(c_n,c)=0$.  Hence the denominator identity
reduces to $q=sq_n$.  Since $q\ne0$, also $s\ne0$, and multiplicativity of
integral quadrance gives $\Norm q=\Norm s\Norm{q_n}\ge\Norm{q_n}$, a
contradiction.  Every lower-height competitor therefore has positive error,
which is strictly larger than the terminal zero error.

The stage recursion decreases the natural quadrance of its complete
denominator and stops exactly at zero remainder.  Structural recursion
therefore supplies a terminal member of every computed stage list.  Its
current pair is best by the preceding argument and hence belongs to
$\operatorname{AllBest}(z)$ by
\eqref{eq:all-best-membership}.
\end{proof}

\leanref{RFAP.HurwitzTrace.exists_terminal_mem_stages;
RFAP.HurwitzGeometry.Stage.current_error_eq_zero_of_terminal;
RFAP.HurwitzGeometry.Stage.isBest_of_generated_of_terminal;
RFAP.HurwitzGeometry.allBest_nonempty;
RFAP.ApproximationPaper.terminalHurwitzConvergent_isBest;
RFAP.ApproximationPaper.allBestGaussianApproximations_nonempty}

\section{Explicit arithmetic envelopes}
\label{sec:arithmetic-envelopes}

This section puts finite bounds around the RFAP-2 calculations named in the
glossary.  It first gives the normalized rational working layer used by the
repeated geometric constructions, then treats the square-threshold request,
grid record, homogeneous integral evaluation, finite Hurwitz recursion,
continuants, and protected quotient target.  These are bounds computed from
the displayed input, not maxima extracted after running the calculation.
The later envelopes begin after RFAP-1 has supplied the Newton-band anchor and
radius; they include the normalized Horner and Newton arithmetic performed
from that anchor.  Eliminant, Sturm, winding, and complete
repeated-subdivision costs still belong to the larger RFAP-1 resource problem.
The optional all-best selector is bounded separately by the exact rectangular
count in \cref{thm:all-best-selector}; it is not consumed by the Newton or
Hurwitz trace calculations bounded in this section.

\begin{definition}[component width]
\label{def:component-width}
For $n\in\N$, put
\[
                 \ell(n)=\lfloor\log_2n\rfloor+1,
\]
where the executable natural logarithm has $\log_2 0=0$.  Thus zero still
occupies one displayed bit.  For $x\in\Z$, put
$\ell_{\Z}(x)=\ell(|x|)$, and for $u=a+ib\in\Zi$ put
\[
                 \ell_{\Zi}(u)=\max\{\ell_{\Z}(a),\ell_{\Z}(b)\}.
\]
The width of $(a+ib)/D$ is the maximum of the two numerator-coordinate widths
and $\ell(D)$.  Since $|a|,|b|\le\Norm{a+ib}$, an upper bound $K$ on the
quadrance of a Gaussian integer gives the component-width bound $\ell(K)$.
\end{definition}

The Lean definition of $\ell$ is primitive natural recursion rather than a
real logarithm.  Its monotonicity and the component bounds are proved by
integer arithmetic.

\leanref{RFAP.CalculationSize.natWidth;
RFAP.CalculationSize.gaussianIntWidth;
RFAP.CalculationSize.normNat_eq_component_squares;
RFAP.CalculationSize.gaussianIntWidth_le_budget}

\subsection{Normalized working rationals}

\begin{definition}[normalized rational operations]
\label{def:normalized-rational-operations}
A working rational is a displayed record $a/b$ with $b>0$ and
$\gcd(|a|,b)=1$.  This canonical record is the rational scalar used by the
algebraic statements: two such records are equal exactly when their
numerators and denominators are structurally equal, equivalently when their
cross products are equal.  Given reduced records $a/b$ and $c/d$, put
\[
 g=\gcd(b,d),\qquad
 R_+=\frac{a(d/g)+c(b/g)}{(b/g)d}.
\]
The working sum is the constructive gcd reduction of $R_+$.  For
multiplication put
\[
 g_1=\gcd(|a|,d),\qquad g_2=\gcd(|c|,b),\qquad
 R_\times=\frac{(a/g_1)(c/g_2)}{(b/g_2)(d/g_1)},
\]
and again reduce the resulting record.  Thus certain cancellation occurs
before either product is formed.  The reciprocal of a nonzero record swaps
$|a|$ and $b$ and moves the sign to the numerator; zero maps to zero.  In
particular it avoids the squared-denominator intermediate produced by a
generic total-inverse construction.

Negation changes only the numerator sign; subtraction is normalized
negation followed by normalized addition.  Division is normalized reciprocal
followed by the cross-cancelled multiplication.  These six operations are
the complete scalar layer used below.

For the logical reference implementation of the natural gcd, define
\[
 G(0,n)=1,\qquad G(m,n)=1+G(n\bmod m,m)\quad(m>0).
\]
This counts the terminal clause and satisfies the elementary finite bound
$G(m,n)\le m+1$.
\end{definition}

\begin{lemma}[normalized-operation envelopes]
\label{lem:normalized-operation-envelopes}
Erasing each operation in \cref{def:normalized-rational-operations} gives
exactly the ordinary rational sum, product, or reciprocal.  If $A_+$ and
$B_+$ are the absolute numerator and positive denominator of the reduced
sum, then
\[
 A_+\le |a|d+|c|b,
 \qquad B_+\le bd.
\]
For the reduced product,
\[
 A_\times\le |a|\,|c|,
 \qquad B_\times\le bd.
\]
Consequently the returned records have respective bit-width bounds
\[
 \ell\!\left(\max\{|a|d+|c|b,bd\}\right),
 \qquad
 \ell\!\left(\max\{|a|\,|c|,bd\}\right).
\]
The exact logical gcd-clause counts are one denominator gcd plus the final
reduction for addition, and two cross-cancellation gcds plus the final
reduction for multiplication.  Each term is bounded by the displayed
$G(m,n)\le m+1$ recurrence.  Negation and reciprocal have one final
normalization; subtraction and division add that cost to the corresponding
binary operation.  Their output-component envelopes are respectively
$\max\{|a|,b\}$, the addition envelope, $\max\{|a|,b\}$, and the
multiplication envelope.
\end{lemma}

\begin{proof}
Write $b=(b/g)g$ and $d=(d/g)g$ and cross-multiply $R_+$ with the ordinary
sum; the two integer expressions agree coefficientwise.  The multiplication
identity follows in the same way after reconstructing each cancelled factor.
Constructive gcd normalization preserves these values and can only decrease
the absolute numerator and positive denominator.  The triangle inequality
for integer absolute value gives the addition bound; exact divisibility and
monotonicity of natural multiplication give the remaining component bounds.
Induction on the first argument of the displayed Euclidean recursion proves
the clause bound, since $n\bmod m<m$ whenever $m>0$.
\end{proof}

\leanref{RFAP.NormalizedFraction.addRaw;
RFAP.NormalizedFraction.mulRaw;
RFAP.NormalizedFraction.normalizedAdd_eq;
RFAP.NormalizedFraction.normalizedMul_eq;
RFAP.NormalizedFraction.normalizedSub_eq;
RFAP.NormalizedFraction.normalizedDiv_eq;
RFAP.NormalizedFraction.add_componentMagnitude_le;
RFAP.NormalizedFraction.mul_componentMagnitude_le;
RFAP.NormalizedFraction.subSteps_le_envelope;
RFAP.NormalizedFraction.divSteps_le_envelope;
RFAP.NormalizedFraction.gcdClauses_le;
RFAP.ApproximationPaper.normalizedCoordinateAdd_width_le_budget;
RFAP.ApproximationPaper.normalizedCoordinateMul_width_le_budget;
RFAP.ApproximationPaper.normalizedCoordinateSub_width_le_budget;
RFAP.ApproximationPaper.normalizedCoordinateDiv_width_le_budget}

Safe-cut candidates, dyadic square-bracket midpoints, grid scaling and
division, and the rational Newton centre now use this layer.  The Newton
centre first reduces the two coordinates of the polynomial and derivative
values, performs the Gaussian division through normalized scalar operations,
and persists a reduced pair.  The corresponding erasure theorem states the
same Newton formula used throughout RFAP-1.  Hurwitz continuants are not
normalized after each recurrence: their integral recurrence is
representation-sensitive, and the adjacent unit determinant already supplies
the required Gaussian coprimality without introducing a choice of unit.

\leanref{RFAP.SafeCandidates.interiorCandidate;
RFAP.DyadicBracket.SquareInterval.midpoint;
RFAP.HeightApproximation.gridCoordinate;
RFAP.LocalExtraction.anchorNewtonCenter_eq;
RFAP.NormalizedGaussian.normalizedDiv_eq}

\begin{definition}[normalized Horner and Newton envelopes]
\label{def:normalized-horner-envelope}
For a reduced scalar record let $m$ be its largest numerator-or-denominator
magnitude, and for a reduced Gaussian record take the maximum of its two
scalar values.  Put
\[
 A(x,y)=2xy,\qquad M(x,y)=xy,
\]
and define the Gaussian envelopes
\[
 G_+(x,y)=A(x,y),\qquad
 G_\times(x,y)=A(M(x,y),M(x,y)),
\]
\[
 G_N(x)=A(M(x,x),M(x,x)),\quad
 G_{\rm inv}(x)=M(x,G_N(x)),\quad
 G_{/}(x,y)=G_\times(x,G_{\rm inv}(y)).
\]
These are deliberately uniform coordinate bounds; cancellation may make the
actual records much smaller.

For a high-to-low coefficient list $c_0,\ldots,c_d$ evaluated at $z$, start
with $H_0=m(0)$ and set
\[
 H_{j+1}=\max\{H_j,
   G_+(G_\times(H_j,m(z)),m(c_j))\}.
\]
Alongside it, sum the multiplication and addition clause envelopes from
\cref{def:normalized-rational-operations}.  Call the resulting recurrences
$H(z;c_0,\ldots,c_d)$ and $T_H(z;c_0,\ldots,c_d)$.

For the Newton update, compute these two recurrences for $f$ and its displayed
coefficientwise derivative.  Extend the magnitude recurrence through one
Gaussian division and one subtraction, retaining the maximum of the two
Horner bounds, quotient bound, and returned-centre bound.  Extend the step
recurrence by the corresponding division and subtraction clause envelopes.
Call the final magnitude maximum $E_N(f,z)$, and put
$B_N(f,z)=\ell(E_N(f,z))$; call the step sum $T_N(f,z)$.
\end{definition}

\begin{lemma}[normalized Horner and Newton traces]
\label{lem:normalized-horner-envelope}
Normalized Horner evaluation erases to ordinary coefficient-list evaluation.
Every accumulator in its executed trace has component width at most
$\ell(H)$, and its accounted arithmetic-clause count is at most $T_H$.
The normalized Newton record erases to
\[
                z-\frac{f(z)}{f'(z)},
\]
every accumulator in both Horner traces and the returned centre has component
width at most $B_N(f,z)$, and the combined clause count is at most
$T_N(f,z)$.
\end{lemma}

\begin{proof}
At one Horner clause, the multiplication envelope bounds both reduced
coordinate products.  The addition envelope then bounds both coordinate
sums, while the outer maximum retains the preceding accumulator.  Induction
on the displayed coefficient list proves both the trace bound and the sum of
the per-clause costs.  Applying the same induction to the coefficientwise
derivative gives the second trace.  The Gaussian division and subtraction
lemmas then give the final two bounds.  The erasure identities at each clause
prove the two semantic formulae by the same finite induction.
\end{proof}

\leanref{RFAP.NormalizedEvaluation.value_eval;
RFAP.NormalizedEvaluation.trace_magnitude_le_envelope;
RFAP.NormalizedEvaluation.steps_le_envelope;
RFAP.NormalizedEvaluation.value_newtonCenter;
RFAP.CalculationSize.normalizedNewtonCenter_width_le_budget;
RFAP.CalculationSize.normalizedNewton_steps_le_budget;
RFAP.ApproximationPaper.normalizedNewtonCenter_eq}

\begin{lemma}[target-scan and grid-record envelopes]
\label{lem:target-grid-envelope}
Let $C,\delta>0$ be displayed rationals and put
\[
 n=\left\lceil\frac C\delta\right\rceil,
 \qquad m=\min\{k\in\N_{>0}:n<k^2\}.
\]
The bounded square scan uses at most $n+1$ comparison-bearing clauses and
satisfies
\[
                         (m-1)^2\le n<m^2.
\]

For a displayed centre $c\in\Qi$ and radius $r$, put
\[
 D=\left\lfloor\frac{1/2}{r}\right\rfloor+1,
 \qquad h=\frac1{2D},
 \qquad q=\frac UD
\]
where $U\in\Zi$ is obtained by coordinatewise nearest rounding of $Dc$.
Then
\[
 \Norm{q-c}\le2h^2,
 \qquad
 \Norm U\le
 K_{\rm grid}:=\left\lceil
       D^2\,2\bigl(\Norm c+2h^2\bigr)\right\rceil.
\]
Consequently every component of the returned fraction record has at most
\[
                   \max\{\ell(K_{\rm grid}),\ell(D)\}
\]
bits.
\end{lemma}

\begin{proof}
Minimality of the successful square candidate gives the first bracket; the
implementation has fuel $n+1$.  Nearest rounding gives a coordinate error in
$[-h,h]$, so its two squared errors sum to at most $2h^2$.  The two-leg
quadrance inequality gives
$\Norm q\le2(\Norm c+2h^2)$.  Finally, clearing the displayed denominator
gives the exact identity $\Norm U=D^2\Norm q$, and taking the rational ceiling
produces the natural bound above.
\end{proof}

\leanref{RFAP.CalculationSize.targetDenominator_bracket;
RFAP.CalculationSize.targetComparisonBudget;
RFAP.CalculationSize.gridFraction_offset_normSq_le;
RFAP.CalculationSize.gridFraction_numerator_norm_le_budget;
RFAP.ApproximationPaper.gridFraction_width_le_budget}

\begin{definition}[homogeneous-evaluation envelope]
\label{def:homogeneous-envelope}
For $U,V\in\Zi$ and a high-to-low integral coefficient list, define
\[
 E_{U,V}([])=0,\qquad
 E_{U,V}(A::L)=2\left(
    \Norm A\,\Norm U^{|L|}+\Norm V\,E_{U,V}(L)\right).
\]
Let $E^{\max}_{U,V}(L)$ be the maximum of $E_{U,V}$ over all suffixes of
$L$.  This is itself a finite recurrence, computed before evaluation.
The displayed implementation uses $|L|$ coefficient clauses and has the
multiplication budget
\[
 T_{\times}([])=0,\qquad
 T_{\times}(A::L)=|L|+2+T_{\times}(L).
\]
The term $|L|$ counts the direct finite power calculation in that clause.
\end{definition}

\begin{lemma}[homogeneous working values]
\label{lem:homogeneous-envelope}
Every Gaussian integer returned by a recursive suffix of the homogeneous
evaluation has quadrance at most $E^{\max}_{U,V}(L)$, hence component width at
most $\ell(E^{\max}_{U,V}(L))$.
\end{lemma}

\begin{proof}
Induct on the coefficient list.  Multiplicativity of Gaussian quadrance and
\[
                  \Norm{x+y}\le2(\Norm x+\Norm y)
\]
give exactly the recurrence in \cref{def:homogeneous-envelope}.  Taking the
maximum with the recursively computed suffix bound retains every intermediate
suffix value.
\end{proof}

\leanref{RFAP.CalculationSize.homogeneousNormEnvelope;
RFAP.CalculationSize.homogeneousWorkingEnvelope;
RFAP.CalculationSize.homogeneousTrace_norm_le_workingEnvelope;
RFAP.CalculationSize.homogeneousMultiplicationBudget;
RFAP.ApproximationPaper.homogeneousEvaluation_norm_le_envelope}

\begin{theorem}[finite Hurwitz arithmetic envelope]
\label{thm:hurwitz-arithmetic-envelope}
For a displayed source $z=N/D$ with $D\ne0$, let
\[
 L=\ell(\Norm D),\qquad
 A=\max\{2(\Norm N+\Norm D),4\Norm D\}.
\]
The executable Hurwitz encoding has at most $L$ digit clauses, every digit
$a_j$ satisfies $\Norm{a_j}\le A$, and the numerator and denominator of its
final continuant each have quadrance at most
\begin{equation}
\label{eq:continuant-envelope}
                         \bigl(2(A+1)\bigr)^L.
\end{equation}
Thus $L$, $\ell(A)$, and the binary width of
\eqref{eq:continuant-envelope} are explicit bounds for the recursion count,
digit components, and final continuant components respectively.
The same bound holds for every recursive suffix continuant visited by the
calculation.  Its multiplication-and-addition count is at most $2L$, in
addition to the $L$ encoding clauses.
\end{theorem}

\begin{proof}
If $R=N-aD$ is one nonterminal integral remainder, the fundamental-square
bound gives the stronger natural inequality
\[
                         2\Norm R\le\Norm D.
\]
Each transition therefore consumes at least one binary digit of denominator
quadrance, proving the bound $L$.  The division identity and the two-leg
quadrance inequality give
$\Norm a\le2(\Norm N+\Norm D)$.  After the first step, the new numerator is
the old denominator and the new denominator has smaller quadrance; the second
term in $A$ makes the same bound invariant down the complete trace.

For a digit list with $\Norm a_j\le A$, induction on the continuant recurrence
and the same two-leg inequality bound both continuant entries by
$[2(A+1)]^j$ after $j$ digits.  Monotonicity in $j\le L$ gives
\eqref{eq:continuant-envelope} for every suffix.  Each digit clause performs
one Gaussian-integer multiplication and one addition, giving $2L$.
\end{proof}

\leanref{RFAP.HurwitzRational.twice_remainder_norm_le_denominator;
RFAP.HurwitzRational.encode_length_le_width;
RFAP.HurwitzRational.encode_digit_norm_le;
RFAP.CalculationSize.continuant_norm_le_envelope;
RFAP.CalculationSize.continuantTrace_norm_le_envelope;
RFAP.CalculationSize.encodedContinuant_steps_le_budget;
RFAP.ApproximationPaper.hurwitzEncoding_length_le_budget;
RFAP.ApproximationPaper.hurwitzContinuant_norm_le_budget}

\begin{proposition}[protected quotient target]
\label{prop:protected-target-envelope}
At a protected quotient site, replace $C$ in
\cref{lem:target-grid-envelope} by the computed coefficient
$2L_1(g,A)/\eta_h$.  The same comparison budget and exact predecessor bracket
hold.  No extra unbounded search is introduced by the numerator or pole
branch.
\end{proposition}

\leanref{RFAP.CalculationSize.protectedTargetComparisonBudget;
RFAP.CalculationSize.protectedTargetDenominator_bracket}

\begin{theorem}[RFAP-2 calculation-size certificate]
\label{thm:stage-two-resource-certificate}
Suppose the displayed input consists of a polynomial $f$, an RFAP-1
Newton-band anchor and radius, the already cleared Gaussian-integral data
$U,V$ and coefficient list used for homogeneous evaluation, a displayed
Gaussian fraction for the finite Hurwitz calculation, and the two rationals
$C,\delta$ for a target scan.  Finite recurrences in precisely this data
compute natural numbers $T$ and $B$ such that:
\begin{enumerate}
\item the normalized polynomial and derivative Horner traces, Newton
  division, grid construction, homogeneous evaluation, Hurwitz encoding and
  continuant construction, and target scan use at most $T$ accounted
  arithmetic or recursive clauses;
\item every accumulator in the two normalized Horner traces, the returned
  Newton record, all three normalized grid-operation results, and the returned
  grid fraction have component width at most $B$;
\item every homogeneous-Horner suffix value, every Hurwitz digit, and both
  entries of every recursive continuant suffix have component width at most
  $B$; and
\item the selected target denominator has width at most $B$.
\end{enumerate}
Here $T$ is the sum of the source recurrences already displayed in this
section.  The number $B$ is the maximum of their Horner--Newton, grid,
homogeneous, digit, continuant, and target-width envelopes.  In particular,
neither is defined by running the calculation and then taking a maximum of
the trace.
\end{theorem}

\begin{proof}
Apply
\cref{lem:normalized-horner-envelope,lem:target-grid-envelope,lem:homogeneous-envelope,thm:hurwitz-arithmetic-envelope}
to the corresponding
displayed inputs.  Addition of their clause inequalities gives the stated
$T$.  Monotonicity of $\ell$ converts each natural component or quadrance
bound into a width bound; inclusion into the explicitly nested maximum gives
$B$.  The target-denominator term uses $m\le n+1$, which follows from the
minimal square scan.  These finite inequalities are assembled into one
proof-bearing record by the implementation.
\end{proof}

\leanref{RFAP.CalculationSize.StageTwoInput;
RFAP.CalculationSize.stageTwoStepBudget;
RFAP.CalculationSize.stageTwoBitBudget;
RFAP.CalculationSize.stageTwo_resource_certificate;
RFAP.ApproximationPaper.stageTwoResourceCertificate}

This is the RFAP-2 side of the resource statement: its input boundary is the
displayed RFAP-1 Newton-band data and cleared integral coefficient list.  The
companion paper now supplies the complementary certificate for constructing
the normalized eliminant and B\'ezout data, clearing the coefficients, and
performing every Sturm winding, safe cut, and requested subdivision round
\cite{RFAPOne}.  Its budgets are
\texttt{stageThreeStepBudget} and \texttt{stageThreeBitBudget}.  Consequently
the two step budgets add and the two width budgets combine by their maximum.
The interface remains explicit---the RFAP-2 theorem does not silently charge
its post-processing bound for work performed by RFAP-1---but there is no
longer an unbounded eliminant, Sturm, winding, clearing, or subdivision phase
between the two certificates.

\leanref{RFAP.Paper.stageThreeStepBudget;
RFAP.Paper.stageThreeBitBudget;
RFAP.Paper.stageThreeResourceCertificate}

\section{Worked calculations}
\label{sec:cubic}

This section feeds one RFAP-1 track through both branches of the glossary.
The cubic calculation exhibits the coefficient scale, height obstruction,
Newton band, and grid fraction.  Its first displayed sample is then treated
as a Gaussian fraction: the nearest-integer program produces its complete
finite trace, the continuants give the generated stages, and the finite
Lakein inequality compares every listed division-free residual with all
smaller-denominator competitors.

\subsection{The cubic height calculation}

\begin{example}[displayed cubic data]
\label{ex:cubic-data}
Consider the RFAP-1 polynomial
\[
                         f(Z)=Z^3-(1+i).
\]
Its coefficients are Gaussian integers, so $S_f=1$.  The companion
calculation returns three initial winding-one boxes with rational centres
\[
 q_1=\frac{139+37i}{128},\qquad
 q_2=\frac{-203+203i}{256},\qquad
 q_3=\frac{-37-139i}{128}.
\]
The common-denominator presentations have no common positive-integer scale,
and
\[
 H(q_1)=H(q_3)=128^2=16384,
 \qquad H(q_2)=256^2=65536.
\]
Their canonical Gaussian reductions are
\[
 q_1=\frac{51+88i}{64+64i},\qquad
 q_2=\frac{-203}{128+128i},\qquad
 q_3=\frac{51-88i}{64+64i},
\]
so their intrinsic heights are respectively
\[
                     8192,\qquad32768,\qquad8192.
\]
Thus reduction by the Gaussian factor $1+i$, which no positive-integer gcd
detects, improves every displayed common-denominator height by a factor of
two.
\end{example}

\begin{proposition}[exact cubic height obstruction]
\label{prop:cubic-obstruction}
The three residual quadrances are
\begin{align*}
 \Norm{f(q_1)}=\Norm{f(q_3)}
   &=\frac{39952889}{549755813888}
     \approx 7.27\mathbin{\times}10^{-5},\\
 \Norm{f(q_2)}
   &=\frac{537358761}{35184372088832}
     \approx 1.53\mathbin{\times}10^{-5}.
\end{align*}
Their exact scaled obstruction values are
\[
 \begin{aligned}
 H(q_1)^3\Norm{f(q_1)}
   &=H(q_3)^3\Norm{f(q_3)}=319623112,\\
 H(q_2)^3\Norm{f(q_2)}&=4298870088.
 \end{aligned}
\]
After canonical Gaussian reduction, the intrinsic obstruction values are
\[
 \begin{aligned}
 h(q_1)^3\Norm{f(q_1)}
   &=h(q_3)^3\Norm{f(q_3)}=39952889,\\
 h(q_2)^3\Norm{f(q_2)}&=537358761.
 \end{aligned}
\]
In particular, all three displayed values satisfy the lower bound of
\cref{lem:height-lower}.
\end{proposition}

\begin{proof}
For $q=(u/D)$, homogeneous evaluation is the Gaussian-integer calculation
\[
             W_f(u,D)=u^3-(1+i)D^3=D^3f(q).
\]
At the three displayed inputs it gives
\[
 \begin{aligned}
 W_f(139+37i,128)&=17594-3174i,\\
 W_f(-203+203i,256)&=-46362-46362i,\\
 W_f(-37-139i,128)&=-3174+17594i.
 \end{aligned}
\]
Squaring and adding their integer coordinates gives respectively
$319623112$, $4298870088$, and $319623112$.  The identity
$\Norm{W_f}=H(q)^3\Norm{f(q)}$ gives the exact residuals and scaled values.
For each input, coefficientwise Gaussian division by $1+i$, followed by the
unit $i$ which places the denominator in the chosen sector, gives the three
canonical fractions displayed above.  Their heights are half the respective
common-denominator heights.  Cubing that factor divides the three scaled
values by eight, giving $39952889$, $537358761$, and $39952889$.
Each residual numerator is nonzero, so its quadrance is at least one.
\end{proof}

\leanref{RFAP.DenominatorHeight.residualNumerator;
RFAP.DenominatorHeight.residualNumerator_identity;
RFAP.DenominatorHeight.scaled_height_mul_residual_eq_normNat;
RFAP.DenominatorHeight.one_le_scaled_height_mul_residual;
RFAP.GaussianDenominatorHeight.scaled_intrinsicHeight_mul_residual_eq_normNat;
RFAP.GaussianDenominatorHeight.one_le_scaled_intrinsicHeight_mul_residual}

Equivalently, the universal floors at these displayed denominators are
\[
 \Norm{f(q_1)},\Norm{f(q_3)}\ge\frac1{128^6}
       \approx2.27\mathbin{\times}10^{-13},
 \qquad
 \Norm{f(q_2)}\ge\frac1{256^6}
       \approx3.55\mathbin{\times}10^{-15}.
\]
Canonical reduction strengthens them to
\[
 \Norm{f(q_1)},\Norm{f(q_3)}\ge\frac1{8192^3}
       \approx1.82\mathbin{\times}10^{-12},
 \qquad
 \Norm{f(q_2)}\ge\frac1{32768^3}
       \approx2.84\mathbin{\times}10^{-14}.
\]

\begin{example}[three states of one shrinking track]
\label{ex:cubic-track}
The first track has the following finite refinement states.  Each centre is
the rational midpoint of its displayed box.

\begin{center}
\small
\begin{tabularx}{\textwidth}{@{}c>{\raggedright\arraybackslash}Xcc@{}}
\toprule
state & centre & $\diamq(B)$ & $H(q)$\\
\midrule
$0$ & $(139+37i)/128$
    & $1/2048$ & $128^2=16384$\\
$1$ & $(2221+595i)/2048$
    & $1/524288$ & $2048^2=4194304$\\
$2$ & $(142111+38079i)/131072$
    & $1/2147483648$ & $131072^2=17179869184$\\
\bottomrule
\end{tabularx}
\end{center}

Their exact residual quadrances are
\[
 \begin{aligned}
 \frac{39952889}{549755813888}
   &\approx7.27\mathbin{\times}10^{-5},\\
 \frac{8767783105145}{9223372036854775808}
   &\approx9.51\mathbin{\times}10^{-7},\\
 \frac{546115072853076822625}{633825300114114700748351602688}
   &\approx8.62\mathbin{\times}10^{-10}.
 \end{aligned}
\]
The corresponding nonzero-residual floors are
\[
 2^{-42}\approx2.27\mathbin{\times}10^{-13},\qquad
 2^{-66}\approx1.36\mathbin{\times}10^{-20},\qquad
 2^{-102}\approx1.97\mathbin{\times}10^{-31}.
\]
\end{example}

Every displayed residual in \cref{ex:cubic-track} lies above its obstruction
floor, as it must.  The substantial gap is useful intuition: the height
obstruction is universal, not a claim that a particular box midpoint is
optimal.

\begin{proposition}[the sufficient-height invocation]
\label{prop:cubic-sufficient}
For $A=[-7,7]^2$, the RFAP-1 coefficient calculation gives
$L_1(f,A)=460992$.  For every displayed $\delta>0$, a finite rational scan
can choose $\eps>0$ satisfying
\[
                          921984\eps^2\le\delta.
\]
Newton extraction then computes a band radius $r$ and returns a grid fraction
$q_r$ satisfying
\[
 H(q_r)=\left(\left\lfloor\frac{1/2}{r}\right\rfloor+1\right)^2,
 \qquad \Norm{f(q_r)}\le921984\eps^2\le\delta.
\]
\end{proposition}

\begin{proof}
Here $2L_1(f,A)=921984$.  For the chosen $\delta$, compute
\[
 n=\left\lceil\frac{921984}{\delta}\right\rceil
\]
and scan upward for the least positive $m$ with $n<m^2$; then put
$\eps=1/m$.  By \cref{cor:height-target}, this gives
$921984\eps^2\le\delta$.  Theorem~\ref{thm:height-upper} computes $r$ and
the exact nearest-grid denominator.  The returned grid fraction need not be
any of the three box midpoints above.
\end{proof}

\leanref{RFAP.HeightApproximation.approximate;
RFAP.HeightApproximation.Output.height_eq;
RFAP.HeightApproximation.Output.residual_le;
RFAP.HeightApproximation.approximateTarget;
RFAP.HeightApproximation.approximateTarget_residual_le}

\subsection{A finite Hurwitz trace of the first sample}

\begin{definition}[cleared error for the sample]
\label{def:cubic-cleared-error}
Regard the first midpoint as the displayed Gaussian fraction
\[
                          z=\frac{139+37i}{128}.
\]
For an integral continuant $(P_j,Q_j)$ define
\[
          E_j=(139+37i)Q_j-128P_j\in\Zi.
\]
Then coefficientwise cross multiplication gives
\begin{equation}
\label{eq:cubic-cleared-error}
                  \Norm{zQ_j-P_j}=\frac{\Norm{E_j}}{128^2}.
\end{equation}
\end{definition}

\begin{proposition}[the computed five-stage trace]
\label{prop:cubic-hurwitz-trace}
The deterministic nearest-Gaussian-integer program returns the digit list
\[
              \bigl[1,\ 1-3i,\ -2+5i,\ 2-i,\ 1-2i\bigr].
\]
Its generated continuants, Gaussian denominator heights, and cleared errors
are as follows.

\begin{center}
\scriptsize
\begin{tabular}{@{}ccccc@{}}
\toprule
$j$ & $a_j$ & $(P_j,Q_j)$ & $\Norm{Q_j}$
    & $\Norm{E_j}/128^2$\\
\midrule
$0$ & $1$ & $(1,1)$ & $1$
    & $1490/16384\approx9.09\mathbin{\times}10^{-2}$\\
$1$ & $1-3i$ & $(2-3i,1-3i)$ & $10$
    & $52/16384\approx3.17\mathbin{\times}10^{-3}$\\
$2$ & $-2+5i$ & $(12+16i,14+11i)$ & $317$
    & $10/16384\approx6.10\mathbin{\times}10^{-4}$\\
$3$ & $2-i$ & $(42+17i,40+5i)$ & $1625$
    & $2/16384\approx1.22\mathbin{\times}10^{-4}$\\
$4$ & $1-2i$ & $(88-51i,64-64i)$ & $8192$
    & $0$\\
\bottomrule
\end{tabular}
\end{center}
\end{proposition}

\begin{proof}
The complete-quotient numerator--denominator records evolve by the exact
Euclidean calculation
\[
 \begin{aligned}
 (139+37i,128)&\longmapsto(128,11+37i)\\
 &\longmapsto(11+37i,6-4i)\\
 &\longmapsto(6-4i,3-i)\\
 &\longmapsto(3-i,1+i).
 \end{aligned}
\]
At each arrow, $a_j$ is the coordinatewise nearest Gaussian integer and the
new denominator is $N_j-a_jD_j$.  At the last record,
$(3-i)-(1-2i)(1+i)=0$, so the trace terminates.  Starting from
$(P_{-1},Q_{-1})=(1,0)$ and $(P_0,Q_0)=(1,1)$, the recurrence
\eqref{eq:continuant-recurrence} gives the five displayed continuants.
Squaring the coordinates of each $Q_j$ gives
$1,10,317,1625,8192$; substituting each pair in
\cref{def:cubic-cleared-error} gives
$\Norm{E_j}=1490,52,10,2,0$.
\end{proof}

\leanref{RFAP.HurwitzRational.encode;
RFAP.HurwitzRational.continuant;
RFAP.HurwitzTrace.stages;
RFAP.HurwitzGeometry.generated_of_mem_stages}

\paragraph{Arithmetic envelopes for the displayed sample.}
The new resource recurrences can be evaluated without changing the example.
For homogeneous evaluation at $q_1$, the data are
\[
 \Norm{139+37i}=20690,\qquad \Norm{128}=16384,\qquad
 [A_0,A_1,A_2,A_3]=[1,0,0,-1-i].
\]
Reading the coefficient list from its last suffix to the full list, the
values of $E_{139+37i,128}$ are
\[
 4,\qquad 131072,\qquad 4294967296,\qquad
 158451277373328.
\]
Thus the certified suffix-record envelope has $48$ bits per integral
component.  The displayed implementation uses four coefficient clauses,
four Gaussian-integer additions, and the direct-power multiplication budget
$T_{\times}=14$.  The actual final value
$W_f=17594-3174i$ has $15$-bit components; the difference records the
deliberate slack in the uniform two-leg recurrence.

For the Hurwitz encoder,
\[
 L=\ell(16384)=15,\qquad
 A=\max\{2(20690+16384),4\cdot16384\}=74148.
\]
The theorem therefore certifies at most fifteen digit clauses, $17$-bit
digit components, and at most $258$ bits per component of the final
continuant, since
\[
 (2(A+1))^L=148298^{15}
 \approx3.69\mathbin{\times}10^{77}.
\]
The actual trace has five clauses, and the largest absolute integer
coordinate displayed in its complete quotients and continuants is $139$, an
eight-bit value.  These gaps are useful: the certified figures are
source-only worst-case envelopes, while the table records this particular
execution.

\leanref{RFAP.CalculationSize.homogeneousWorkingEnvelope;
RFAP.CalculationSize.homogeneousMultiplicationBudget;
RFAP.HurwitzRational.digitNormBudget;
RFAP.ApproximationPaper.hurwitzEncoding_length_le_budget;
RFAP.ApproximationPaper.hurwitzContinuant_norm_le_budget}

The last row also displays an explicit difference between the two height
notions.  In $\Zi$,
\[
 139+37i=(1+i)(88-51i),\qquad
 128=(1+i)(64-64i).
\]
Thus the original common-denominator record has height $16384$, while the
terminal continuant has displayed Gaussian denominator height $8192$.  It is
already reduced; multiplication of numerator and denominator by $i$ gives
the canonical fraction
\[
                         \frac{51+88i}{64+64i}.
\]
Hence the intrinsic height of the original sample is $8192$.  The common
factor appears as output of the finite Euclidean trace, the unit
normalization is an explicit four-case calculation, and the final cross
product is exactly zero.  Applying
\cref{thm:gaussian-height-lower} sharpens the
nonzero polynomial-residual floor for $q_1$ from
$2^{-42}\approx2.27\mathbin{\times}10^{-13}$ to
\[
                   \frac1{8192^3}=2^{-39}
                     \approx1.82\mathbin{\times}10^{-12}.
\]
The actual residual remains the exact value displayed in
\cref{prop:cubic-obstruction}; the stronger floor now depends on the exact
value through its canonical reduction rather than on a fortuitously chosen
presentation.

\begin{corollary}[the stagewise Lakein comparisons]
\label{cor:cubic-lakein}
For each row $j$ of \cref{prop:cubic-hurwitz-trace}, every displayed
$(P,Q)\in\Zi^2$ satisfying $Q\ne0$ and
$\Norm Q<\Norm{Q_j}$ obeys
\[
             \Norm{zQ-P}\ge\frac{\Norm{E_j}}{16384}.
\]
For example, if $\Norm Q<1625$, then
\[
             \Norm{zQ-P}\ge\frac1{8192}
                     \approx1.22\mathbin{\times}10^{-4}.
\]
\end{corollary}

\begin{proof}
Every row is a member of the computed stage trace, hence is generated.  Apply
\cref{thm:finite-lakein-full} and replace the current residual by the exact
identity \eqref{eq:cubic-cleared-error}.  For $j=3$ the table gives
$\Norm{E_3}=2$, so $2/16384=1/8192$.  At the terminal row the asserted lower
bound is zero.
\end{proof}

\leanref{RFAP.HurwitzGeometry.generated_of_mem_stages;
RFAP.HurwitzGeometry.Stage.atLeastAsGoodAs_of_generated;
RFAP.ApproximationPaper.hurwitzConvergent_atLeastAsGoodAs}

\section{Discussion}
\label{sec:discussion}

In the glossary's terms, least coefficient clearing gives exact displayed
common- and Gaussian-denominator identities, and executable Gaussian
Euclidean reduction turns the latter into an intrinsic height obstruction;
nearest-grid Newton bands construct a sufficient height certificate at
square-root target scale.  For
continued-fraction comparison, the finite trace and its continuants give
generated stages; successor geometry and the root-form invariant certify
every Lakein factor.  The source-lattice cylinder then characterises best
pairs, and the source-bounded rectangle computes their complete finite list.
The main work of the paper was not discovering a new completed object, but
turning each customary appeal to a tail, a picture, a symmetry, a lattice
sequence, or a finite choice into displayed data and a checkable certificate.

\subsection{Replacing the infinite argument by a finite one}

Classical accounts of Hurwitz continued fractions naturally discuss an
infinite expansion, its complete quotients, and the complex value represented
by the expansion~\cite{HurwitzComplexCF,LeVequeGaussianCF,LakeinApproximation}.
For the rational inputs of this paper those objects are unnecessary, but
removing them is not a word-for-word translation.  A reference to ``the
tail'' must become one displayed integral remainder and its proof of
nonzeroness.  A backwards argument along an infinite sequence must become
induction on the constructors of a finite generated trace.  A limiting
identity must become the exact stage-reconstruction identity
\eqref{eq:stage-represents}.

This does not dispute the mathematical conclusions of the earlier proofs.
It changes their dependency boundary.  The result proved here requires no
completed source value, infinite digit stream, convergence theorem, or
selection of a limiting object.  Once those conveniences are unavailable,
however, several steps that look routine in prose have to be supplied as
substantial pieces of mathematics.

The same issue appears in the height argument.  Density of rational points
in a completed plane would say too little: it would neither construct the
denominator nor relate it to the residual certificate.  The grid fraction of
\cref{def:newton-grid} instead records the exact denominator chosen from the
Newton-band radius, and \cref{lem:grid-mesh} proves the exact uniform
minimality property of that choice.  The bounded square-threshold scan turns
a residual request into a reciprocal-integer $\eps$ without an appeal to a
real square root.  At the other end, the lower bound is not a distance from
an assumed root; it is the quadrance floor of the explicitly computed
Gaussian integer $W_f(U,V)$.  The necessary obstruction and sufficient
construction are therefore two different finite certificates rather than one
asymptotic existence statement.

\subsection{What the geometric shorthand concealed}

The succession part of the classical proof is often conveyed by a diagram:
a complete quotient, or its inverse, lies outside one of several discs.  A
diagram does not identify the certificate needed by the present theorem.
For each orientation we had to name the relevant Gaussian unit, clear an
explicitly nonzero denominator, expand the resulting quadrance
coefficientwise, and prove its sign by rational inequalities.  Distances and
square roots contribute nothing useful here; quadrances preserve the exact
finite algebra.

The deterministic tie rule makes this more than bureaucratic expansion.  Our
fundamental square is half-open:
\[
                         [-1/2,1/2)^2.
\]
Some successor inequalities are therefore strict while their reflected
counterparts are weak.  A phrase such as ``the remaining cases follow by
symmetry'' can erase precisely the strict inequality that excludes a
boundary digit.  The proof in \cref{sec:succession-geometry} retains every
orientation.  When a weak boundary remains, it supplies a second exterior
certificate rather than pretending that reflection preserved the tie rule.

Likewise, saying that only finitely many Gaussian integers remain is not yet
a constructive classification.  The proof first derives coordinate bounds,
then displays the bounded integers, and finally assigns every case to an
explicit certificate.  This is why the twenty-eight forbidden successors in
\cref{def:forbidden-successors} are mathematical data, not an incidental
implementation table.

\subsection{The exceptional state machine}

The quadrance-five case was the point at which informal reconstruction was
most unreliable.  Numerical exploration suggested the eight Hermitian root
forms and was useful for finding the finite state machine.  It was not a
proof.  In an early exploration, a pulled-back form whose quadratic
coefficient vanished was mistakenly treated as infeasible.  Such forms are
feasible, and some occur in actual finite Hurwitz traces.

The corrected proof never discards those states.  It establishes the exact
pullback identity, classifies the twenty-four exceptional first pullbacks,
and checks the $24\cdot32=768$ bounded second-step substitutions.  A result is
accepted only if it is an original root form, an explicitly nonnegative
augmentation, or a form with a displayed slack certificate.  The remaining
branch supplies one of the forbidden adjacent digit pairs and is ruled out
by the exact succession law.  Induction then transports the eight root
inequalities along the finite trace.

This episode gives a useful standard for computer-assisted exploration.
Floating-point searches may reveal a plausible list, but they do not belong
to the theorem's trust boundary.  Every transition used here is justified by
an integral coefficient identity; every discarded transition has a rational
inequality or a named forbidden-successor witness; and the finite tables are
checked by the Lean kernel without a new native-decision axiom.

\subsection{Formalisation and the audited boundary}

Our broader complaints about Lean's authoring language, unnecessary
choice-bearing routes through Mathlib, and the interface exposed by
\texttt{Quot.sound} are elaborated in the companion paper~\cite{RFAPOne}.
That paper also describes the eventual replacement of the quotient interface
by canonical reduced records.  The complaints need not be repeated here.

The paper-specific point is that the additional Hurwitz and all-best
machinery preserves that audited boundary.  Executable Gaussian division,
continuant recursion, rational square inequalities, finite Hermitian-form
classification, cylinder minimality, and bounded candidate filtering leave
the public closure at \texttt{[propext]}; neither
\texttt{Classical.choice}, \texttt{Quot.sound}, nor a native decision axiom
appears.

\subsection{Scope and next questions}

The denominator-height obstruction is universal.  Its coefficient scale is
least for the normalized coefficient list, its scaled residual is exactly an
integral quadrance, and \cref{thm:height-lower-sharp} shows joint sharpness in
the displayed scale, height, and degree.  The Gaussian-denominator theorem,
the two explicit presentation conversions, and executable gcd reduction
remove the earlier gap between the common, displayed Gaussian, and intrinsic
height records.  Unit normalization makes the reduced record canonical, so
equal displayed values have equal intrinsic heights without selecting from a
quotient.

The Newton denominator is least for the uniform strict half-mesh condition,
not necessarily for one particular centre or polynomial.  Likewise,
\cref{lem:newton-band-half-sharp} makes the geometric factor $1/2$ exact
under the stated band constraints, but does not claim that a particular
eliminant-selected band is widest or that its ambient Lipschitz majorant is
sharp.  The analogous Horner majorants and B\'ezout protection margins at a
rational-function site can contain further problem-specific slack.

The calculation-size question is now closed at the intended level in both
papers.  RFAP-I supplies source-computed $T$ and $B$ bounds for its normalized
eliminant, B\'ezout, coefficient-clearing, Sturm, winding, safe-cut, and
requested-subdivision calculations; \cref{thm:stage-two-resource-certificate}
does the same for the normalized Newton and RFAP-II post-processing family.
The bounds are finite certificates, not sharp complexity estimates.  A
logarithmic gcd analysis, tighter subresultant and Sturm growth, common-grid
subdivision, or a more detailed machine-cost model could reduce them, but no
such refinement is needed to bound the intermediate displayed records used
by the two papers.

The Hurwitz--Lakein theorem is also precise about what it does not claim.  It
compares the division-free residual of each generated continuant with every
competitor of strictly smaller Gaussian denominator quadrance.  The new
all-best theorem is a separate result: for a displayed rational source it
proves the minimal-vector equivalence, derives a finite search rectangle, and
returns every best pair.  It also decides the exact consecutive occupied
denominator levels by filtering that rectangle.  Finite descent recovers the
equivalent empty-mixed-cylinder condition; the Euclidean complement and
half-square boundary calculation show that every consecutive determinant
quadrance is one or two; and the Cramer certificate explains the determinant
algebra without a quotient carrier.
The terminal Hurwitz pair is best, but the paper does not claim that the
ordinary Hurwitz trace enumerates the rest.
Chevallier's algorithm supplies the stronger structural answer for completed
complex inputs by advancing through consecutive minimal vectors
\cite{ChevallierGaussLattices}.  Our coefficient-local generator now derives
a complete bounded candidate list from one adjacent pair and validates it by
comparisons confined to that list.  Its executable update is therefore local,
although it uses a bounded coefficient square rather than formalising
Chevallier's more economical structural update formula.  Our finite
Euclidean and half-square calculation proves his index-one/index-two determinant
classification for displayed rational sources.  The coordinate-parity
calculation also proves the centered-coset classification and the
unit-and-divisible-coefficient normal form for a putative repeated index-two
transition.  Four centered candidates give the factor-four error estimate,
and an exact Cauchy calculation restricts the remaining coefficient to the
twelve Gaussian integers of quadrance one, two, or four.  Exact normalization,
unit-associate reduction, and the three finite cylinder calculations then
rule out every value: an index-two transition is necessarily followed by an
index-one transition.  This completes the displayed-rational form of
Chevallier's repeated-index lemma and proves the no-skipped-level property of
the coefficient-local generator and validator.  The normalized
nearest-coefficient calculation in \cref{thm:finite-dirichlet-window} gives
the generic root-free quadratic gap.  Five small coefficient tests, after
one of eight finite coordinate symmetries, split into four rational regions
and strengthen the gap to $205/256$, bounding an index-one window by
$256/205$.  Four centered unit certificates force the two normalized
coordinate quadrances above $15/56$; retaining their orientation and the
affine minorant $2(U+1)/7$ raises the index-two gap to $205/128$ and gives the
same window bound.  Replacing the two fixed affine parameters by
$s=\tau/3$ and $m=s/(1-s^2)$ proves the exact rational-cut family
\[
 \frac34\le\tau<1,\quad\tau^2-12\tau+9\ge0
 \quad\Longrightarrow\quad
 \tau\mathcal D_j\le1,\quad
 \tau\sum_{j<n}\mathcal D_j\le n.
\]
In the conventional completed notation, the reciprocal of the smaller
quadratic endpoint is Chevallier's optimal squared constant
$(2+\sqrt3)/3\approx1.24$.  The converse finite theorem now proves that for
every displayed rational $\tau$ in the same range with
$\tau^2-12\tau+9<0$, and every natural scale $T$, the Pell recurrence
\eqref{eq:critical-pell-recurrence} constructs a displayed rational source
and consecutive index-one window beyond $T$ satisfying
$1<\tau\mathcal D$.  Thus every displayed rational strengthening of the cut
has explicit counterexamples at unbounded requested scales.  The arbitrary
finite same-source construction strengthens this to any supplied run length;
one proper rational box preserves the whole run.  At every displayed rational
source in that box, stopped coefficient-local successor dynamics reaches a
selected critical state with probability one in blocks bounded by the source
denominator quadrance plus one.  The only non-successor boundary operations
are the explicit terminal reset and absorbing seed hold.  The universal half
is now global over displayed rational data as well: for every natural squared
request $R>1$, finite maximization of the best-pair table returns a denominator
quadrance below $R$ with $\tau R\Norm{qz-p}^{2}\le1$.

What remains absent is Chevallier's stronger quantifier order: almost every
fixed completed complex input defeats every improved constant at arbitrarily
large scales.  That statement uses the infinite completed setting and
ergodic arguments outside the present discipline; none of the finite reset,
stopping, or box-content statements is presented as a substitute for Haar
invariance.

Within that boundary, the two parts of the paper now complement one another.
The height theory measures what a finite root or pole approximation costs;
the Hurwitz theory proves that a particular finite Gaussian Euclidean trace
uses its denominator budget efficiently in the exact second-kind sense.  Both
conclusions are calculations on supplied records, and neither requires the
completed values that those records may conventionally be said to approach.

\begin{thebibliography}{99}
\bibitem{RFAPOne}
R. Lawrence,
\emph{A finite constructive Fundamental Theorem of Algebra and
rational-function argument principle over the Gaussian rationals},
companion manuscript and Lean formalisation, 2026.

\bibitem{Richman}
F. Richman,
\emph{The fundamental theorem of algebra: a constructive development without
choice},
Pacific Journal of Mathematics 196 (2000), 213--230.

\bibitem{PerrucciRoyQuantitative}
D. Perrucci and M.-F. Roy,
\emph{Quantitative fundamental theorem of algebra},
Quarterly Journal of Mathematics 70 (2019), 1009--1037;
arXiv:1803.04358.

\bibitem{HurwitzComplexCF}
A. Hurwitz,
\emph{\"Uber die Entwicklung complexer Gr\"ossen in Kettenbr\"uche},
Acta Mathematica 11 (1887), 187--200;
doi:10.1007/BF02612324.

\bibitem{LeVequeGaussianCF}
W. J. LeVeque,
\emph{Continued fractions and approximations in $k(i)$. I, II},
Nederl. Akad. Wetensch. Proc. Ser. A 55 = Indagationes Mathematicae 14
(1952), 526--535 and 536--545.

\bibitem{LakeinApproximation}
R. B. Lakein,
\emph{Approximation properties of some complex continued fractions},
Monatshefte f\"ur Mathematik 77 (1973), 396--403;
doi:10.1007/BF01295317.

\bibitem{DaniNogueira}
S. G. Dani and A. Nogueira,
\emph{Continued fractions for complex numbers and values of binary quadratic
forms},
Transactions of the American Mathematical Society 366 (2014), 3553--3583.

\bibitem{EiItoNakadaNatsui}
H. Ei, S. Ito, H. Nakada, and R. Natsui,
\emph{On the construction of the natural extension of the Hurwitz complex
continued fraction map},
Monatshefte f\"ur Mathematik 188 (2019), 37--86.

\bibitem{ChevallierGaussLattices}
N. Chevallier,
\emph{Gauss lattices and complex continued fractions},
arXiv:2101.05480, 2021.

\bibitem{EberlContinuedFractions}
M. Eberl,
\emph{Continued Fractions},
Archive of Formal Proofs, 2024,
\url{https://isa-afp.org/entries/Continued_Fractions.html}.
\end{thebibliography}

\end{document}
